% !TEX program = lualatex
% ---------------------------------------------------------------------------------------------
%  Readback note 30: The tower method at a place: floor measures, gauges and the bounds on each floor
%
%  Build:  python3 tools/rb.py note <slug>
%
%  One unit of a Lean artifact, rewritten as ordinary mathematics. The prose is hand-written. Lean names,
%  line ranges and counts come from data.tex; quoted Lean statements are copied from the artifact by the build.
%
%  Rules of the text: standard terminology; every symbol is explained where it first appears and listed again
%  in the notation at the end; no letter has two meanings; a proof is complete or says what it takes for granted;
%  Lean names appear only in the LEAN lines and in the concordance.
%
%  The signature rule (the build enforces it): a function, relation or predicate that the artifact defines is
%  written only through its notation macro, which the build makes from the spec (`notation`) with one slot for
%  each explicit argument of the Lean definition. No argument is fixed and dropped, also not "for brevity" and
%  not for one section. Every argument and every hypothesis of a stated declaration stands inside its
%  statement, or in a Setting that the statement cites by number; only implicit and instance arguments (the
%  ambient objects) may stand in the conventions.
% ---------------------------------------------------------------------------------------------
\documentclass[11pt,a4paper]{article}
\ifdefined\pdfvariable\pdfvariable suppressoptionalinfo \numexpr 1+2+4+8\relax\fi
\usepackage[a4paper,hmargin=35mm,top=22mm,bottom=27mm,footskip=12mm,nohead]{geometry}

\usepackage{fontspec}
\directlua{
  luaotfload.add_fallback("gfallback", {"STIXTwoMath-Regular.otf:mode=node;", "DejaVuSans.ttf:mode=node;"})
  luaotfload.add_fallback("gmonofallback", {"DejaVuSansMono.ttf:mode=node;", "STIXTwoMath-Regular.otf:mode=node;"})
}
\setmainfont{NewCM10}[Extension=.otf,UprightFont=*-Book,ItalicFont=*-BookItalic,
  BoldFont=*-Bold,BoldItalicFont=*-BoldItalic,SmallCapsFont=*-Book,SmallCapsFeatures={Letters=SmallCaps},
  RawFeature={fallback=gfallback}]
\setsansfont{NewCMSans10}[Extension=.otf,UprightFont=*-Book,ItalicFont=*-BookOblique,
  BoldFont=*-Bold,BoldItalicFont=*-BoldOblique,RawFeature={fallback=gfallback}]
\setmonofont{JetBrainsMono}[Extension=.otf,UprightFont=*-Regular,ItalicFont=*-Italic,
  BoldFont=*-SemiBold,BoldItalicFont=*-SemiBoldItalic,Scale=MatchLowercase,
  RawFeature={-calt;-liga;fallback=gmonofallback}]
\usepackage{amsmath}
\usepackage{amsthm}
\usepackage[math-style=TeX]{unicode-math}
\setmathfont{NewCMMath-Book.otf}
\emergencystretch=2.5em\relax

\usepackage{xcolor}
\usepackage{tikz}
\usetikzlibrary{calc,arrows.meta,positioning,backgrounds}
\usepackage[protrusion=true,expansion=true]{microtype}
\usepackage{enumitem}
\usepackage{needspace}
\usepackage{booktabs}
\usepackage{array}
\usepackage{longtable}
\usepackage{fancyhdr}
\usepackage{lastpage}
\usepackage{multicol}
\usepackage[hidelinks,pdfusetitle]{hyperref}

\definecolor{ink}{HTML}{0F172A}
\definecolor{inksoft}{HTML}{475569}
\definecolor{inkfaint}{HTML}{94A3B8}
\definecolor{rulegrey}{HTML}{CBD5E1}
\definecolor{arith}{HTML}{C2410C}
\definecolor{est}{HTML}{7C3AED}
\definecolor{ana}{HTML}{0891B2}
\definecolor{ok}{HTML}{047857}
\definecolor{leancol}{HTML}{3730A3}

\title{The tower method at a place: floor measures, gauges and the bounds on each floor}
\author{grh-compact}
\date{}

% ------------------------------------------------------------------------------------- macros
\newcommand{\leanfont}{\ttfamily\fontsize{8.6}{10.5}\selectfont\color{leancol}}
\newcommand{\lean}[1]{\csname LN@#1\endcsname}
\newcommand{\fillin}[1]{{\color{arith}\itshape[#1]}}  % a place of the template that is not written yet; the build reports every one that is left
\newcommand{\leanquote}[1]{\csname LQ@#1\endcsname}
% generated by rb.py note; do not edit
\newcommand{\ConcDecls}{79}
\newcommand{\Pin}{72fead10}
\newcommand{\PinText}{that of commit {\leanfont 72fead10}}
\newcommand{\ArtifactName}{grh-compact}
\expandafter\newcommand\csname LN@hasSchurBoundofconvolveeq\endcsname{\href{https://github.com/lsernau/grh-lean/blob/72fead102f663d5c7a407964a13dc0cf9d977207/grh-compact/src/New/RealPlace/BesselCut.lean\#L27-L64}{{\leanfont IsBesselFunction.\allowbreak{}hasSchurBound\_of\_convolve\_eq}}}
\expandafter\newcommand\csname LN@existscutisCutPair\endcsname{\href{https://github.com/lsernau/grh-lean/blob/72fead102f663d5c7a407964a13dc0cf9d977207/grh-compact/src/New/RealPlace/BesselCut.lean\#L72-L157}{{\leanfont HasPeakAt.\allowbreak{}exists\_cut\_isCutPair}}}
\expandafter\newcommand\csname LN@orthonormal\endcsname{\href{https://github.com/lsernau/grh-lean/blob/72fead102f663d5c7a407964a13dc0cf9d977207/grh-compact/src/New/RealPlace/BesselOrthonormal.lean\#L10-L26}{{\leanfont BesselSystem.\allowbreak{}orthonormal}}}
\expandafter\newcommand\csname LN@openCellChart\endcsname{\href{https://github.com/lsernau/grh-lean/blob/72fead102f663d5c7a407964a13dc0cf9d977207/grh-compact/src/New/RealPlace/CellNeighbourhood.lean\#L15-L17}{{\leanfont openCellChart}}}
\expandafter\newcommand\csname LN@chartCoords\endcsname{\href{https://github.com/lsernau/grh-lean/blob/72fead102f663d5c7a407964a13dc0cf9d977207/grh-compact/src/New/RealPlace/CellNeighbourhood.lean\#L30-L34}{{\leanfont chartCoords}}}
\expandafter\newcommand\csname LN@besselTangent\endcsname{\href{https://github.com/lsernau/grh-lean/blob/72fead102f663d5c7a407964a13dc0cf9d977207/grh-compact/src/New/RealPlace/CellNeighbourhood.lean\#L38-L41}{{\leanfont besselTangent}}}
\expandafter\newcommand\csname LN@logDerivUunipotentOfmulembedH\endcsname{\href{https://github.com/lsernau/grh-lean/blob/72fead102f663d5c7a407964a13dc0cf9d977207/grh-compact/src/New/RealPlace/CellNeighbourhood.lean\#L43-L53}{{\leanfont logDerivU\_unipotentOf\_mul\_embedH}}}
\expandafter\newcommand\csname LN@existssymmeqchartCoordsskewOfReal\endcsname{\href{https://github.com/lsernau/grh-lean/blob/72fead102f663d5c7a407964a13dc0cf9d977207/grh-compact/src/New/RealPlace/CellNeighbourhood.lean\#L55-L61}{{\leanfont exists\_symm\_eq\_chartCoords\_skewOfReal}}}
\expandafter\newcommand\csname LN@openCellChartjacobian\endcsname{\href{https://github.com/lsernau/grh-lean/blob/72fead102f663d5c7a407964a13dc0cf9d977207/grh-compact/src/New/RealPlace/CellNeighbourhood.lean\#L67-L114}{{\leanfont openCellChart\_jacobian}}}
\expandafter\newcommand\csname LN@openCellChartinjOn\endcsname{\href{https://github.com/lsernau/grh-lean/blob/72fead102f663d5c7a407964a13dc0cf9d977207/grh-compact/src/New/RealPlace/CellNeighbourhood.lean\#L117-L131}{{\leanfont openCellChart\_injOn}}}
\expandafter\newcommand\csname LN@memopenCellandboundsofnear\endcsname{\href{https://github.com/lsernau/grh-lean/blob/72fead102f663d5c7a407964a13dc0cf9d977207/grh-compact/src/New/RealPlace/CellNeighbourhood.lean\#L138-L163}{{\leanfont mem\_openCell\_and\_bounds\_of\_near}}}
\expandafter\newcommand\csname LN@euclideanenormsqeqsum\endcsname{\href{https://github.com/lsernau/grh-lean/blob/72fead102f663d5c7a407964a13dc0cf9d977207/grh-compact/src/New/RealPlace/CellNeighbourhood.lean\#L165-L166}{{\leanfont euclidean\_enorm\_sq\_eq\_sum}}}
\expandafter\newcommand\csname LN@RealPlaceSizes\endcsname{\href{https://github.com/lsernau/grh-lean/blob/72fead102f663d5c7a407964a13dc0cf9d977207/grh-compact/src/New/RealPlace/ComplexGauge.lean\#L18-L25}{{\leanfont RealPlaceSizes}}}
\expandafter\newcommand\csname LN@haarRatio\endcsname{\href{https://github.com/lsernau/grh-lean/blob/72fead102f663d5c7a407964a13dc0cf9d977207/grh-compact/src/New/RealPlace/ComplexGauge.lean\#L45-L46}{{\leanfont RealPlaceSizes.\allowbreak{}haarRatio}}}
\expandafter\newcommand\csname LN@haarRatiospec\endcsname{\href{https://github.com/lsernau/grh-lean/blob/72fead102f663d5c7a407964a13dc0cf9d977207/grh-compact/src/New/RealPlace/ComplexGauge.lean\#L50-L111}{{\leanfont RealPlaceSizes.\allowbreak{}haarRatio\_spec}}}
\expandafter\newcommand\csname LQ@haarRatiospec\endcsname{\leanquoteline{0}{lemma RealPlaceSizes.haarRatio\_spec : (hF.haarRatio ≠ 0 ∧ hF.haarRatio ≠ ∞) ∧}\leanquoteline{2}{hF.complexLevel.groupMeasure = hF.haarRatio • ℙ ∧ (∀ f, Measurable f → besselU.FixesLeft f → ∫⁻ g,}\leanquoteline{2}{f g ∂hF.complexLevel.upperMeasure 0 = hF.haarRatio * ∫⁻ g, f g ∂besselQuotMeasure) ∧ .ofReal}\leanquoteline{2}{(volU 𝔰.k / volU 𝔰.r) * hF.haarRatio ≤ 1}}
\expandafter\newcommand\csname LN@existsorthonormalkerlintegralle\endcsname{\href{https://github.com/lsernau/grh-lean/blob/72fead102f663d5c7a407964a13dc0cf9d977207/grh-compact/src/New/RealPlace/ComplexGauge.lean\#L126-L152}{{\leanfont exists\_orthonormal\_ker\_lintegral\_le}}}
\expandafter\newcommand\csname LN@tenle\endcsname{\href{https://github.com/lsernau/grh-lean/blob/72fead102f663d5c7a407964a13dc0cf9d977207/grh-compact/src/New/RealPlace/ComplexGauge.lean\#L154-L155}{{\leanfont RealPlaceSizes.\allowbreak{}ten\_le}}}
\expandafter\newcommand\csname LN@gauge\endcsname{\href{https://github.com/lsernau/grh-lean/blob/72fead102f663d5c7a407964a13dc0cf9d977207/grh-compact/src/New/RealPlace/ComplexGauge.lean\#L158-L214}{{\leanfont RealPlaceSizes.\allowbreak{}gauge}}}
\expandafter\newcommand\csname LN@continuousframeVector\endcsname{\href{https://github.com/lsernau/grh-lean/blob/72fead102f663d5c7a407964a13dc0cf9d977207/grh-compact/src/New/RealPlace/FrameNorm.lean\#L12-L13}{{\leanfont continuous\_frameVector}}}
\expandafter\newcommand\csname LN@frameVectormul\endcsname{\href{https://github.com/lsernau/grh-lean/blob/72fead102f663d5c7a407964a13dc0cf9d977207/grh-compact/src/New/RealPlace/FrameNorm.lean\#L15-L17}{{\leanfont frameVector\_mul}}}
\expandafter\newcommand\csname LN@frameVectoreqsum\endcsname{\href{https://github.com/lsernau/grh-lean/blob/72fead102f663d5c7a407964a13dc0cf9d977207/grh-compact/src/New/RealPlace/FrameNorm.lean\#L19-L23}{{\leanfont frameVector\_eq\_sum}}}
\expandafter\newcommand\csname LN@normframeVectorle\endcsname{\href{https://github.com/lsernau/grh-lean/blob/72fead102f663d5c7a407964a13dc0cf9d977207/grh-compact/src/New/RealPlace/FrameNorm.lean\#L25-L26}{{\leanfont norm\_frameVector\_le}}}
\expandafter\newcommand\csname LN@normsumlemulnormsumframeVector\endcsname{\href{https://github.com/lsernau/grh-lean/blob/72fead102f663d5c7a407964a13dc0cf9d977207/grh-compact/src/New/RealPlace/FrameNorm.lean\#L28-L31}{{\leanfont norm\_sum\_le\_mul\_norm\_sum\_frameVector}}}
\expandafter\newcommand\csname LN@besselSizele\endcsname{\href{https://github.com/lsernau/grh-lean/blob/72fead102f663d5c7a407964a13dc0cf9d977207/grh-compact/src/New/RealPlace/FrameNorm.lean\#L33-L34}{{\leanfont besselSize\_le}}}
\expandafter\newcommand\csname LN@existsnormleofbesselSizele\endcsname{\href{https://github.com/lsernau/grh-lean/blob/72fead102f663d5c7a407964a13dc0cf9d977207/grh-compact/src/New/RealPlace/FrameNorm.lean\#L38-L42}{{\leanfont exists\_norm\_le\_of\_besselSize\_le}}}
\expandafter\newcommand\csname LN@besselSizebesselSmul\endcsname{\href{https://github.com/lsernau/grh-lean/blob/72fead102f663d5c7a407964a13dc0cf9d977207/grh-compact/src/New/RealPlace/FrameNorm.lean\#L44-L45}{{\leanfont besselSize\_besselS\_mul}}}
\expandafter\newcommand\csname LN@besselSizelenormmul\endcsname{\href{https://github.com/lsernau/grh-lean/blob/72fead102f663d5c7a407964a13dc0cf9d977207/grh-compact/src/New/RealPlace/FrameNorm.lean\#L47-L50}{{\leanfont besselSize\_le\_norm\_mul}}}
\expandafter\newcommand\csname LN@continuousbesselSize\endcsname{\href{https://github.com/lsernau/grh-lean/blob/72fead102f663d5c7a407964a13dc0cf9d977207/grh-compact/src/New/RealPlace/FrameNorm.lean\#L52-L55}{{\leanfont continuous\_besselSize}}}
\expandafter\newcommand\csname LN@norminclMatrixleone\endcsname{\href{https://github.com/lsernau/grh-lean/blob/72fead102f663d5c7a407964a13dc0cf9d977207/grh-compact/src/New/RealPlace/FrameNorm.lean\#L57-L59}{{\leanfont norm\_inclMatrix\_le\_one}}}
\expandafter\newcommand\csname LN@normJ\endcsname{\href{https://github.com/lsernau/grh-lean/blob/72fead102f663d5c7a407964a13dc0cf9d977207/grh-compact/src/New/RealPlace/FrameNorm.lean\#L61-L61}{{\leanfont norm\_J₁}}}
\expandafter\newcommand\csname LN@normleofmemframeBox\endcsname{\href{https://github.com/lsernau/grh-lean/blob/72fead102f663d5c7a407964a13dc0cf9d977207/grh-compact/src/New/RealPlace/FrameNorm.lean\#L64-L75}{{\leanfont norm\_le\_of\_mem\_frameBox}}}
\expandafter\newcommand\csname LN@norminnerJle\endcsname{\href{https://github.com/lsernau/grh-lean/blob/72fead102f663d5c7a407964a13dc0cf9d977207/grh-compact/src/New/RealPlace/FrameNorm.lean\#L77-L78}{{\leanfont norm\_inner\_J₁\_le}}}
\expandafter\newcommand\csname LN@realFrameGrameqinner\endcsname{\href{https://github.com/lsernau/grh-lean/blob/72fead102f663d5c7a407964a13dc0cf9d977207/grh-compact/src/New/RealPlace/FrameNorm.lean\#L80-L82}{{\leanfont realFrameGram\_eq\_inner}}}
\expandafter\newcommand\csname LN@realFrameGrammulinvariant\endcsname{\href{https://github.com/lsernau/grh-lean/blob/72fead102f663d5c7a407964a13dc0cf9d977207/grh-compact/src/New/RealPlace/FrameNorm.lean\#L84-L86}{{\leanfont realFrameGram\_mul\_invariant}}}
\expandafter\newcommand\csname LN@realFrameGrameqframeMatrixmul\endcsname{\href{https://github.com/lsernau/grh-lean/blob/72fead102f663d5c7a407964a13dc0cf9d977207/grh-compact/src/New/RealPlace/FrameNorm.lean\#L88-L90}{{\leanfont realFrameGram\_eq\_frameMatrix\_mul}}}
\expandafter\newcommand\csname LN@normrealFrameGramlemul\endcsname{\href{https://github.com/lsernau/grh-lean/blob/72fead102f663d5c7a407964a13dc0cf9d977207/grh-compact/src/New/RealPlace/FrameNorm.lean\#L92-L93}{{\leanfont norm\_realFrameGram\_le\_mul}}}
\expandafter\newcommand\csname LN@besselSizeleofdual\endcsname{\href{https://github.com/lsernau/grh-lean/blob/72fead102f663d5c7a407964a13dc0cf9d977207/grh-compact/src/New/RealPlace/FrameNorm.lean\#L97-L127}{{\leanfont besselSize\_le\_of\_dual}}}
\expandafter\newcommand\csname LN@sectionMatrixmulframeCols\endcsname{\href{https://github.com/lsernau/grh-lean/blob/72fead102f663d5c7a407964a13dc0cf9d977207/grh-compact/src/New/RealPlace/GramSection.lean\#L17-L18}{{\leanfont sectionMatrix\_mul\_frameCols}}}
\expandafter\newcommand\csname LN@frameColsgramsectionMatrix\endcsname{\href{https://github.com/lsernau/grh-lean/blob/72fead102f663d5c7a407964a13dc0cf9d977207/grh-compact/src/New/RealPlace/GramSection.lean\#L20-L21}{{\leanfont frameCols\_gram\_sectionMatrix}}}
\expandafter\newcommand\csname LN@gramSection\endcsname{\href{https://github.com/lsernau/grh-lean/blob/72fead102f663d5c7a407964a13dc0cf9d977207/grh-compact/src/New/RealPlace/GramSection.lean\#L25-L40}{{\leanfont gramSection}}}
\expandafter\newcommand\csname LN@measurablegramSection\endcsname{\href{https://github.com/lsernau/grh-lean/blob/72fead102f663d5c7a407964a13dc0cf9d977207/grh-compact/src/New/RealPlace/GramSection.lean\#L42-L45}{{\leanfont measurable\_gramSection}}}
\expandafter\newcommand\csname LN@diagLastsqrtmulinv\endcsname{\href{https://github.com/lsernau/grh-lean/blob/72fead102f663d5c7a407964a13dc0cf9d977207/grh-compact/src/New/RealPlace/GramSection.lean\#L49-L50}{{\leanfont diagLast\_sqrt\_mul\_inv}}}
\expandafter\newcommand\csname LN@coegramSection\endcsname{\href{https://github.com/lsernau/grh-lean/blob/72fead102f663d5c7a407964a13dc0cf9d977207/grh-compact/src/New/RealPlace/GramSection.lean\#L55-L55}{{\leanfont coe\_gramSection}}}
\expandafter\newcommand\csname LN@gramSectioneventuallyEq\endcsname{\href{https://github.com/lsernau/grh-lean/blob/72fead102f663d5c7a407964a13dc0cf9d977207/grh-compact/src/New/RealPlace/GramSection.lean\#L57-L58}{{\leanfont gramSection\_eventuallyEq}}}
\expandafter\newcommand\csname LN@realFrameGramgramSection\endcsname{\href{https://github.com/lsernau/grh-lean/blob/72fead102f663d5c7a407964a13dc0cf9d977207/grh-compact/src/New/RealPlace/GramSection.lean\#L61-L62}{{\leanfont realFrameGram\_gramSection}}}
\expandafter\newcommand\csname LN@existseqmulgramSectionmul\endcsname{\href{https://github.com/lsernau/grh-lean/blob/72fead102f663d5c7a407964a13dc0cf9d977207/grh-compact/src/New/RealPlace/GramSection.lean\#L66-L80}{{\leanfont exists\_eq\_mul\_gramSection\_mul}}}
\expandafter\newcommand\csname LN@invCorner\endcsname{\href{https://github.com/lsernau/grh-lean/blob/72fead102f663d5c7a407964a13dc0cf9d977207/grh-compact/src/New/RealPlace/GramSectionBounds.lean\#L18-L18}{{\leanfont invCorner}}}
\expandafter\newcommand\csname LN@openCell\endcsname{\href{https://github.com/lsernau/grh-lean/blob/72fead102f663d5c7a407964a13dc0cf9d977207/grh-compact/src/New/RealPlace/GramSectionBounds.lean\#L22-L22}{{\leanfont Shape.\allowbreak{}openCell}}}
\expandafter\newcommand\csname LN@cornerDefect\endcsname{\href{https://github.com/lsernau/grh-lean/blob/72fead102f663d5c7a407964a13dc0cf9d977207/grh-compact/src/New/RealPlace/GramSectionBounds.lean\#L26-L26}{{\leanfont cornerDefect}}}
\expandafter\newcommand\csname LN@IsBoundedInCell\endcsname{\href{https://github.com/lsernau/grh-lean/blob/72fead102f663d5c7a407964a13dc0cf9d977207/grh-compact/src/New/RealPlace/GramSectionBounds.lean\#L30-L30}{{\leanfont IsBoundedInCell}}}
\expandafter\newcommand\csname LN@diagLast\endcsname{\href{https://github.com/lsernau/grh-lean/blob/72fead102f663d5c7a407964a13dc0cf9d977207/grh-compact/src/New/RealPlace/GramSectionBounds.lean\#L33-L33}{{\leanfont diagLast}}}
\expandafter\newcommand\csname LN@sectionCoeff\endcsname{\href{https://github.com/lsernau/grh-lean/blob/72fead102f663d5c7a407964a13dc0cf9d977207/grh-compact/src/New/RealPlace/GramSectionBounds.lean\#L39-L40}{{\leanfont sectionCoeff}}}
\expandafter\newcommand\csname LN@sectionFrame\endcsname{\href{https://github.com/lsernau/grh-lean/blob/72fead102f663d5c7a407964a13dc0cf9d977207/grh-compact/src/New/RealPlace/GramSectionBounds.lean\#L44-L44}{{\leanfont sectionFrame}}}
\expandafter\newcommand\csname LN@sectionMatrix\endcsname{\href{https://github.com/lsernau/grh-lean/blob/72fead102f663d5c7a407964a13dc0cf9d977207/grh-compact/src/New/RealPlace/GramSectionBounds.lean\#L48-L49}{{\leanfont sectionMatrix}}}
\expandafter\newcommand\csname LN@gramDiscRealeq\endcsname{\href{https://github.com/lsernau/grh-lean/blob/72fead102f663d5c7a407964a13dc0cf9d977207/grh-compact/src/New/RealPlace/GramSectionBounds.lean\#L57-L58}{{\leanfont gramDiscReal\_eq}}}
\expandafter\newcommand\csname LN@memopenCelliff\endcsname{\href{https://github.com/lsernau/grh-lean/blob/72fead102f663d5c7a407964a13dc0cf9d977207/grh-compact/src/New/RealPlace/GramSectionBounds.lean\#L60-L61}{{\leanfont mem\_openCell\_iff}}}
\expandafter\newcommand\csname LN@continuousgramDiscReal\endcsname{\href{https://github.com/lsernau/grh-lean/blob/72fead102f663d5c7a407964a13dc0cf9d977207/grh-compact/src/New/RealPlace/GramSectionBounds.lean\#L63-L64}{{\leanfont continuous\_gramDiscReal}}}
\expandafter\newcommand\csname LN@isOpenopenCell\endcsname{\href{https://github.com/lsernau/grh-lean/blob/72fead102f663d5c7a407964a13dc0cf9d977207/grh-compact/src/New/RealPlace/GramSectionBounds.lean\#L67-L68}{{\leanfont isOpen\_openCell}}}
\expandafter\newcommand\csname LN@differentiableAtsectionMatrixuncurry\endcsname{\href{https://github.com/lsernau/grh-lean/blob/72fead102f663d5c7a407964a13dc0cf9d977207/grh-compact/src/New/RealPlace/GramSectionBounds.lean\#L72-L79}{{\leanfont differentiableAt\_sectionMatrix\_uncurry}}}
\expandafter\newcommand\csname LN@cornerDefectpos\endcsname{\href{https://github.com/lsernau/grh-lean/blob/72fead102f663d5c7a407964a13dc0cf9d977207/grh-compact/src/New/RealPlace/GramSectionBounds.lean\#L81-L81}{{\leanfont cornerDefect\_pos}}}
\expandafter\newcommand\csname LN@sqrteqsqrtcornerDefect\endcsname{\href{https://github.com/lsernau/grh-lean/blob/72fead102f663d5c7a407964a13dc0cf9d977207/grh-compact/src/New/RealPlace/GramSectionBounds.lean\#L83-L84}{{\leanfont sqrt\_eq\_sqrt\_cornerDefect}}}
\expandafter\newcommand\csname LN@differentiableAtsectionMatrix\endcsname{\href{https://github.com/lsernau/grh-lean/blob/72fead102f663d5c7a407964a13dc0cf9d977207/grh-compact/src/New/RealPlace/GramSectionBounds.lean\#L89-L91}{{\leanfont differentiableAt\_sectionMatrix}}}
\expandafter\newcommand\csname LN@normsectionMatrixle\endcsname{\href{https://github.com/lsernau/grh-lean/blob/72fead102f663d5c7a407964a13dc0cf9d977207/grh-compact/src/New/RealPlace/GramSectionBounds.lean\#L96-L114}{{\leanfont norm\_sectionMatrix\_le}}}
\expandafter\newcommand\csname LN@opencell\endcsname{\href{https://github.com/lsernau/grh-lean/blob/72fead102f663d5c7a407964a13dc0cf9d977207/grh-compact/src/New/RealPlace/OpenCellBound.lean\#L16-L100}{{\leanfont BesselSystem.\allowbreak{}open\_cell}}}
\expandafter\newcommand\csname LQ@opencell\endcsname{\leanquoteline{0}{lemma BesselSystem.open\_cell \{𝔰 B f ρ γ y\} (X : BesselSystem 𝔰 B f) (hγ : 0 < gramDiscReal}\leanquoteline{2}{(realFrameGram γ)) (hy : y ∈ frameBox ρ) (hz : γ⁻¹ * y ∈ frameBox ρ) (hρ : B ≤ ρ) : ∃ A n m, ∃ t :}\leanquoteline{2}{U (1 : Matrix 𝔰.κ 𝔰.κ ℂ) × (𝔰.middleGram (1 : ℂ)).besselU0 → \_, IsCompact A ∧ Measurable t ∧}\leanquoteline{2}{IsL2Contraction m (volume.restrict A) n (fun (i, j) x y ↦}\leanquoteline{2}{X.W .full (γ * t x * y) i * conj (X.W .full (t x * y) j)) ∧ IsProbabilityMeasure m ∧}\leanquoteline{2}{IsFiniteMeasure n ∧ n.real .univ ≤ cellMass γ X.deg ρ (card X.I) ∧ ∀ᵐ x ∂m, t x ∈ besselU0 ⊔}\leanquoteline{2}{besselH ∧ ∀ y ∈ frameBox ρ, besselSize (γ * t x * y) ≤ ρ → y ∈ A}}
\expandafter\newcommand\csname LN@lelintegralofgramSection\endcsname{\href{https://github.com/lsernau/grh-lean/blob/72fead102f663d5c7a407964a13dc0cf9d977207/grh-compact/src/New/RealPlace/SobolevCell.lean\#L19-L33}{{\leanfont le\_lintegral\_of\_gramSection}}}
\expandafter\newcommand\csname LN@lintegralmulleofframeMatrixle\endcsname{\href{https://github.com/lsernau/grh-lean/blob/72fead102f663d5c7a407964a13dc0cf9d977207/grh-compact/src/New/RealPlace/SobolevCell.lean\#L38-L54}{{\leanfont lintegral\_mul\_le\_of\_frameMatrix\_le}}}
\expandafter\newcommand\csname LN@lintegralgramSectionintegralle\endcsname{\href{https://github.com/lsernau/grh-lean/blob/72fead102f663d5c7a407964a13dc0cf9d977207/grh-compact/src/New/RealPlace/SobolevCell.lean\#L62-L134}{{\leanfont lintegral\_gramSection\_integral\_le}}}
\expandafter\newcommand\csname LN@tsuminvindexpowleofcongr\endcsname{\href{https://github.com/lsernau/grh-lean/blob/72fead102f663d5c7a407964a13dc0cf9d977207/grh-compact/src/New/RealPlace/VolumeKernel.lean\#L20-L48}{{\leanfont Filling.\allowbreak{}tsum\_inv\_index\_pow\_le\_of\_congr}}}
\expandafter\newcommand\csname LN@middlesharedvectorle\endcsname{\href{https://github.com/lsernau/grh-lean/blob/72fead102f663d5c7a407964a13dc0cf9d977207/grh-compact/src/New/RealPlace/VolumeKernel.lean\#L55-L61}{{\leanfont Filling.\allowbreak{}HasLevel.\allowbreak{}middle\_shared\_vector\_le}}}
\expandafter\newcommand\csname LN@coleqzeroofnormltone\endcsname{\href{https://github.com/lsernau/grh-lean/blob/72fead102f663d5c7a407964a13dc0cf9d977207/grh-compact/src/New/RealPlace/VolumeKernel.lean\#L63-L68}{{\leanfont Filling.\allowbreak{}col\_eq\_zero\_of\_norm\_lt\_one}}}
\expandafter\newcommand\csname LN@eqofframeImageeq\endcsname{\href{https://github.com/lsernau/grh-lean/blob/72fead102f663d5c7a407964a13dc0cf9d977207/grh-compact/src/New/RealPlace/VolumeKernel.lean\#L70-L81}{{\leanfont Filling.\allowbreak{}IsLabelledTerm.\allowbreak{}eq\_of\_frameImage\_eq}}}
\expandafter\newcommand\csname LN@existsframeImagetoRowseqadd\endcsname{\href{https://github.com/lsernau/grh-lean/blob/72fead102f663d5c7a407964a13dc0cf9d977207/grh-compact/src/New/RealPlace/VolumeKernel.lean\#L87-L122}{{\leanfont Filling.\allowbreak{}HasLevel.\allowbreak{}exists\_frameImage\_toRows₂\_eq\_add}}}
\expandafter\newcommand\csname LN@frameVolume\endcsname{\href{https://github.com/lsernau/grh-lean/blob/72fead102f663d5c7a407964a13dc0cf9d977207/grh-compact/src/New/RealPlace/VolumeKernel.lean\#L131-L131}{{\leanfont frameVolume}}}
\expandafter\newcommand\csname LN@frameXPart\endcsname{\href{https://github.com/lsernau/grh-lean/blob/72fead102f663d5c7a407964a13dc0cf9d977207/grh-compact/src/New/RealPlace/VolumeKernel.lean\#L135-L135}{{\leanfont frameX₀Part}}}
\expandafter\newcommand\csname LN@frameVolumepos\endcsname{\href{https://github.com/lsernau/grh-lean/blob/72fead102f663d5c7a407964a13dc0cf9d977207/grh-compact/src/New/RealPlace/VolumeKernel.lean\#L138-L144}{{\leanfont frameVolume\_pos}}}
\expandafter\newcommand\csname LN@frameVolumemulleft\endcsname{\href{https://github.com/lsernau/grh-lean/blob/72fead102f663d5c7a407964a13dc0cf9d977207/grh-compact/src/New/RealPlace/VolumeKernel.lean\#L154-L156}{{\leanfont frameVolume\_mul\_left}}}
\expandafter\newcommand\csname LN@continuousframeVolume\endcsname{\href{https://github.com/lsernau/grh-lean/blob/72fead102f663d5c7a407964a13dc0cf9d977207/grh-compact/src/New/RealPlace/VolumeKernel.lean\#L158-L158}{{\leanfont continuous\_frameVolume}}}
\expandafter\newcommand\csname LN@existskernelframeVolumele\endcsname{\href{https://github.com/lsernau/grh-lean/blob/72fead102f663d5c7a407964a13dc0cf9d977207/grh-compact/src/New/RealPlace/VolumeKernel.lean\#L164-L235}{{\leanfont exists\_kernel\_frameVolume\_le}}}
\expandafter\newcommand\csname LN@frameXPartunipotentOfinvmul\endcsname{\href{https://github.com/lsernau/grh-lean/blob/72fead102f663d5c7a407964a13dc0cf9d977207/grh-compact/src/New/RealPlace/VolumeKernel.lean\#L237-L239}{{\leanfont frameX₀Part\_unipotentOf\_inv\_mul}}}
\newcommand{\ConcordanceRows}{\crow{hasSchurBoundofconvolveeq}{\href{https://github.com/lsernau/grh-lean/blob/72fead102f663d5c7a407964a13dc0cf9d977207/grh-compact/src/New/RealPlace/BesselCut.lean\#L27-L64}{{\leanfont IsBesselFunction.\allowbreak{}hasSchurBound\_of\_convolve\_eq}}}{BesselCut.lean}{27--64}
\crow{existscutisCutPair}{\href{https://github.com/lsernau/grh-lean/blob/72fead102f663d5c7a407964a13dc0cf9d977207/grh-compact/src/New/RealPlace/BesselCut.lean\#L72-L157}{{\leanfont HasPeakAt.\allowbreak{}exists\_cut\_isCutPair}}}{BesselCut.lean}{72--157}
\crow{orthonormal}{\href{https://github.com/lsernau/grh-lean/blob/72fead102f663d5c7a407964a13dc0cf9d977207/grh-compact/src/New/RealPlace/BesselOrthonormal.lean\#L10-L26}{{\leanfont BesselSystem.\allowbreak{}orthonormal}}}{BesselOrthonormal.lean}{10--26}
\crow{openCellChart}{\href{https://github.com/lsernau/grh-lean/blob/72fead102f663d5c7a407964a13dc0cf9d977207/grh-compact/src/New/RealPlace/CellNeighbourhood.lean\#L15-L17}{{\leanfont openCellChart}}}{CellNeighbourhood.lean}{15--17}
\crow{chartCoords}{\href{https://github.com/lsernau/grh-lean/blob/72fead102f663d5c7a407964a13dc0cf9d977207/grh-compact/src/New/RealPlace/CellNeighbourhood.lean\#L30-L34}{{\leanfont chartCoords}}}{CellNeighbourhood.lean}{30--34}
\crow{besselTangent}{\href{https://github.com/lsernau/grh-lean/blob/72fead102f663d5c7a407964a13dc0cf9d977207/grh-compact/src/New/RealPlace/CellNeighbourhood.lean\#L38-L41}{{\leanfont besselTangent}}}{CellNeighbourhood.lean}{38--41}
\crow{logDerivUunipotentOfmulembedH}{\href{https://github.com/lsernau/grh-lean/blob/72fead102f663d5c7a407964a13dc0cf9d977207/grh-compact/src/New/RealPlace/CellNeighbourhood.lean\#L43-L53}{{\leanfont logDerivU\_unipotentOf\_mul\_embedH}}}{CellNeighbourhood.lean}{43--53}
\crow{existssymmeqchartCoordsskewOfReal}{\href{https://github.com/lsernau/grh-lean/blob/72fead102f663d5c7a407964a13dc0cf9d977207/grh-compact/src/New/RealPlace/CellNeighbourhood.lean\#L55-L61}{{\leanfont exists\_symm\_eq\_chartCoords\_skewOfReal}}}{CellNeighbourhood.lean}{55--61}
\crow{openCellChartjacobian}{\href{https://github.com/lsernau/grh-lean/blob/72fead102f663d5c7a407964a13dc0cf9d977207/grh-compact/src/New/RealPlace/CellNeighbourhood.lean\#L67-L114}{{\leanfont openCellChart\_jacobian}}}{CellNeighbourhood.lean}{67--114}
\crow{openCellChartinjOn}{\href{https://github.com/lsernau/grh-lean/blob/72fead102f663d5c7a407964a13dc0cf9d977207/grh-compact/src/New/RealPlace/CellNeighbourhood.lean\#L117-L131}{{\leanfont openCellChart\_injOn}}}{CellNeighbourhood.lean}{117--131}
\crow{memopenCellandboundsofnear}{\href{https://github.com/lsernau/grh-lean/blob/72fead102f663d5c7a407964a13dc0cf9d977207/grh-compact/src/New/RealPlace/CellNeighbourhood.lean\#L138-L163}{{\leanfont mem\_openCell\_and\_bounds\_of\_near}}}{CellNeighbourhood.lean}{138--163}
\crow{euclideanenormsqeqsum}{\href{https://github.com/lsernau/grh-lean/blob/72fead102f663d5c7a407964a13dc0cf9d977207/grh-compact/src/New/RealPlace/CellNeighbourhood.lean\#L165-L166}{{\leanfont euclidean\_enorm\_sq\_eq\_sum}}}{CellNeighbourhood.lean}{165--166}
\crow{RealPlaceSizes}{\href{https://github.com/lsernau/grh-lean/blob/72fead102f663d5c7a407964a13dc0cf9d977207/grh-compact/src/New/RealPlace/ComplexGauge.lean\#L18-L25}{{\leanfont RealPlaceSizes}}}{ComplexGauge.lean}{18--25}
\crow{haarRatio}{\href{https://github.com/lsernau/grh-lean/blob/72fead102f663d5c7a407964a13dc0cf9d977207/grh-compact/src/New/RealPlace/ComplexGauge.lean\#L45-L46}{{\leanfont RealPlaceSizes.\allowbreak{}haarRatio}}}{ComplexGauge.lean}{45--46}
\crow{haarRatiospec}{\href{https://github.com/lsernau/grh-lean/blob/72fead102f663d5c7a407964a13dc0cf9d977207/grh-compact/src/New/RealPlace/ComplexGauge.lean\#L50-L111}{{\leanfont RealPlaceSizes.\allowbreak{}haarRatio\_spec}}}{ComplexGauge.lean}{50--111}
\crow{existsorthonormalkerlintegralle}{\href{https://github.com/lsernau/grh-lean/blob/72fead102f663d5c7a407964a13dc0cf9d977207/grh-compact/src/New/RealPlace/ComplexGauge.lean\#L126-L152}{{\leanfont exists\_orthonormal\_ker\_lintegral\_le}}}{ComplexGauge.lean}{126--152}
\crow{tenle}{\href{https://github.com/lsernau/grh-lean/blob/72fead102f663d5c7a407964a13dc0cf9d977207/grh-compact/src/New/RealPlace/ComplexGauge.lean\#L154-L155}{{\leanfont RealPlaceSizes.\allowbreak{}ten\_le}}}{ComplexGauge.lean}{154--155}
\crow{gauge}{\href{https://github.com/lsernau/grh-lean/blob/72fead102f663d5c7a407964a13dc0cf9d977207/grh-compact/src/New/RealPlace/ComplexGauge.lean\#L158-L214}{{\leanfont RealPlaceSizes.\allowbreak{}gauge}}}{ComplexGauge.lean}{158--214}
\crow{continuousframeVector}{\href{https://github.com/lsernau/grh-lean/blob/72fead102f663d5c7a407964a13dc0cf9d977207/grh-compact/src/New/RealPlace/FrameNorm.lean\#L12-L13}{{\leanfont continuous\_frameVector}}}{FrameNorm.lean}{12--13}
\crow{frameVectormul}{\href{https://github.com/lsernau/grh-lean/blob/72fead102f663d5c7a407964a13dc0cf9d977207/grh-compact/src/New/RealPlace/FrameNorm.lean\#L15-L17}{{\leanfont frameVector\_mul}}}{FrameNorm.lean}{15--17}
\crow{frameVectoreqsum}{\href{https://github.com/lsernau/grh-lean/blob/72fead102f663d5c7a407964a13dc0cf9d977207/grh-compact/src/New/RealPlace/FrameNorm.lean\#L19-L23}{{\leanfont frameVector\_eq\_sum}}}{FrameNorm.lean}{19--23}
\crow{normframeVectorle}{\href{https://github.com/lsernau/grh-lean/blob/72fead102f663d5c7a407964a13dc0cf9d977207/grh-compact/src/New/RealPlace/FrameNorm.lean\#L25-L26}{{\leanfont norm\_frameVector\_le}}}{FrameNorm.lean}{25--26}
\crow{normsumlemulnormsumframeVector}{\href{https://github.com/lsernau/grh-lean/blob/72fead102f663d5c7a407964a13dc0cf9d977207/grh-compact/src/New/RealPlace/FrameNorm.lean\#L28-L31}{{\leanfont norm\_sum\_le\_mul\_norm\_sum\_frameVector}}}{FrameNorm.lean}{28--31}
\crow{besselSizele}{\href{https://github.com/lsernau/grh-lean/blob/72fead102f663d5c7a407964a13dc0cf9d977207/grh-compact/src/New/RealPlace/FrameNorm.lean\#L33-L34}{{\leanfont besselSize\_le}}}{FrameNorm.lean}{33--34}
\crow{existsnormleofbesselSizele}{\href{https://github.com/lsernau/grh-lean/blob/72fead102f663d5c7a407964a13dc0cf9d977207/grh-compact/src/New/RealPlace/FrameNorm.lean\#L38-L42}{{\leanfont exists\_norm\_le\_of\_besselSize\_le}}}{FrameNorm.lean}{38--42}
\crow{besselSizebesselSmul}{\href{https://github.com/lsernau/grh-lean/blob/72fead102f663d5c7a407964a13dc0cf9d977207/grh-compact/src/New/RealPlace/FrameNorm.lean\#L44-L45}{{\leanfont besselSize\_besselS\_mul}}}{FrameNorm.lean}{44--45}
\crow{besselSizelenormmul}{\href{https://github.com/lsernau/grh-lean/blob/72fead102f663d5c7a407964a13dc0cf9d977207/grh-compact/src/New/RealPlace/FrameNorm.lean\#L47-L50}{{\leanfont besselSize\_le\_norm\_mul}}}{FrameNorm.lean}{47--50}
\crow{continuousbesselSize}{\href{https://github.com/lsernau/grh-lean/blob/72fead102f663d5c7a407964a13dc0cf9d977207/grh-compact/src/New/RealPlace/FrameNorm.lean\#L52-L55}{{\leanfont continuous\_besselSize}}}{FrameNorm.lean}{52--55}
\crow{norminclMatrixleone}{\href{https://github.com/lsernau/grh-lean/blob/72fead102f663d5c7a407964a13dc0cf9d977207/grh-compact/src/New/RealPlace/FrameNorm.lean\#L57-L59}{{\leanfont norm\_inclMatrix\_le\_one}}}{FrameNorm.lean}{57--59}
\crow{normJ}{\href{https://github.com/lsernau/grh-lean/blob/72fead102f663d5c7a407964a13dc0cf9d977207/grh-compact/src/New/RealPlace/FrameNorm.lean\#L61-L61}{{\leanfont norm\_J₁}}}{FrameNorm.lean}{61--61}
\crow{normleofmemframeBox}{\href{https://github.com/lsernau/grh-lean/blob/72fead102f663d5c7a407964a13dc0cf9d977207/grh-compact/src/New/RealPlace/FrameNorm.lean\#L64-L75}{{\leanfont norm\_le\_of\_mem\_frameBox}}}{FrameNorm.lean}{64--75}
\crow{norminnerJle}{\href{https://github.com/lsernau/grh-lean/blob/72fead102f663d5c7a407964a13dc0cf9d977207/grh-compact/src/New/RealPlace/FrameNorm.lean\#L77-L78}{{\leanfont norm\_inner\_J₁\_le}}}{FrameNorm.lean}{77--78}
\crow{realFrameGrameqinner}{\href{https://github.com/lsernau/grh-lean/blob/72fead102f663d5c7a407964a13dc0cf9d977207/grh-compact/src/New/RealPlace/FrameNorm.lean\#L80-L82}{{\leanfont realFrameGram\_eq\_inner}}}{FrameNorm.lean}{80--82}
\crow{realFrameGrammulinvariant}{\href{https://github.com/lsernau/grh-lean/blob/72fead102f663d5c7a407964a13dc0cf9d977207/grh-compact/src/New/RealPlace/FrameNorm.lean\#L84-L86}{{\leanfont realFrameGram\_mul\_invariant}}}{FrameNorm.lean}{84--86}
\crow{realFrameGrameqframeMatrixmul}{\href{https://github.com/lsernau/grh-lean/blob/72fead102f663d5c7a407964a13dc0cf9d977207/grh-compact/src/New/RealPlace/FrameNorm.lean\#L88-L90}{{\leanfont realFrameGram\_eq\_frameMatrix\_mul}}}{FrameNorm.lean}{88--90}
\crow{normrealFrameGramlemul}{\href{https://github.com/lsernau/grh-lean/blob/72fead102f663d5c7a407964a13dc0cf9d977207/grh-compact/src/New/RealPlace/FrameNorm.lean\#L92-L93}{{\leanfont norm\_realFrameGram\_le\_mul}}}{FrameNorm.lean}{92--93}
\crow{besselSizeleofdual}{\href{https://github.com/lsernau/grh-lean/blob/72fead102f663d5c7a407964a13dc0cf9d977207/grh-compact/src/New/RealPlace/FrameNorm.lean\#L97-L127}{{\leanfont besselSize\_le\_of\_dual}}}{FrameNorm.lean}{97--127}
\crow{sectionMatrixmulframeCols}{\href{https://github.com/lsernau/grh-lean/blob/72fead102f663d5c7a407964a13dc0cf9d977207/grh-compact/src/New/RealPlace/GramSection.lean\#L17-L18}{{\leanfont sectionMatrix\_mul\_frameCols}}}{GramSection.lean}{17--18}
\crow{frameColsgramsectionMatrix}{\href{https://github.com/lsernau/grh-lean/blob/72fead102f663d5c7a407964a13dc0cf9d977207/grh-compact/src/New/RealPlace/GramSection.lean\#L20-L21}{{\leanfont frameCols\_gram\_sectionMatrix}}}{GramSection.lean}{20--21}
\crow{gramSection}{\href{https://github.com/lsernau/grh-lean/blob/72fead102f663d5c7a407964a13dc0cf9d977207/grh-compact/src/New/RealPlace/GramSection.lean\#L25-L40}{{\leanfont gramSection}}}{GramSection.lean}{25--40}
\crow{measurablegramSection}{\href{https://github.com/lsernau/grh-lean/blob/72fead102f663d5c7a407964a13dc0cf9d977207/grh-compact/src/New/RealPlace/GramSection.lean\#L42-L45}{{\leanfont measurable\_gramSection}}}{GramSection.lean}{42--45}
\crow{diagLastsqrtmulinv}{\href{https://github.com/lsernau/grh-lean/blob/72fead102f663d5c7a407964a13dc0cf9d977207/grh-compact/src/New/RealPlace/GramSection.lean\#L49-L50}{{\leanfont diagLast\_sqrt\_mul\_inv}}}{GramSection.lean}{49--50}
\crow{coegramSection}{\href{https://github.com/lsernau/grh-lean/blob/72fead102f663d5c7a407964a13dc0cf9d977207/grh-compact/src/New/RealPlace/GramSection.lean\#L55-L55}{{\leanfont coe\_gramSection}}}{GramSection.lean}{55--55}
\crow{gramSectioneventuallyEq}{\href{https://github.com/lsernau/grh-lean/blob/72fead102f663d5c7a407964a13dc0cf9d977207/grh-compact/src/New/RealPlace/GramSection.lean\#L57-L58}{{\leanfont gramSection\_eventuallyEq}}}{GramSection.lean}{57--58}
\crow{realFrameGramgramSection}{\href{https://github.com/lsernau/grh-lean/blob/72fead102f663d5c7a407964a13dc0cf9d977207/grh-compact/src/New/RealPlace/GramSection.lean\#L61-L62}{{\leanfont realFrameGram\_gramSection}}}{GramSection.lean}{61--62}
\crow{existseqmulgramSectionmul}{\href{https://github.com/lsernau/grh-lean/blob/72fead102f663d5c7a407964a13dc0cf9d977207/grh-compact/src/New/RealPlace/GramSection.lean\#L66-L80}{{\leanfont exists\_eq\_mul\_gramSection\_mul}}}{GramSection.lean}{66--80}
\crow{invCorner}{\href{https://github.com/lsernau/grh-lean/blob/72fead102f663d5c7a407964a13dc0cf9d977207/grh-compact/src/New/RealPlace/GramSectionBounds.lean\#L18-L18}{{\leanfont invCorner}}}{GramSectionBounds.lean}{18--18}
\crow{openCell}{\href{https://github.com/lsernau/grh-lean/blob/72fead102f663d5c7a407964a13dc0cf9d977207/grh-compact/src/New/RealPlace/GramSectionBounds.lean\#L22-L22}{{\leanfont Shape.\allowbreak{}openCell}}}{GramSectionBounds.lean}{22--22}
\crow{cornerDefect}{\href{https://github.com/lsernau/grh-lean/blob/72fead102f663d5c7a407964a13dc0cf9d977207/grh-compact/src/New/RealPlace/GramSectionBounds.lean\#L26-L26}{{\leanfont cornerDefect}}}{GramSectionBounds.lean}{26--26}
\crow{IsBoundedInCell}{\href{https://github.com/lsernau/grh-lean/blob/72fead102f663d5c7a407964a13dc0cf9d977207/grh-compact/src/New/RealPlace/GramSectionBounds.lean\#L30-L30}{{\leanfont IsBoundedInCell}}}{GramSectionBounds.lean}{30--30}
\crow{diagLast}{\href{https://github.com/lsernau/grh-lean/blob/72fead102f663d5c7a407964a13dc0cf9d977207/grh-compact/src/New/RealPlace/GramSectionBounds.lean\#L33-L33}{{\leanfont diagLast}}}{GramSectionBounds.lean}{33--33}
\crow{sectionCoeff}{\href{https://github.com/lsernau/grh-lean/blob/72fead102f663d5c7a407964a13dc0cf9d977207/grh-compact/src/New/RealPlace/GramSectionBounds.lean\#L39-L40}{{\leanfont sectionCoeff}}}{GramSectionBounds.lean}{39--40}
\crow{sectionFrame}{\href{https://github.com/lsernau/grh-lean/blob/72fead102f663d5c7a407964a13dc0cf9d977207/grh-compact/src/New/RealPlace/GramSectionBounds.lean\#L44-L44}{{\leanfont sectionFrame}}}{GramSectionBounds.lean}{44--44}
\crow{sectionMatrix}{\href{https://github.com/lsernau/grh-lean/blob/72fead102f663d5c7a407964a13dc0cf9d977207/grh-compact/src/New/RealPlace/GramSectionBounds.lean\#L48-L49}{{\leanfont sectionMatrix}}}{GramSectionBounds.lean}{48--49}
\crow{gramDiscRealeq}{\href{https://github.com/lsernau/grh-lean/blob/72fead102f663d5c7a407964a13dc0cf9d977207/grh-compact/src/New/RealPlace/GramSectionBounds.lean\#L57-L58}{{\leanfont gramDiscReal\_eq}}}{GramSectionBounds.lean}{57--58}
\crow{memopenCelliff}{\href{https://github.com/lsernau/grh-lean/blob/72fead102f663d5c7a407964a13dc0cf9d977207/grh-compact/src/New/RealPlace/GramSectionBounds.lean\#L60-L61}{{\leanfont mem\_openCell\_iff}}}{GramSectionBounds.lean}{60--61}
\crow{continuousgramDiscReal}{\href{https://github.com/lsernau/grh-lean/blob/72fead102f663d5c7a407964a13dc0cf9d977207/grh-compact/src/New/RealPlace/GramSectionBounds.lean\#L63-L64}{{\leanfont continuous\_gramDiscReal}}}{GramSectionBounds.lean}{63--64}
\crow{isOpenopenCell}{\href{https://github.com/lsernau/grh-lean/blob/72fead102f663d5c7a407964a13dc0cf9d977207/grh-compact/src/New/RealPlace/GramSectionBounds.lean\#L67-L68}{{\leanfont isOpen\_openCell}}}{GramSectionBounds.lean}{67--68}
\crow{differentiableAtsectionMatrixuncurry}{\href{https://github.com/lsernau/grh-lean/blob/72fead102f663d5c7a407964a13dc0cf9d977207/grh-compact/src/New/RealPlace/GramSectionBounds.lean\#L72-L79}{{\leanfont differentiableAt\_sectionMatrix\_uncurry}}}{GramSectionBounds.lean}{72--79}
\crow{cornerDefectpos}{\href{https://github.com/lsernau/grh-lean/blob/72fead102f663d5c7a407964a13dc0cf9d977207/grh-compact/src/New/RealPlace/GramSectionBounds.lean\#L81-L81}{{\leanfont cornerDefect\_pos}}}{GramSectionBounds.lean}{81--81}
\crow{sqrteqsqrtcornerDefect}{\href{https://github.com/lsernau/grh-lean/blob/72fead102f663d5c7a407964a13dc0cf9d977207/grh-compact/src/New/RealPlace/GramSectionBounds.lean\#L83-L84}{{\leanfont sqrt\_eq\_sqrt\_cornerDefect}}}{GramSectionBounds.lean}{83--84}
\crow{differentiableAtsectionMatrix}{\href{https://github.com/lsernau/grh-lean/blob/72fead102f663d5c7a407964a13dc0cf9d977207/grh-compact/src/New/RealPlace/GramSectionBounds.lean\#L89-L91}{{\leanfont differentiableAt\_sectionMatrix}}}{GramSectionBounds.lean}{89--91}
\crow{normsectionMatrixle}{\href{https://github.com/lsernau/grh-lean/blob/72fead102f663d5c7a407964a13dc0cf9d977207/grh-compact/src/New/RealPlace/GramSectionBounds.lean\#L96-L114}{{\leanfont norm\_sectionMatrix\_le}}}{GramSectionBounds.lean}{96--114}
\crow{opencell}{\href{https://github.com/lsernau/grh-lean/blob/72fead102f663d5c7a407964a13dc0cf9d977207/grh-compact/src/New/RealPlace/OpenCellBound.lean\#L16-L100}{{\leanfont BesselSystem.\allowbreak{}open\_cell}}}{OpenCellBound.lean}{16--100}
\crow{lelintegralofgramSection}{\href{https://github.com/lsernau/grh-lean/blob/72fead102f663d5c7a407964a13dc0cf9d977207/grh-compact/src/New/RealPlace/SobolevCell.lean\#L19-L33}{{\leanfont le\_lintegral\_of\_gramSection}}}{SobolevCell.lean}{19--33}
\crow{lintegralmulleofframeMatrixle}{\href{https://github.com/lsernau/grh-lean/blob/72fead102f663d5c7a407964a13dc0cf9d977207/grh-compact/src/New/RealPlace/SobolevCell.lean\#L38-L54}{{\leanfont lintegral\_mul\_le\_of\_frameMatrix\_le}}}{SobolevCell.lean}{38--54}
\crow{lintegralgramSectionintegralle}{\href{https://github.com/lsernau/grh-lean/blob/72fead102f663d5c7a407964a13dc0cf9d977207/grh-compact/src/New/RealPlace/SobolevCell.lean\#L62-L134}{{\leanfont lintegral\_gramSection\_integral\_le}}}{SobolevCell.lean}{62--134}
\crow{tsuminvindexpowleofcongr}{\href{https://github.com/lsernau/grh-lean/blob/72fead102f663d5c7a407964a13dc0cf9d977207/grh-compact/src/New/RealPlace/VolumeKernel.lean\#L20-L48}{{\leanfont Filling.\allowbreak{}tsum\_inv\_index\_pow\_le\_of\_congr}}}{VolumeKernel.lean}{20--48}
\crow{middlesharedvectorle}{\href{https://github.com/lsernau/grh-lean/blob/72fead102f663d5c7a407964a13dc0cf9d977207/grh-compact/src/New/RealPlace/VolumeKernel.lean\#L55-L61}{{\leanfont Filling.\allowbreak{}HasLevel.\allowbreak{}middle\_shared\_vector\_le}}}{VolumeKernel.lean}{55--61}
\crow{coleqzeroofnormltone}{\href{https://github.com/lsernau/grh-lean/blob/72fead102f663d5c7a407964a13dc0cf9d977207/grh-compact/src/New/RealPlace/VolumeKernel.lean\#L63-L68}{{\leanfont Filling.\allowbreak{}col\_eq\_zero\_of\_norm\_lt\_one}}}{VolumeKernel.lean}{63--68}
\crow{eqofframeImageeq}{\href{https://github.com/lsernau/grh-lean/blob/72fead102f663d5c7a407964a13dc0cf9d977207/grh-compact/src/New/RealPlace/VolumeKernel.lean\#L70-L81}{{\leanfont Filling.\allowbreak{}IsLabelledTerm.\allowbreak{}eq\_of\_frameImage\_eq}}}{VolumeKernel.lean}{70--81}
\crow{existsframeImagetoRowseqadd}{\href{https://github.com/lsernau/grh-lean/blob/72fead102f663d5c7a407964a13dc0cf9d977207/grh-compact/src/New/RealPlace/VolumeKernel.lean\#L87-L122}{{\leanfont Filling.\allowbreak{}HasLevel.\allowbreak{}exists\_frameImage\_toRows₂\_eq\_add}}}{VolumeKernel.lean}{87--122}
\crow{frameVolume}{\href{https://github.com/lsernau/grh-lean/blob/72fead102f663d5c7a407964a13dc0cf9d977207/grh-compact/src/New/RealPlace/VolumeKernel.lean\#L131-L131}{{\leanfont frameVolume}}}{VolumeKernel.lean}{131--131}
\crow{frameXPart}{\href{https://github.com/lsernau/grh-lean/blob/72fead102f663d5c7a407964a13dc0cf9d977207/grh-compact/src/New/RealPlace/VolumeKernel.lean\#L135-L135}{{\leanfont frameX₀Part}}}{VolumeKernel.lean}{135--135}
\crow{frameVolumepos}{\href{https://github.com/lsernau/grh-lean/blob/72fead102f663d5c7a407964a13dc0cf9d977207/grh-compact/src/New/RealPlace/VolumeKernel.lean\#L138-L144}{{\leanfont frameVolume\_pos}}}{VolumeKernel.lean}{138--144}
\crow{frameVolumemulleft}{\href{https://github.com/lsernau/grh-lean/blob/72fead102f663d5c7a407964a13dc0cf9d977207/grh-compact/src/New/RealPlace/VolumeKernel.lean\#L154-L156}{{\leanfont frameVolume\_mul\_left}}}{VolumeKernel.lean}{154--156}
\crow{continuousframeVolume}{\href{https://github.com/lsernau/grh-lean/blob/72fead102f663d5c7a407964a13dc0cf9d977207/grh-compact/src/New/RealPlace/VolumeKernel.lean\#L158-L158}{{\leanfont continuous\_frameVolume}}}{VolumeKernel.lean}{158--158}
\crow{existskernelframeVolumele}{\href{https://github.com/lsernau/grh-lean/blob/72fead102f663d5c7a407964a13dc0cf9d977207/grh-compact/src/New/RealPlace/VolumeKernel.lean\#L164-L235}{{\leanfont exists\_kernel\_frameVolume\_le}}}{VolumeKernel.lean}{164--235}
\crow{frameXPartunipotentOfinvmul}{\href{https://github.com/lsernau/grh-lean/blob/72fead102f663d5c7a407964a13dc0cf9d977207/grh-compact/src/New/RealPlace/VolumeKernel.lean\#L237-L239}{{\leanfont frameX₀Part\_unipotentOf\_inv\_mul}}}{VolumeKernel.lean}{237--239}}
% the notations of the defined functions: one slot for each explicit argument of the Lean declaration
\newcommand{\RealPlaceSizes}[4]{\operatorname{RPS}(#1,#2,#3,#4)}
\newcommand{\haarRatio}[1]{\lambda(#1)}
\newcommand{\gauge}[2]{\Gamma(#1;#2)}
\newcommand{\gramSection}[1]{\gamma(#1)}
\newcommand{\invCorner}[1]{\operatorname{ic}(#1)}
\newcommand{\cornerDefect}[1]{\operatorname{cd}(#1)}
\newcommand{\IsBoundedInCell}[2]{#1\in\operatorname{Cell}_{\le #2}}
\newcommand{\diagLast}[1]{\operatorname{d}_{\varrho}(#1)}
\newcommand{\sectionCoeff}[2]{\mathsf{c}(#1,#2)}
\newcommand{\sectionFrame}[3]{\mathsf{f}(#1,#2,#3)}
\newcommand{\sectionMatrix}[3]{\mathsf{m}(#1,#2,#3)}
\newcommand{\frameVolume}[1]{\Delta(#1)}
\newcommand{\frameXPart}[1]{\mathsf{x}_0(#1)}
\newcommand{\UnitMembers}{239}
\newcommand{\UnitInNote}{79}
\newcommand{\UnitProofs}{108}
\newcommand{\UnitSteps}{639}
\newcommand{\UnitAllSteps}{665}


\linespread{1.08}                       % 11 pt on 14.7 pt
\setlength{\parindent}{0pt}
\setlength{\parskip}{5.5pt plus 1.5pt minus 0.5pt}
\setlength{\abovedisplayskip}{7.5pt plus 2pt minus 3pt}
\setlength{\belowdisplayskip}{7.5pt plus 2pt minus 3pt}
\setlength{\abovedisplayshortskip}{3pt plus 2pt}
\setlength{\belowdisplayshortskip}{6pt plus 2pt minus 2pt}
\emergencystretch=1.5em
\clubpenalty=10000 \widowpenalty=10000 \displaywidowpenalty=10000
\relpenalty=7000 \binoppenalty=9000
\raggedbottom
\pagestyle{fancy}\fancyhf{}\renewcommand{\headrulewidth}{0pt}
\fancyfoot[C]{\sffamily\fontsize{7.6}{9}\selectfont\color{inkfaint}\thepage\,/\,\pageref*{LastPage}}
\setlist{topsep=2pt,partopsep=0pt,parsep=0pt,itemsep=4pt,leftmargin=2.1em,labelsep=0.55em}

\newcommand{\Q}{\mathbb{Q}}
\newcommand{\Z}{\mathbb{Z}}
\newcommand{\F}{\mathbb{F}}
\newcommand{\len}{\operatorname{len}}
\newcommand{\cok}{\operatorname{cok}}
\newcommand{\End}{\operatorname{End}}
\newcommand{\Hom}{\operatorname{Hom}}
\newcommand{\Gal}{\operatorname{Gal}}
\newcommand{\diag}{\operatorname{diag}}
\newcommand{\kw}[1]{\emph{#1}}
\newcommand{\m}[1]{\mbox{$#1$}}        % a short formula that must not be broken across lines
\newcommand{\head}[1]{\par\addvspace{17pt plus 4pt minus 2pt}\Needspace{6\baselineskip}%
  \noindent{\fontsize{13.2}{16}\selectfont\bfseries\color{ink}#1}\par\nopagebreak\vspace{-1pt}}
\newcommand{\sect}[2]{\refstepcounter{section}\label{#2}\par\addvspace{19pt plus 4pt minus 2pt}\Needspace{7\baselineskip}%
  \noindent\llap{{\fontsize{13.2}{16}\selectfont\bfseries\color{inkfaint}\thesection}\hspace{9pt}}%
  {\fontsize{13.2}{16}\selectfont\bfseries\color{ink}#1}\par\nopagebreak\vspace{-1pt}}
\newcommand{\subhead}[1]{\par\addvspace{9pt plus 2pt}\Needspace{4\baselineskip}\noindent\textbf{#1}\enspace\ignorespaces}
% where a statement lives in the Lean archive
\newcommand{\leantag}{{\sffamily\bfseries\fontsize{6.3}{8}\selectfont\color{leancol}\addfontfeatures{LetterSpace=7}LEAN}}
\newcommand{\sep}{\unskip\nobreak\hspace{5pt}{\color{inkfaint}\textperiodcentered}\hspace{5pt plus 2pt}\ignorespaces}
\newcommand{\leanline}[1]{\par\nopagebreak\vspace{-2.5pt}{\fontsize{8.9}{12.4}\selectfont\color{inksoft}\raggedright
  \noindent\llap{\leantag\hspace{9pt}}#1\par}\vspace{1pt}}
% a displayed statement in a tinted box: #1 colour, #2 text
\newcommand{\statement}[2]{\par\addvspace{7pt}\begingroup\setlength{\fboxsep}{8.5pt}\setlength{\fboxrule}{0.6pt}%
  \noindent\fcolorbox{#1!30!white}{#1!3!white}{\parbox{\dimexpr\linewidth-2\fboxsep-2\fboxrule\relax}{#2}}%
  \endgroup\par\addvspace{7pt}}
\newcommand{\nt}[2]{\par\noindent\hangindent=1.1em\hangafter=1 #1\enspace{\color{inksoft}#2}\par\vspace{1pt}}
\newcommand{\legdot}[1]{\tikz[baseline=-0.62ex]\fill[#1] (0,0) circle (2.5pt);\hspace{3pt}}
\newcommand{\crow}[4]{\csname CS@#1\endcsname & #2 & \textcolor{inksoft}{#3} & \textcolor{inksoft}{#4}\\}
\newcommand{\leanquoteline}[2]{\par\noindent\hangindent=\dimexpr#1\fontdimen2\font+2.2em\relax\hangafter=1
  \hspace*{#1\fontdimen2\font}#2}
\newcommand{\quoted}[1]{\par\vspace{3pt}{\ttfamily\fontsize{7.5}{10.4}\selectfont\color{leancol}\raggedright
  \setlength{\parskip}{0pt}\leanquote{#1}\par}\vspace{4pt}}
\newcommand{\bx}[3]{{\sffamily\bfseries\fontsize{7.3}{9.2}\selectfont\color{#1!82!black}#2}\\[1.5pt]#3}
\newcommand{\cs}[2]{\expandafter\newcommand\csname CS@#1\endcsname{#2}}

% theorems
\newtheoremstyle{plainx}{7pt plus 2pt}{7pt plus 2pt}{\itshape}{}{\bfseries}{.}{0.5em}{}
\newtheoremstyle{uprightx}{7pt plus 2pt}{7pt plus 2pt}{}{}{\bfseries}{.}{0.5em}{}
\theoremstyle{plainx}
\newtheorem{theorem}{Theorem}[section]
\newtheorem{lemma}[theorem]{Lemma}
\newtheorem{proposition}[theorem]{Proposition}
\newtheorem{corollary}[theorem]{Corollary}
\theoremstyle{uprightx}
\newtheorem{definition}[theorem]{Definition}
\newtheorem{setting}[theorem]{Setting}          % hypotheses shared by several statements; each of them cites it by number
\newtheorem{remark}[theorem]{Remark}
\newtheorem{example}[theorem]{Example}
\renewcommand{\qedsymbol}{{\setlength{\fboxsep}{0pt}\setlength{\fboxrule}{0.5pt}\fbox{\rule{0pt}{0.46em}\hspace{0.46em}}}}
% what holds throughout the note: the ambient objects only (implicit and instance arguments), never a hypothesis
\newenvironment{conventions}{\par\addvspace{5pt}\noindent\textbf{Conventions.}\enspace\ignorespaces}{\par\addvspace{3pt}}
\makeatletter
\renewenvironment{proof}[1][Proof]{\par\pushQED{\qed}\normalfont\topsep0pt\partopsep0pt
  \trivlist\item[\hskip\labelsep\itshape #1.]\ignorespaces}{\popQED\endtrivlist\@endpefalse}
\makeatother

% notations of definitions of unit 132 that this note uses and another note of the unit states
% (made by shared/mkspec.py from shared/notation-132.json; the same TeX as in the note that states them)
\newcommand{\sieveVector}[2]{\Xi(#1;#2)}
\newcommand{\IsCutBessel}[7]{\operatorname{CB}(#1,#2,#3,#4,#5,#6,#7)}
\newcommand{\formConst}[1]{c_{\mathrm{f}}(#1)}
\newcommand{\HasFibreBound}[5]{\operatorname{Fib}_{#2,#4}(#1;#3)\le #5}
\newcommand{\Levels}[3]{\operatorname{Lev}(#1,#2,#3)}
\newcommand{\floorMeasure}[2]{\theta_{#2}(#1)}
\newcommand{\stackMeasure}[1]{\mu^{#1}}
\newcommand{\groupMeasure}[1]{\mathsf{M}(#1)}
\newcommand{\upperMeasure}[2]{\mathsf{M}_{\ge #2}(#1)}
\newcommand{\IsNear}[3]{#3\in\operatorname{Near}_{#2}(#1)}
\newcommand{\orthonormalPart}[1]{\mathsf{O}(#1)}
\newcommand{\closePairs}[1]{\operatorname{Cl}(#1)}
\newcommand{\floorValues}[4]{\operatorname{Val}_{#2}(#1;#3,#4)}
\newcommand{\volumeExponent}[1]{\mathsf{s}(#1)}
\newcommand{\floorWeight}[4]{\mathsf{w}_{#2}(#1;#3,#4)}
\newcommand{\floorModulus}[3]{\mathsf{d}_{#2}(#1;#3)}
\newcommand{\FloorMassLe}[4]{\operatorname{Mass}_{#2}(#1;#4)\le #3}
\newcommand{\upperPairing}[5]{\mathsf{Av}_{#2}(#1;#3)(#4,#5)}
\newcommand{\IsPairingKernel}[4]{(#3,#4)\in\operatorname{PK}_{#2}(#1)}
\newcommand{\invWeight}[3]{\mathsf{w}^{-1}_{#2}(#1;#3)}
\newcommand{\pivotModulus}[5]{\mathsf{Pv}_{#2}(#1;#3)(#4,#5)}
\newcommand{\weightedSquare}[5]{\mathsf{Sq}_{#2}(#1;#3,#4)(#5)}
\newcommand{\weightedSquareSum}[5]{\mathsf{Sm}_{#4}(#1;#2,#3)(#5)}
\newcommand{\selfNorm}[4]{\operatorname{SN}_{#2}(#1;#3,#4)}
\newcommand{\Gauges}[1]{\operatorname{Gau}(#1)}
\newcommand{\positionsOfSize}[2]{\operatorname{Pos}(#1;#2)}
\newcommand{\headDensity}[4]{\mathsf{h}_{#1}(#4;#2,#3)}
\newcommand{\translatePairing}[5]{\langle #1(#4),#2(#3#5)\rangle}
\newcommand{\IsBessel}[2]{#1\in\operatorname{Bes}(#2)}
\newcommand{\IsCut}[1]{#1\in\operatorname{Cut}}
\newcommand{\pairingNorm}[2]{\operatorname{PN}(#1,#2)}
\newcommand{\formIntegrandOne}[6]{\mathsf{I}^{1}_{#1}(#2;#3,#4,#5)(#6)}
\newcommand{\formIntegrand}[7]{\mathsf{I}_{#1}(#2,#3;#4,#5,#6)(#7)}
\newcommand{\floorForm}[5]{\mathcal{I}_{#1}(#2,#3;#4,#5)}
\newcommand{\formModulus}[6]{\mathcal{I}^{\mathrm{mod}}_{#1}(#2,#3;#4,#5,#6)}
\newcommand{\formMajorant}[6]{\mathcal{I}^{\mathrm{maj}}_{#1}(#2,#3;#4,#5,#6)}
\newcommand{\selfNormProduct}[7]{\operatorname{SNP}_{#1}(#7;#2,#3;#4,#5,#6)}
\newcommand{\Sizes}[1]{\operatorname{Siz}(#1)}
\newcommand{\weylRatio}[3]{w_{#1}(#3)^{-1}w_{#1}(#2)}
\newcommand{\vecsOfSize}[1]{\mathrm{Vec}_{#1}}
\newcommand{\HasMatSize}[2]{\operatorname{sz}(#2)\le #1}
\newcommand{\IsHead}[2]{#2\in\operatorname{Hd}_{#1}}
\newcommand{\floorConst}[2]{C_{#2}(#1)}
\newcommand{\globalConst}[1]{C_{\mathrm{gl}}(#1)}

% hand macros of the three notes of unit 132 (objects of the background library; letters of shared/PLAN-132.md)
\newcommand{\cL}{\mathcal{L}}
\newcommand{\CC}{\mathbb{C}}
\newcommand{\RR}{\mathbb{R}}
\newcommand{\NN}{\mathbb{N}}
\newcommand{\dd}{\mathrm{d}}
\newcommand{\GT}[1]{G(#1)}
\newcommand{\JT}[1]{J(#1)}
\newcommand{\Hl}[2]{\mathcal{H}_{#1}(#2)}
\newcommand{\Ql}[2]{\mathcal{Q}_{#1}(#2)}
\newcommand{\Nb}[2]{\mathcal{N}_{<#1}(#2)}
\newcommand{\Nf}[2]{\mathcal{N}_{#1}(#2)}
\newcommand{\UT}[1]{\mathcal{U}(#1)}
\newcommand{\modA}{\operatorname{mod}_A}
\newcommand{\ZT}[2]{Z_{#1}(#2)}
\newcommand{\sd}{\operatorname{sd}}
\newcommand{\wl}[2]{w_{#1}(#2)}
\newcommand{\wo}[2]{w^{\circ}_{#1}(#2)}
\newcommand{\ip}[2]{\langle #1,#2\rangle}
\newcommand{\gL}[1]{g_{#1}}
\newcommand{\BoxL}[1]{\operatorname{Box}(#1)}
\newcommand{\Posn}[1]{\operatorname{Posn}(#1)}
\newcommand{\Ball}[1]{\mathfrak{B}(#1)}
\newcommand{\etab}[1]{\eta_{<#1}}
\newcommand{\vol}{\operatorname{vol}}
\newcommand{\cLC}[1]{\mathcal{L}_{\CC}(#1)}
\newcommand{\Gone}[1]{G_1(#1)}
\newcommand{\Jone}[1]{J_1(#1)}
\newcommand{\volU}[1]{\vol\mathrm{U}(#1)}
\newcommand{\hw}[3]{\operatorname{hw}_{#1,#2}(#3)}
\newcommand{\damp}[2]{\ell_{#1}(#2)}
\newcommand{\bandc}[1]{\Theta(#1)}
\newcommand{\cellc}[1]{c_G(#1)}
\newcommand{\corr}[3]{J^{#1}_{#2}(#3)}
\newcommand{\boost}[2]{\chi_{#1}(#2)}
\newcommand{\quadric}[2]{\delta^{(#2)}(#1)}
\newcommand{\Inp}[1]{\textup{(#1)}}
\newcommand{\emb}[1]{\underline{#1}}
\newcommand{\uq}[1]{\mathrm{u}(#1)}
\newcommand{\muz}{\mu^{0}}
\newcommand{\bs}[1]{\operatorname{bs}(#1)}
\newcommand{\frbox}[1]{\operatorname{FB}(#1)}
\newcommand{\cay}[1]{\operatorname{cay}(#1)}
\newcommand{\cayd}[1]{\mathsf{j}(#1)}
\newcommand{\gd}[1]{\mathsf{g}(#1)}
\newcommand{\Ocell}[1]{\mathcal{C}(#1)}
\newcommand{\chart}{\operatorname{ch}}
\newcommand{\chartc}{\operatorname{cc}}
\newcommand{\btan}{\operatorname{bt}}
% where each declaration of the concordance is treated
\cs{invCorner}{Definition~\ref{d:cell}}
\cs{cornerDefect}{Definition~\ref{d:cell}}
\cs{openCell}{Definition~\ref{d:cell}}
\cs{IsBoundedInCell}{Definition~\ref{d:cell}}
\cs{gramDiscRealeq}{Definition~\ref{d:cell}}
\cs{memopenCelliff}{Definition~\ref{d:cell}}
\cs{continuousgramDiscReal}{Definition~\ref{d:cell}}
\cs{isOpenopenCell}{Definition~\ref{d:cell}}
\cs{cornerDefectpos}{Definition~\ref{d:cell}}
\cs{diagLast}{Definition~\ref{d:section}}
\cs{sectionCoeff}{Definition~\ref{d:section}}
\cs{sectionFrame}{Definition~\ref{d:section}}
\cs{sectionMatrix}{Definition~\ref{d:section}}
\cs{sqrteqsqrtcornerDefect}{Definition~\ref{d:section}}
\cs{diagLastsqrtmulinv}{Definition~\ref{d:section}}
\cs{sectionMatrixmulframeCols}{Definition~\ref{d:section}}
\cs{frameColsgramsectionMatrix}{Definition~\ref{d:section}}
\cs{gramSection}{Definition~\ref{d:gsec}}
\cs{coegramSection}{Definition~\ref{d:gsec}}
\cs{gramSectioneventuallyEq}{Definition~\ref{d:gsec}}
\cs{realFrameGramgramSection}{Definition~\ref{d:gsec}}
\cs{openCellChart}{Definition~\ref{d:chart}}
\cs{chartCoords}{Definition~\ref{d:chart}}
\cs{besselTangent}{Definition~\ref{d:chart}}
\cs{RealPlaceSizes}{Definition~\ref{d:rps}}
\cs{tenle}{Definition~\ref{d:rps}}
\cs{haarRatio}{Definition~\ref{d:ratio}}
\cs{gauge}{Definition~\ref{d:cgauge}}
\cs{frameVolume}{Definition~\ref{d:vol}}
\cs{frameXPart}{Definition~\ref{d:vol}}
\cs{continuousframeVolume}{Definition~\ref{d:vol}}
\cs{frameXPartunipotentOfinvmul}{Definition~\ref{d:vol}}
\cs{frameVectormul}{Lemma~\ref{l:fv}}
\cs{frameVectoreqsum}{Lemma~\ref{l:fv}}
\cs{normsumlemulnormsumframeVector}{Lemma~\ref{l:fv}}
\cs{continuousframeVector}{Lemma~\ref{l:fv}}
\cs{normframeVectorle}{Lemma~\ref{l:fv}}
\cs{existsnormleofbesselSizele}{Lemma~\ref{l:bsnorm}}
\cs{besselSizelenormmul}{Lemma~\ref{l:bsnorm}}
\cs{continuousbesselSize}{Lemma~\ref{l:bsnorm}}
\cs{besselSizele}{Lemma~\ref{l:bsnorm}}
\cs{besselSizebesselSmul}{Lemma~\ref{l:bsnorm}}
\cs{normleofmemframeBox}{Lemma~\ref{l:framebox}}
\cs{norminclMatrixleone}{Lemma~\ref{l:framebox}}
\cs{normJ}{Lemma~\ref{l:framebox}}
\cs{norminnerJle}{Lemma~\ref{l:framebox}}
\cs{realFrameGrameqinner}{Lemma~\ref{l:rfg}}
\cs{normrealFrameGramlemul}{Lemma~\ref{l:rfg}}
\cs{realFrameGrammulinvariant}{Lemma~\ref{l:rfg}}
\cs{realFrameGrameqframeMatrixmul}{Lemma~\ref{l:rfg}}
\cs{besselSizeleofdual}{Lemma~\ref{l:dual}}
\cs{measurablegramSection}{Lemma~\ref{l:gsec}}
\cs{differentiableAtsectionMatrixuncurry}{Lemma~\ref{l:secnorm}}
\cs{differentiableAtsectionMatrix}{Lemma~\ref{l:secnorm}}
\cs{normsectionMatrixle}{Lemma~\ref{l:secnorm}}
\cs{existseqmulgramSectionmul}{Lemma~\ref{l:double}}
\cs{memopenCellandboundsofnear}{Lemma~\ref{l:near}}
\cs{logDerivUunipotentOfmulembedH}{Lemma~\ref{l:logd}}
\cs{existssymmeqchartCoordsskewOfReal}{Lemma~\ref{l:coords}}
\cs{orthonormal}{Lemma~\ref{l:orth}}
\cs{euclideanenormsqeqsum}{Lemma~\ref{l:orth}}
\cs{lintegralmulleofframeMatrixle}{Lemma~\ref{l:fmle}}
\cs{lintegralgramSectionintegralle}{Lemma~\ref{l:sob}}
\cs{existsorthonormalkerlintegralle}{Lemma~\ref{l:ker}}
\cs{frameVolumepos}{Lemma~\ref{l:volpos}}
\cs{frameVolumemulleft}{Lemma~\ref{l:volpos}}
\cs{tsuminvindexpowleofcongr}{Lemma~\ref{l:count}}
\cs{middlesharedvectorle}{Lemma~\ref{l:frames}}
\cs{coleqzeroofnormltone}{Lemma~\ref{l:frames}}
\cs{eqofframeImageeq}{Lemma~\ref{l:frames}}
\cs{existsframeImagetoRowseqadd}{Lemma~\ref{l:translate}}
\cs{openCellChartjacobian}{Theorem A (i)}
\cs{openCellChartinjOn}{Theorem A (ii)}
\cs{lelintegralofgramSection}{Theorem A (iii)}
\cs{haarRatiospec}{Theorem B}
\cs{opencell}{Theorem C}
\cs{hasSchurBoundofconvolveeq}{Theorem D (i)}
\cs{existscutisCutPair}{Theorem D (ii)}
\cs{existskernelframeVolumele}{Theorem E}

\begin{document}
\color{ink}

{\raggedright
{\sffamily\bfseries\fontsize{7.2}{9}\selectfont\color{leancol}\addfontfeatures{LetterSpace=8}READBACK NOTE 30\par}
\vspace{5pt}
{\fontsize{21.5}{25.5}\selectfont\bfseries The tower method at a place: floor measures, gauges and the bounds on each floor\par}
\vspace{5pt}
{\fontsize{12.4}{16}\selectfont\itshape\color{inksoft} One unit of the Lean artifact \ArtifactName, read back into
mathematics\par}}
\vspace{8pt}
{\color{rulegrey}\hrule height 0.6pt}
\vspace{7pt}
\thispagestyle{empty}

\textbf{Summary.} At the complex place the unitary group $\Gone{\mathfrak{s}}$ has two descriptions that this note
compares. The \emph{tower} of the companion Note~T writes an element by floors; the \emph{open-cell chart} writes
it as $u\,\emb{k}\,\gramSection{Z}\,(v\,\emb{h})^{-1}$, with $Z$ the frame Gram matrix, $\gramSection{Z}$ an explicit section, $u,v$ in the
unipotent group that fixes the frame and $k,h$ in the compact group $\mathrm{U}(\kappa)$. The Jacobian of the chart
is $\gd{Z}^{k-1}$ times two Cayley weights, and the chart is injective; hence the tower measure is a multiple
$\haarRatio{\mathfrak{h}}$ of Haar measure with $\haarRatio{\mathfrak{h}}\le\volU{r}/\volU{k}$. The same chart gives a mean-square bound for
products of Bessel functions on the open cell. The note also normalises and cuts off a peaked Bessel function,
with Schur bounds for its derivatives, and bounds the volume weight $\Delta$ by a kernel of Hua type.

\statement{leancol}{\fontsize{9.6}{13.2}\selectfont\textbf{What this note is.} The artifact \ArtifactName{} is a
body of Lean code whose final theorems are the Riemann hypothesis and the generalized Riemann hypothesis for Dirichlet
$L$-functions. Its authors report that it compiles and that Lean's kernel accepts these theorems from the three
standard axioms; nothing was built for this note. That is a statement about a formal text: no mathematician has read
the argument through, and this note does not say that the Riemann hypothesis is proved. The note treats one part of
one unit of the artifact (the folder of the analysis at the real places): \UnitInNote{} of the unit's \UnitMembers{}
declarations, those of ten files; two companion notes treat the others (Note~T, \emph{The tower method at a place:
floor measures, gauges and the bounds on each floor}; Note~L, \emph{The local theorem for floor forms and the band
limits at the complex place}). Every result of the ten files is stated here and was compared, argument by argument,
with its Lean text. The proofs are complete for the lemmas on norms, on the section and on the volume weight; for
Theorems A to~E they follow the formal proofs, which were read line by line, and give every step that carries an
idea, while long matrix identities and comparisons of exponents are named and left to the formal text. Classical
results of the background library are quoted (Section~\ref{s:inputs}). The Lean text read is \PinText; Appendix~A is
the concordance of \ConcDecls{} declarations.}

\sect{The statements}{s:intro}
Notation is fixed in Section~\ref{s:setting}: $\Gone{\mathfrak{s}}$ is the unitary group of the hermitian space
$V=\CC^{r}$ with $\varrho$ hyperbolic planes, a positive part $\langle x_i\rangle_{i\in\kappa}$ and a negative
vector~$z$; $Z(y)$ is the frame Gram matrix of $y$; $\gd{Z}=|\det Z|^2-|D_\varrho(Z)|^2$, with $D_\varrho$ the leading
principal minor of order $\varrho$; $\Ocell{\mathfrak{s}}$ is the open cell $\{\gd{Z}>0\}$; $\mathcal{U}^0$ is the
group of the unipotent elements $\uq{q}$ that fix the frame, with its Haar measure $\muz$; $\emb{h}$ is the image of
$h\in\mathrm{U}(\kappa)$.

\statement{arith}{\phantomsection\label{t:A}\textbf{Theorem A (the open-cell chart).}\enspace Let $\mathfrak{s}$ be
a shape, and let $\chart$ be the map that sends a parameter $p=(Z,\theta,q,\eta,q')$ to
$\uq{q}\,\emb{\cay{\tilde\theta}}\,\gramSection{Z}\,\bigl(\uq{q'}\,\emb{\cay{\eta}}\bigr)^{-1}\in\Gone{\mathfrak{s}}$
(Definition~\ref{d:chart}).

\textbf{(i)} There is a real-linear isomorphism $e$ from the space of the parameters onto the space of the real
$r\times r$ matrices that preserves Lebesgue measure and such that, at every $p$ with $Z\in\Ocell{\mathfrak{s}}$,
the chart is differentiable and its logarithmic derivative $\mathrm{D}^{\log}\chart(p)$ satisfies
\[\bigl|\det\bigl(e^{-1}\circ\mathrm{sk}\circ\mathrm{D}^{\log}\chart(p)\bigr)\bigr|=\gd{Z}^{\,k-1}\cdot\cayd{\tilde\theta}\cdot\cayd{\eta}.\]

\textbf{(ii)} The chart is injective on the set of the parameters with $Z\in\Ocell{\mathfrak{s}}$.

\textbf{(iii)} Let $f\colon\Gone{\mathfrak{s}}\to[0,\infty]$ be measurable and $L\in[0,\infty]$. If for every
$k\in\mathrm{U}(\kappa)$
\[L\le\int_{\Ocell{\mathfrak{s}}}\cellc{\mathfrak{s}}\,\gd{Z}^{k-1}\int_{\mathrm{U}(\kappa)}\iint
f\bigl(u\,\emb{k}\,\gramSection{Z}\,(v\,\emb{h})^{-1}\bigr)\,\dd\muz(u)\,\dd\muz(v)\,\dd h\,\dd Z,\]
then $L\le\int_{\Gone{\mathfrak{s}}}f(y)\,\dd y$.}
\leanline{\lean{openCellChartjacobian}\sep\lean{openCellChartinjOn}\sep\lean{lelintegralofgramSection}}

\statement{arith}{\phantomsection\label{t:B}\textbf{Theorem B (the tower measure at the complex place).}\enspace Let
$\mathfrak{s}$ be a shape, $\rho$ and $b$ real numbers, $s_0$ a natural number, and $\mathfrak{h}$ real-place sizes
in $\RealPlaceSizes{\mathfrak{s}}{\rho}{s_0}{b}$ (Definition~\ref{d:rps}), with their level $\cLC{\mathfrak{h}}$
(Definition~\ref{d:clevel}). Then the constant $\haarRatio{\mathfrak{h}}$ of Definition~\ref{d:ratio} satisfies:
\begin{enumerate}[label=\textup{(\roman*)}]
\item $\haarRatio{\mathfrak{h}}\ne0$ and $\haarRatio{\mathfrak{h}}\ne\infty$;
\item $\groupMeasure{\cLC{\mathfrak{h}}}=\haarRatio{\mathfrak{h}}\cdot\dd y$, the Haar measure of
  $\Gone{\mathfrak{s}}$; in particular the tower measure is a Haar measure;
\item for every measurable $f\colon\Gone{\mathfrak{s}}\to[0,\infty]$ that is invariant under left translations by
  $\UT{T}$: $\int f\,\dd\upperMeasure{\cLC{\mathfrak{h}}}{0}=\haarRatio{\mathfrak{h}}\int f\,\dd\dot y$;
\item $\dfrac{\volU{k}}{\volU{r}}\cdot\haarRatio{\mathfrak{h}}\le1$.
\end{enumerate}}
\leanline{\lean{haarRatiospec}}

\statement{arith}{\phantomsection\label{t:C}\textbf{Theorem C (the open-cell bound).}\enspace Let $\mathfrak{s}$ be a
shape, $B$ a real number, $f$ a natural number and $X$ a Bessel system for $(\mathfrak{s},B,f)$
(Section~\ref{s:inputs}), with full function $\Phi_0\colon\Gone{\mathfrak{s}}\to\CC^I$ and formal degree $\deg$. Let
$\gamma,y\in\Gone{\mathfrak{s}}$ and $\rho$ real with $\gd{Z(\gamma)}>0$, $y\in\frbox{\rho}$,
$\gamma^{-1}y\in\frbox{\rho}$ and $B\le\rho$. Then there are a compact set $A\subseteq\Gone{\mathfrak{s}}$, a
finite measure $\nu$ on $\Gone{\mathfrak{s}}$, a probability measure $m$ on $\mathrm{U}(\kappa)\times\mathcal{U}^0$
and a measurable map $t\colon\mathrm{U}(\kappa)\times\mathcal{U}^0\to\Gone{\mathfrak{s}}$ such that:
\begin{enumerate}[label=\textup{(\roman*)}]
\item for every measurable bounded $\varphi\colon\Gone{\mathfrak{s}}\to\CC$,
\[\int\sum_{i,i'\in I}\Bigl|\int_A\varphi(y')\,\Phi_{0,i}(\gamma\,t(x)\,y')\,\overline{\Phi_{0,i'}(t(x)\,y')}\,\dd y'\Bigr|^{2}\dd m(x)\le\int|\varphi|^2\,\dd\nu;\]
\item $\nu(\Gone{\mathfrak{s}})\le\Bigl(\dfrac{\rho^{9n}}{\gd{Z(\gamma)}}\Bigr)^{10^{5}n^{4}}\cdot
  \dfrac{|I|^{2}}{\deg\cdot\cellc{\mathfrak{s}}\cdot(\gd{Z(\gamma)}/2)^{k-1}}$;
\item for $m$-almost every $x$: $t(x)$ lies in the group generated by $\mathcal{U}^0$ and $\Hl{n}{T}$, and every
  $y'\in\frbox{\rho}$ with $\bs{\gamma\,t(x)\,y'}\le\rho$ lies in $A$.
\end{enumerate}}
\leanline{\lean{opencell}}

\statement{arith}{\phantomsection\label{t:D}\textbf{Theorem D (the Bessel cut).}\enspace Let $\mathfrak{s}$ be a
shape, $I$ a finite set, $\varepsilon$ a family of signs with a frame $\mathsf{p}$ of $\Jone{\mathfrak{s}}$
(Section~\ref{s:inputs}), $\mathrm{rep}$ a homomorphism from $\mathrm{U}(\kappa)$ to the unitary group of $\CC^I$,
$A_0$ and $L$ natural numbers, and $l$ a parameter with a large gap $A_0$ and with $|l_a|\le L$ for all $a$; write
$\deg_l$ for its formal degree and $\psi_l$ for its coefficient in the frame.

\textbf{(i)} Let $\Phi_0\colon\Gone{\mathfrak{s}}\to\CC^I$ be a Bessel function for $\mathrm{rep}$ that is integrable
for $\dd\dot y$, and $B$ a real number. If $\Phi_0*(\deg_l\psi_l)=\Phi_0$ and
$(r+1)^{4}(L+r+1+2n^{2})\le B$, then $\Phi_0$ has the Schur bounds with respect to $\mathrm{rep}$, $B$ and the
constant $\deg_l^{-1}\int|\Phi_0|^2\,\dd\dot y$.

\textbf{(ii)} Let $f\ne0$ and $N$ be natural numbers, $b$ a real number, $x\in\Gone{\mathfrak{s}}$ and
$Y\colon\Gone{\mathfrak{s}}\to\CC^I$ with a peak at $x$ of constant $b^{N}$ (for $l$ and $\mathrm{rep}$). Suppose
$\bs{x}\le b^{5n}$, $2000\le b$, $N+99r^{2}\le r^{3}$, $r+L+1\le b^{4}$ and $3r^{3}f\le A_0$. Then there are
continuous $\Phi_0,\Phi\colon\Gone{\mathfrak{s}}\to\CC^I$ with $\Phi_0*(\deg_l\psi_l)=\Phi_0$ such that
$(\Phi_0,\Phi)$ is a cut pair with respect to $b$, $f$ and $\mathrm{rep}$.}
\leanline{\lean{hasSchurBoundofconvolveeq}\sep\lean{existscutisCutPair}}

\statement{arith}{\phantomsection\label{t:E}\textbf{Theorem E (a kernel for the volume weight).}\enspace Let
$\mathfrak{s}$ be a shape, $y\in\Gone{\mathfrak{s}}$ and $\theta$ a real number with $\bs{y}\le\theta$. Let $j$ be a
natural number with $j\le n$, $\sigma$ a finite set with an injection $e\colon\sigma\to\{0,\dots,j-1\}$, $R$ a
complex $\sigma\times\sigma$ matrix with $\det R\ne0$, $Z$ a complex $n\times n$ matrix whose leading $j\times j$
block $Z_j$ has $\det Z_j\ne0$, and $V$ a complex $\kappa\times j$ matrix. Then there is a measurable function
$K$ on the complex $\kappa\times\sigma$ matrices, with values in $[0,\infty]$, such that
\[\int K(a)\,\dd a\ \le\ \Bigl(\frac{(2\theta^{2})^{j}}{|\det Z_j|}\Bigr)^{k+n}\cdot
\Bigl(\frac{\pi^{|\sigma|}}{|\sigma|!\,\bigl(\theta^{|\sigma|}\,|\det R|\bigr)^{2}}\Bigr)^{k},\]
and such that for all $\gamma\in\Gone{\mathfrak{s}}$ and all complex $\kappa\times\sigma$ matrices $A$ with
$Z(\gamma)=Z$, with $\gamma$ in normal form at $j$, and with
$\frameXPart{\gamma}\,\mathsf{I}_j=V+A\,R\,\mathsf{E}^{*}Z_j$,
\[\Bigl(\frac{\frameVolume{y}}{\frameVolume{\gamma^{-1}y}}\Bigr)^{k+n}\le K(A).\]
Here $\mathsf{I}_j$ is the $n\times j$ matrix of the inclusion of the first $j$ frame positions and $\mathsf{E}$ the
$j\times\sigma$ matrix of $e$.}
\leanline{\lean{existskernelframeVolumele}}

\sect{The setting}{s:setting}
\begin{conventions}\phantomsection\label{c:conventions}
$\mathfrak{s}=(\varrho,\kappa,k_0)$ is a shape; $n=\varrho+1$, $k=|\kappa|$, $r=2\varrho+k+1$. $I$ is a finite
set~$I$ and $\CC^I$ has its standard inner product $\ip{v}{w}=\sum_i\overline{v_i}w_i$ and norm $|v|$. Matrices are
complex unless said otherwise; $X^{*}$ is the conjugate transpose and $\|X\|$ the operator norm for the Euclidean
norms. Finite index sets carry the counting structure, spaces of real or complex matrices their Lebesgue measure.
In Section~\ref{s:vol} the letters $E$, $\mathsf{S}$ denote a CM field $E$ and a filling $\mathsf{S}$ of the unit on
the second moment. Integrals of functions with values in $[0,\infty]$ are Lebesgue integrals; subtraction of
natural numbers is truncated at~$0$.
\end{conventions}

\subhead{The group at the complex place (background library, units 112, 115, 118, 119).}\phantomsection\label{s:space}
The space $V=\CC^{r}$ has the basis $e_0,\dots,e_{\varrho-1}$, $x_i$ ($i\in\kappa$), $z$,
$e^*_{\varrho-1},\dots,e^*_0$, and the hermitian form $h(v,w)=v^{*}\Jone{\mathfrak{s}}w$ whose Gram matrix
$\Jone{\mathfrak{s}}$ pairs $e_a$ with $e^*_a$, is the identity on the $x_i$ and $-1$ at $z$;
$\Jone{\mathfrak{s}}^{2}=1$. $\Gone{\mathfrak{s}}$ is its unitary group, with a Haar measure $\dd y$. In the
notation of Note~T this is $\GT{T}$ for the place $\CC$ and the form $T$ of $\Jone{\mathfrak{s}}$, and we use the
subgroups $\Hl{l}{T}$, $\Ql{l}{T}$, $\UT{T}$ of that note; $\Hl{n}{T}$ is the image of $\mathrm{U}(\kappa)$, the
unitary group of $\langle x_i\rangle$, and $\emb{h}$ denotes the image of $h$. The \emph{frame} is
$f_0,\dots,f_\varrho$ with $f_a=e_a$, $f_\varrho=z$; $F_{\mathfrak{s}}$ is the $r\times n$ matrix with these
columns, and $F'_{\mathfrak{s}}$ the one with the columns $e^*_0,\dots,e^*_{\varrho-1},x_{k_0}$. For
$y\in\Gone{\mathfrak{s}}$, the \emph{frame vectors} are $y^{-1}f_a\in\CC^{r}$, the \emph{frame Gram matrix} is
$Z(y)=(h(f_a,yf_b))_{a,b}=F_{\mathfrak{s}}^{*}\Jone{\mathfrak{s}}\,yF_{\mathfrak{s}}$, and the \emph{frame matrix}
$\mathrm{fm}(y)$ is the block of $y$ on the frame. For a complex $n\times n$ matrix $Z$,
$\gd{Z}=|\det Z|^{2}-|D_\varrho(Z)|^{2}$. The group $\mathcal{U}^0$ of the elements of $\UT{T}$ that fix every
frame vector consists of the elements $\uq{q}$, $q=(q_1,q_2)$ a complex $\varrho\times\kappa$ matrix and a real
$\varrho\times\varrho$ matrix, and $\muz$ is the image of Lebesgue measure under $q\mapsto\uq{q}$, a Haar measure
of $\mathcal{U}^0$. The group generated by $\mathcal{U}^0$ and $\Hl{n}{T}$ is the stabiliser of the frame; it lies
in $\Ql{n}{T}$, which normalises $\mathcal{U}^0$ and $\UT{T}$, and $\Ql{n}{T}=\UT{T}\cdot\Hl{n}{T}$. The measure
$\dd\dot y$ on $\Gone{\mathfrak{s}}$ computes integrals over $\UT{T}\backslash\Gone{\mathfrak{s}}$:
$\int\!\!\int_{\UT{T}}f(uy)\,\dd u\,\dd\dot y=\int f\,\dd y$ for a fixed Haar measure of $\UT{T}$. The \emph{Bessel
size} of $y$ is $\bs{y}=\inf_{s\in\Ql{n}{T}}(\|sy\|+1)$, and $\frbox{\rho}$ is the set of the $y$ with
$|y^{-1}f_a|\le\rho$ for all $a\le\varrho$. For a real $\kappa\times\kappa$ matrix $\eta$, $\cay{\eta}\in\mathrm{U}(\kappa)$
is the \emph{Cayley chart} $(1+X)(1-X)^{-1}$, $X=\tfrac12\operatorname{skew}(\eta)$, where $\operatorname{skew}$
identifies real matrices with skew-hermitian ones and $\mathrm{sk}$ is its left inverse on all complex matrices; and
$\cayd{\eta}=|\det(1-X)|^{-2k}$ is its \emph{Cayley weight}. A map $\Phi$ into $\Gone{\mathfrak{s}}$ that is
differentiable at $p$ has the \emph{logarithmic derivative}
$\mathrm{D}^{\log}\Phi(p)\colon v\mapsto\Phi(p)^{*}\Jone{\mathfrak{s}}\,\mathrm{d}\Phi_p(v)$.

\subhead{The section.}
\begin{definition}[corner, defect, open cell]\label{d:cell}
For a complex $n\times n$ matrix $Z$: $\invCorner{Z}$ is the last diagonal entry of $Z^{-1}$ (the inverse of a
singular matrix being $0$), and $\cornerDefect{Z}=1-|\invCorner{Z}|^{2}$. The \emph{open cell}
$\Ocell{\mathfrak{s}}$ is the set of the $Z$ with $\det Z\ne0$ and $|\invCorner{Z}|<1$. For a real $P$,
\m{\IsBoundedInCell{Z}{P}} means: $Z\in\Ocell{\mathfrak{s}}$, $\|Z^{-1}\|\le P$, $\|Z\|\le P$ and
$\cornerDefect{Z}^{-1/2}\le P$.
\end{definition}
\leanline{\lean{invCorner}\sep\lean{cornerDefect}\sep\lean{openCell}\sep\lean{IsBoundedInCell}\sep\lean{gramDiscRealeq}\sep\lean{memopenCelliff}\sep\lean{continuousgramDiscReal}\sep\lean{isOpenopenCell}\sep\lean{cornerDefectpos}}
If $\det Z\ne0$ then $\gd{Z}=|\det Z|^{2}\cornerDefect{Z}$, by Cramer's rule; so
$\Ocell{\mathfrak{s}}=\{Z:\gd{Z}>0\}$, an open set since $\gd{\cdot}$ is continuous, and
$\cornerDefect{Z}>0$ on it.
\begin{definition}[the matrices of the section]\label{d:section}
For a complex number $x$, $\diagLast{x}$ is the diagonal $n\times n$ matrix $\operatorname{diag}(1,\dots,1,x)$. For
complex $n\times n$ matrices $Z$, $W$ and a complex number $s$:
\begin{gather*}
\sectionCoeff{Z}{W}=\diagLast{0}-1-\diagLast{0}\,W^{-1}\bigl(1-\diagLast{\tfrac12}\bigr)Z^{-1}\diagLast{2},\\
\sectionFrame{Z}{W}{s}=\bigl(F_{\mathfrak{s}}\,\sectionCoeff{Z}{W}+F'_{\mathfrak{s}}\,\diagLast{s}\bigr)Z,\\
\sectionMatrix{Z}{W}{s}=1+\bigl(\sectionFrame{Z}{W}{s}-F_{\mathfrak{s}}\bigr)F_{\mathfrak{s}}^{*}
+\Bigl(\bigl(F_{\mathfrak{s}}W^{-1}-\sectionFrame{Z}{W}{s}\,\sectionCoeff{W}{Z}\bigr)\diagLast{s^{-1}}-F'_{\mathfrak{s}}\Bigr)F'^{*}_{\mathfrak{s}}.
\end{gather*}
In the Lean text the last two arguments have the default values $W=Z^{*}$ and
$s=\sqrt{1-\invCorner{Z}\invCorner{W}}$ (the principal square root), which is $\sqrt{\cornerDefect{Z}}$ for
$W=Z^{*}$ and $Z\in\Ocell{\mathfrak{s}}$.
\end{definition}
\leanline{\lean{diagLast}\sep\lean{sectionCoeff}\sep\lean{sectionFrame}\sep\lean{sectionMatrix}\sep\lean{sqrteqsqrtcornerDefect}\sep\lean{diagLastsqrtmulinv}\sep\lean{sectionMatrixmulframeCols}\sep\lean{frameColsgramsectionMatrix}}
Two identities follow by multiplying out, for $\det Z\ne0$:
$\sectionMatrix{Z}{W}{s}F_{\mathfrak{s}}=\sectionFrame{Z}{W}{s}$ and
$F_{\mathfrak{s}}^{*}\Jone{\mathfrak{s}}\,\sectionFrame{Z}{W}{s}=Z$.
\begin{definition}[the section]\label{d:gsec}
For a complex $n\times n$ matrix $Z$, $\gramSection{Z}$ is the element of $\Gone{\mathfrak{s}}$ with the matrix
$\sectionMatrix{Z}{Z^{*}}{\sqrt{\cornerDefect{Z}}}$ if $Z\in\Ocell{\mathfrak{s}}$, and the identity otherwise. That
this matrix is unitary for $\Jone{\mathfrak{s}}$ is part of the definition in the Lean text
(Lemma~\ref{l:gsec}).
\end{definition}
\leanline{\lean{gramSection}\sep\lean{coegramSection}\sep\lean{gramSectioneventuallyEq}\sep\lean{realFrameGramgramSection}}
So for $Z\in\Ocell{\mathfrak{s}}$ the matrix of $\gramSection{Z}$ is the section matrix, also near $Z$, and
$Z(\gramSection{Z})=Z$ by the two identities above.
\begin{definition}[the chart and its coordinates]\label{d:chart}
The \emph{parameters} are the tuples $p=(Z,\theta,q,\eta,q')$ of a complex $n\times n$ matrix $Z$, a pair $\theta$ of
a real number and a vector of $\CC^{\kappa\smallsetminus\{k_0\}}$, two parameters $q,q'$ of $\mathcal{U}^0$ and a
real $\kappa\times\kappa$ matrix $\eta$. Let $\tilde\theta$ be the real $\kappa\times\kappa$ matrix that has the
entry and the column of $\theta$ at $k_0$ and is $0$ elsewhere (the volume-preserving identification of the
background library). The \emph{open-cell chart} is
\[\chart(p)=\uq{q}\cdot\emb{\cay{\tilde\theta}}\cdot\gramSection{Z}\cdot\bigl(\uq{q'}\cdot\emb{\cay{\eta}}\bigr)^{-1}.\]
The \emph{chart coordinates} are the real-linear map $\chartc$ from complex $r\times r$ matrices $S$ to parameters
whose $Z$-component is $F'^{*}_{\mathfrak{s}}SF_{\mathfrak{s}}$ and whose other components are read from the blocks of
$\mathrm{sk}(S)$ and of $S$ at the indices of the $x_i$ and of the $e^*_a$. The \emph{Bessel tangent map} is the
real-linear map $\btan$ from triples $(\eta,q_1,q_2)$ of complex $\kappa\times\kappa$, complex $\varrho\times\kappa$
and real $\varrho\times\varrho$ matrices to complex $r\times r$ matrices,
\[\btan(\eta,q_1,q_2)=(W_{\mathfrak{s}}\eta+D_{\mathfrak{s}}q_1)W_{\mathfrak{s}}^{*}+\bigl(D_{\mathfrak{s}}\operatorname{skew}(q_2)-W_{\mathfrak{s}}q_1^{*}\bigr)D_{\mathfrak{s}}^{*},\]
with $W_{\mathfrak{s}}$, $D_{\mathfrak{s}}$ the matrices of the columns $x_i$ and $e^*_a$. None of the three takes
an explicit argument in the Lean text; the shape is implicit.
\end{definition}
\leanline{\lean{openCellChart}\sep\lean{chartCoords}\sep\lean{besselTangent}}

\subhead{Sizes, level and gauge at the complex place.}
\begin{definition}[real-place sizes]\label{d:rps}
For a shape $\mathfrak{s}$, real numbers $\rho$, $b$ and a natural number $s_0$, \emph{real-place sizes}
$\mathfrak{h}$, elements of $\RealPlaceSizes{\mathfrak{s}}{\rho}{s_0}{b}$, are proofs of the seven conditions
\[\rho=b^{6n+2},\quad 8r\le\rho,\quad k=s_0+3n+2,\quad 10^{9}n^{3}<s_0,\quad n^{7}\le k,\quad 10^{8}\le\varrho,\quad 2\le b.\]
\end{definition}
\leanline{\lean{RealPlaceSizes}\sep\lean{tenle}}
They imply $\rho\ge4$.
\begin{definition}[the level at the complex place]\label{d:clevel}
For real-place sizes $\mathfrak{h}$ in $\RealPlaceSizes{\mathfrak{s}}{\rho}{s_0}{b}$, the level
$\cLC{\mathfrak{h}}$ is the level in $\Levels{\kappa\smallsetminus\{k_0\}}{T}{0}$, for the place $\CC$ and the form
$T$ of $\Jone{\mathfrak{s}}$, with: the hyperbolic basis of $V_\varrho$ made of the pair
$\tfrac{1+\mathrm{i}}{2}(x_{k_0}\pm z)$ and the orthonormal vectors $x_i$, $i\ne k_0$; the measure on floor $n$ the
Haar probability measure of $\mathrm{U}(\kappa)$, carried to $\Hl{n}{T}$; the measure on floor $\varrho$ the
convolution of the image of $\frac{2\pi^{k}}{(k-1)!}\max(|c|^{2}-1,0)^{k-1}\dd c$ under
$c\mapsto\boost{\varrho}{c}^{-1}$ with the measure on floor $n$; the unit $1$; the ball $\{v:|v|<\tfrac14\}$; the box
$\{x:\|x\|\le\rho\}$; and all of $\Gone{\mathfrak{s}}$ as positions. Here $\boost{\varrho}{c}\in\Hl{\varrho}{T}$ is
the boost of the background library in the plane of $x_{k_0}$ and $z$. This definition is in the Lean text but not
in the census of the artifact (see Section~\ref{s:notes}); the axioms (L1)--(L3) of a level are verified in it from
results of unit 119.
\end{definition}
\begin{definition}[the Haar ratio]\label{d:ratio}
For real-place sizes $\mathfrak{h}$ in $\RealPlaceSizes{\mathfrak{s}}{\rho}{s_0}{b}$,
$\haarRatio{\mathfrak{h}}\in[0,\infty]$ is a chosen constant with
$\groupMeasure{\cLC{\mathfrak{h}}}=\haarRatio{\mathfrak{h}}\cdot\dd y$; one exists by Theorem~A(i) of Note~T and
the uniqueness of Haar measure (Input~\Inp{J2}).
\end{definition}
\leanline{\lean{haarRatio}}
\begin{definition}[the gauge at the complex place]\label{d:cgauge}
For real-place sizes $\mathfrak{h}$ in $\RealPlaceSizes{\mathfrak{s}}{\rho}{s_0}{b}$ and a real number $R$ with a
proof of $\rho\le R$, $\gauge{\mathfrak{h}}{\rho\le R}$ is the gauge in $\Gauges{\cLC{\mathfrak{h}}}$ with the sets
$\Gamma_i=\{x\in\CC:|x|\le R^{i}\}$ and the constant $\Lambda=R^{2}$. That these data satisfy the axioms of a gauge
is part of the definition in the Lean text (Proposition~\ref{p:gauge}).
\end{definition}
\leanline{\lean{gauge}}

\subhead{The volume weight.}
\begin{definition}[volume weight; the part of the frame in the positive space]\label{d:vol}
For $\gamma\in\Gone{\mathfrak{s}}$,
\[\frameVolume{\gamma}=\bigl|\det\bigl(((1-\Jone{\mathfrak{s}})F_{\mathfrak{s}})^{*}\,\gamma^{-1}F_{\mathfrak{s}}\bigr)\bigr|,
\qquad\frameXPart{\gamma}=W_{\mathfrak{s}}^{*}\,\gamma\,F_{\mathfrak{s}}.\]
The columns of $(1-\Jone{\mathfrak{s}})F_{\mathfrak{s}}$ are $e_a-e^*_a$ ($a<\varrho$) and $2z$, a basis of a
maximal negative subspace; so $\frameVolume{\gamma}$ is the absolute value of the determinant of the Euclidean
products of this basis with the frame vectors of $\gamma$, and $\frameXPart{\gamma}$ is the $\kappa\times n$ matrix
of the components of $\gamma f_0,\dots,\gamma f_\varrho$ on the $x_i$.
\end{definition}
\leanline{\lean{frameVolume}\sep\lean{frameXPart}\sep\lean{continuousframeVolume}\sep\lean{frameXPartunipotentOfinvmul}}
$\frameVolume{\cdot}$ is continuous; and $\frameXPart{\uq{q}^{-1}\gamma}=\frameXPart{\gamma}+q_1^{*}\cdot(\text{the
first $\varrho$ rows of }Z(\gamma))$, by the formula for $\uq{q}$.

\sect{What is taken for granted}{s:inputs}
\subhead{From Note T and Note L (same unit).} The notions of level, gauge, floor measures, tower and upper measures,
and Theorem~A of Note~T (the tower measure is right invariant and equals $\etab{j}$ convolved with floor measures).
Note~L is not used; it uses Theorems B, D and Lemma~\ref{l:bsnorm}.

\subhead{From the unit on the second moment (unit 148) and unit 133.} A \emph{cut pair} with respect to $b$, $f$,
$\mathrm{rep}$ is a pair of continuous functions $\Phi_0$ (full), $\Phi$ (cut) from $\Gone{\mathfrak{s}}$ to $\CC^I$,
both Bessel functions for $\mathrm{rep}$ (measurable, bounded, $F(uy)=\exp(\sd(u)-\overline{\sd(u)})F(y)$ for
$u\in\UT{T}$, $F(\emb{h}y)=\mathrm{rep}(h)F(y)$), with: $\Phi_0$ integrable for $\dd\dot y$;
$\int|\Phi_0|^2\dd\dot y=|I|$; $\int|\Phi|^2\dd\dot y\le|I|$; $\Phi(y)\ne0$ implies $\bs{y}\le b^{6n+2}$;
$\int|\Phi-\Phi_0|^2\dd\dot y\le b^{-10r^{2}}|I|$; and there are a measurable $\UT{T}$-invariant
$\mathsf{A}\ge0$ and a finite measure $\mu$ with $c(y)\,|F(y)|\le\mathsf{A}(y)$ for $F=\Phi_0,\Phi$, where $c(y)=1$
if $\Phi_0(y)=\Phi(y)$ and $c(y)=b^{3r^{3}f}$ otherwise, with $\mathsf{A}(y)\le\mu\{g:\bs{yg}\le b^{6n+2}\}$, and
with $\mu(\Gone{\mathfrak{s}})\int\mathsf{A}\,\dd\dot y\le b^{2r^{3}}|I|$. A function $\Phi_0$ \emph{has Schur's bound
with constant $D$} if $\int|\int\varphi(y)\Phi_0(\gamma y)\dd y|^{2}\dd\dot\gamma\le D\int|\varphi|^2$ for integrable
$\varphi$; it \emph{has the Schur bounds with respect to $\mathrm{rep}$, $B$, $D$} if there are continuous Bessel
functions $V_w$ for $\mathrm{rep}$, indexed by the words $w$ in pairs of basis indices, with $V_\emptyset=\Phi_0$,
with $V_{(\alpha,\beta)w}$ the right derivative of $V_w$ along the elementary matrix at $(\alpha,\beta)$, and with
Schur's bound with constant $B^{2|w|}D$ for $V_w$ when $|w|\le2n^{2}$. A \emph{Bessel system} for
$(\mathfrak{s},B,f)$ consists of a non-empty finite set $I$, a cut pair with respect to $B$, $f$ and an irreducible
$\mathrm{rep}$ (every matrix commuting with all $\mathrm{rep}(h)$ is scalar), a real $\deg>0$, the Schur bounds of
the full function with respect to every $P\ge B^{5}$ and $D=|I|/\deg$, the inequality $8r\le B$, and band limits.
A function $Y$ \emph{has a peak at $x$ of constant $P$} (for $l$, $\mathrm{rep}$) if it is a Bessel function for
$\mathrm{rep}$ with $Y*(\deg_l\psi_l)=Y$, $Y(x)\ne0$,
$|Y(y)|\le\int|\psi_l(x^{-1}sy)|\dd\nu_S(s)$ and $\int(2/\|x^{-1}sy\|)^{8r^{2}+1}\dd\nu_S(s)\le P\,|Y(x)|$ for all
$y$, where $\nu_S$ is the measure of unit 133 on a group of pairs over $\Ql{n}{T}$ with
$\int f\,\dd y=\iint f(x\,s\,y)\,\dd\nu_S(s)\,\dd\dot y$. Here $(F*e)(y)=\int e(a^{-1})F(ya)\,\dd a$. The filling
$\mathsf{S}$ of Section~\ref{s:vol}, its levels, labelled terms, positions and lattices are those of unit 148;
that section recalls what it uses.

\subhead{From the background library.} A \emph{frame} $\mathsf{p}$ for signs $\varepsilon$ is a basis with
$\mathsf{p}^{*}\Jone{\mathfrak{s}}\mathsf{p}=\operatorname{diag}(\varepsilon)$ and $\mathsf{p}$ unitary; a
\emph{parameter} $l$ is an integer vector, twice a Harish-Chandra parameter $\lambda$; it has a \emph{large gap
$A_0$} if $\lambda$ is regular with $|\lambda_i-\lambda_{i'}|\ge500r^{3}+A_0$ when $\varepsilon_i\ne\varepsilon_{i'}$;
$\psi_l(g)$ is the lowest $K$-type coefficient of the discrete series of $\mathrm{U}(\operatorname{diag}\varepsilon)$
at $\mathsf{p}^{-1}g\mathsf{p}$, and $\deg_l=(\int|\psi_l|^2)^{-1}$.

\Inp{J1} \emph{Witt's theorem} (unit 120): over a field with a non-trivial involution, if $X$, $Y$ have left
inverses and equal Gram matrices for an invertible hermitian $J$, then $gX=Y$ for some $g\in\mathrm{U}(J)$. The
second reader relates it to Witt's extension theorem (U.~First, \emph{Witt's extension theorem for quadratic spaces
over semiperfect rings}, Corollary~4.9) as the theorem itself.

\Inp{J2} \emph{Uniqueness of Haar measure} (unit 114): on a second countable locally compact group with a Haar
measure that is also right invariant, every $s$-finite right invariant measure is a multiple, in $[0,\infty]$, of
it. Reference: N.~Bourbaki, \emph{Int\'egration} VII; the second reader's comparison of the $s$-finite form with
the literature is unresolved, reference to be confirmed.

\Inp{J3} \emph{Unfolding} (unit 114): if $(u,g)\mapsto ug$ carries a Haar measure of $\UT{T}$ times a measure $\mu'$
to $\mu$, then $\int f\,\dd\mu'=\int f\,\dd\dot\mu$ for $\UT{T}$-invariant $f\ge0$, with $\dot\mu$ the quotient
measure of $\mu$. This is Weil's integration formula for a quotient; reference to be confirmed.

\Inp{J4} \emph{Haar measure in a chart} (unit 115): if a map $\Phi$ from an open set of parameters to
$\mathrm{U}(J)$, $J^2=1$, is injective and differentiable on a measurable set $A$ with
$|\det(e^{-1}\circ\mathrm{sk}\circ\mathrm{D}^{\log}\Phi(p))|=w(p)$ for a measure-preserving $e$, then for almost
every real $t$ and all measurable $f\ge0$:
$\int_A\frac{w(p)}{\volU{N}}f(\mathrm{e}^{\mathrm{i}t}\Phi(p))\,\dd p\le\int f\,\dd y$ for the Haar probability measure
(and correspondingly for $\dd y$). The Cayley chart of $\mathrm{U}(J)$ has
$|\det(\mathrm{sk}\circ\mathrm{D}^{\log}\operatorname{cay}(\eta))|=\cayd{\eta}$. The catalogue of the write-up names
L.~K.~Hua, \emph{Harmonic analysis of functions of several complex variables in the classical domains} (1963),
ch.~III; place not confirmed.

\Inp{J5} \emph{Bruhat coordinates} (unit 118):
\[\int f(Z)\,\dd Z=\int\prod_{a\le\varrho}|c_a|^{4(\varrho-a)}\iint f(x_-^{*}\operatorname{diag}(c)\,x_+)\,\dd x_-\,\dd x_+\,\dd c\]
over upper unitriangular $x_\pm$ and $c\in\CC^{n}$; the Jacobian of the triangular decomposition, reference to be
confirmed. And (unit 119) the convolution of the laws of the floors $0,\dots,\varrho$ at $\CC$ is the image of
$\prod_a\mathrm{dens}_a(c_a)\,\dd c$ under $c\mapsto(\wl{\varrho}{c}\boost{\varrho}{c_\varrho})^{-1}$, with
\[\mathrm{dens}_a(c)=|c|^{2(k-1)+4(\varrho-a)}\ \ (a<\varrho),\qquad
\mathrm{dens}_\varrho(c)=\frac{2\pi^{k}}{(k-1)!}\max(|c|^2-1,0)^{k-1};\]
and the Heisenberg measures of the floors
convolve to a Haar measure of $\UT{T}$, the image of $\dd x\,\dd\muz$ under $(x,u)\mapsto u\cdot\mathrm{lift}(x)$,
$\mathrm{lift}(x)\in\UT{T}$ with frame matrix the unitriangular matrix of entries $x$.

\Inp{J6} \emph{A Lipschitz bound for holomorphic maps} (unit 10): a holomorphic map bounded by $B$ on a ball of
radius $r$ has derivative of norm at most $2B/r$ at the centre; the second reader relates it to the Schwarz lemma
(W.~Rudin, \emph{Real and complex analysis}, \S12.2) as a weaker form.

\Inp{J7} \emph{A Sobolev chain inequality} (unit 10): for functions $f_w(x,\omega)$ indexed by words $w$ in a finite
alphabet $\mathcal{A}$, on a cube $[0,\epsilon]^{D}\times\Omega'$, $0<\epsilon\le1$, with
$|\partial_df_w|\le\Lambda_0\sum_{l\in\mathcal{A}}|f_{lw}|$ and $\iint|f_w|^2\le K$ for $|w|\le|D|$:
$\int|f_\emptyset(0,\omega)|^2\,\dd\omega\le(2(1+(|\mathcal{A}|\Lambda_0)^2)/\epsilon)^{|D|}K$. Classed as elementary
by the catalogue.

\Inp{J8} \emph{Hua's integral bound} (unit 112): for a positive definite $M$, an invertible $k'\times k'$ matrix
$S$, a finite set $\mathsf{s}$ of rows and $r'\ge k'$,
$\int\det(M+\sum_{i\in\mathsf{s}}y_i^{*}y_i)^{-(|\mathsf{s}|+r')}\dd x\le(\pi^{k'}/(k'!\,|\det S|^2))^{|\mathsf{s}|}\det M^{-r'}$,
where $y_i=x_iS+a_i$ and the rows $x_i$ run over $\CC^{k'}$. The second reader relates it to L.~K.~Hua, \emph{loc.\
cit.}, as a weaker form (an upper bound in place of the exact Gamma factors).

\Inp{J9} \emph{The Ostrowski--Taussky inequality} (unit 10): $\det H\le2^{N}|\det\Theta|$ for a positive semidefinite
$H\le\Theta+\Theta^{*}$; A.~Ostrowski, O.~Taussky, \emph{On the variation of the determinant of a positive definite
matrix}, Indag.\ Math.\ 13 (1951); the second reader: matches. And the \emph{matrix Bessel inequality}:
$X(1-EE^{*})X^{*}\ge0$ for $\|E\|\le1$.

\Inp{J10} \emph{Discrete series} (units 112, 116, 117), for a parameter with large gap $A_0$:
$\deg_l\psi_l$ is an idempotent kernel on $\Gone{\mathfrak{s}}$ (integrable, bounded by $\deg_l$, reproducing under
convolution); $|\psi_l(g)|\le(2/\|g\|)^{c}$ for $c\le500r^{3}+A_0$; $\deg_l(2/\|g\|)^{c}$ has integral at most
$8(2000(r+L+1)^{6})^{2r^{2}}2^{c}$ for $c>8r^{2}$; the module generated by $\psi_l$ under right derivatives is
unitary for the $L^2$ norm, its generator is a smoothing of $\psi_l$, and a derivative along a word $w$ of a vector
of the $|w|$-th step of its filtration has norm at most $((r+1)^{4}(L+r+1+|w|))^{|w|}$; and for a Bessel function
$F$ with $F*e=F$, $e=\deg_l\psi_l$, and $X=e*X$ integrable:
$\int|\int\varphi(y)(F*X)(\gamma y)\dd y|^2\dd\dot\gamma\le\deg_l^{-1}\|X\|_2^2\,\|F\|_2^2\int|\varphi|^2$.
The second reader relates the idempotent kernel to Harish-Chandra, \emph{Discrete series for semisimple Lie
groups II} (1966), pp.~87--88, as a special case whose explicit thresholds are unresolved.

\Inp{J11} \emph{A tube} (unit 118): for $t_0\ge8r$ there are a compact $A\subseteq\{\|y\|\le t_0^{20n}\}$ and a
measurable $\Omega_0\subseteq\mathcal{U}^0$ with $0<\muz(\Omega_0)<\infty$,
$\Omega_0\cdot\{\|y\|\le t_0\}\subseteq A$ and $\int\mathbf{1}_A(uy)\,\dd\muz(u)\le t_0\,\muz(\Omega_0)$ for all $y$.
Conjugation by $s\in\Ql{n}{T}$ preserves $\muz$.

\sect{Frame vectors and the Bessel size}{s:frame}
\begin{lemma}[frame vectors]\label{l:fv}
Let $y,s\in\Gone{\mathfrak{s}}$ and $a\le\varrho$. \textup{(i)} If $s$ lies in the group generated by
$\mathcal{U}^0$ and $\Hl{n}{T}$, then $(sy)^{-1}f_a=y^{-1}f_a$. \textup{(ii)} If $s\in\UT{T}$, then
$y^{-1}f_a=\sum_{b}\mathrm{fm}(s)_{ba}\,(sy)^{-1}f_b$. \textup{(iii)} For a finite set $\iota$, a map $e$ from $\iota$
to the frame positions, $c\colon\iota\to\CC$ and a basis index $i$:
$\bigl|\sum_{\alpha\in\iota}c_\alpha\,s_{i,f_{e(\alpha)}}\bigr|\le\|sy\|\cdot\bigl|\sum_{\alpha\in\iota}c_\alpha\,y^{-1}f_{e(\alpha)}\bigr|$.
\end{lemma}
\begin{proof}
(i) $s^{-1}$ fixes the frame. (ii) $y^{-1}f_a=(sy)^{-1}sf_a$, and $sf_a=\sum_b\mathrm{fm}(s)_{ba}f_b$ because
$s\in\UT{T}$ maps $f_a$ into the span of $f_0,\dots,f_a$. (iii) $\sum c_\alpha sf_{e(\alpha)}=(sy)\sum c_\alpha
y^{-1}f_{e(\alpha)}$, and a coordinate of a vector is at most its norm.
\end{proof}
\leanline{\lean{frameVectormul}\sep\lean{frameVectoreqsum}\sep\lean{normsumlemulnormsumframeVector}\sep\lean{continuousframeVector}\sep\lean{normframeVectorle}}
The frame vectors depend continuously on $y$, and $|y^{-1}f_a|\le\|y^{-1}\|=\|y\|$.

\begin{lemma}[the Bessel size]\label{l:bsnorm}
Let $x,y,y'\in\Gone{\mathfrak{s}}$ and $t_0$ a real number. \textup{(i)} If $\bs{y}\le t_0$, there is $u\in\UT{T}$ with
$\|uy\|\le t_0$. \textup{(ii)} $\bs{y'}\le\|x^{-1}y'\|\cdot\bs{x}$. \textup{(iii)} $y\mapsto\bs{y}$ is continuous.
\end{lemma}
\begin{proof}
For $s\in\Ql{n}{T}$ we have $\bs{y}\le\|sy\|+1$ and $\bs{sy}=\bs{y}$, by definition. (i) Choose $s\in\Ql{n}{T}$
with $\|sy\|+1<t_0+1$, and $h\in\mathrm{U}(\kappa)$ with $\emb{h}s\in\UT{T}$; $\emb{h}$ is unitary, so
$\|\emb{h}sy\|=\|sy\|\le t_0$. (ii) For $s\in\Ql{n}{T}$,
$\|sy'\|+1\le\|sx\|\|x^{-1}y'\|+1\le\|x^{-1}y'\|(\|sx\|+1)$, because every element of $\Gone{\mathfrak{s}}$ has norm
at least $1$; take the infimum. (iii) By (ii) in both directions, $\bs{y'}/\bs{y}$ lies between
$\|y'^{-1}y\|^{-1}$ and $\|y^{-1}y'\|$, which tend to $1$.
\end{proof}
\leanline{\lean{existsnormleofbesselSizele}\sep\lean{besselSizelenormmul}\sep\lean{continuousbesselSize}\sep\lean{besselSizele}\sep\lean{besselSizebesselSmul}}

\begin{lemma}[bounded frame and dual frame]\label{l:framebox}
Let $x\in\Gone{\mathfrak{s}}$ and $M$ a real number with $1\le M$, $x\in\frbox{M}$ and $|x^{-1}e^*_a|\le M$ for all
$a<\varrho$. Then $\|x\|+1\le8(\varrho+1)M$.
\end{lemma}
\begin{proof}
Let $B=x^{-1}$. As $B\Jone{\mathfrak{s}}B^{*}=\Jone{\mathfrak{s}}$,
$BB^{*}=\Jone{\mathfrak{s}}+B(1-\Jone{\mathfrak{s}})B^{*}$ and $\|B\|^2\le1+\|B(1-\Jone{\mathfrak{s}})\|\,\|B\|$.
The non-zero columns of $B(1-\Jone{\mathfrak{s}})$ are $\pm(Be_a-Be^*_a)$ and $2Bz$, of norm at most $2M$; there
are $2\varrho+1$ of them, so $\|B(1-\Jone{\mathfrak{s}})\|\le2M\sqrt{2\varrho+1}$. The quadratic inequality gives
$\|B\|\le8(\varrho+1)M-1$, and $\|x\|=\|B\|$.
\end{proof}
\leanline{\lean{normleofmemframeBox}\sep\lean{norminclMatrixleone}\sep\lean{normJ}\sep\lean{norminnerJle}}
Here $\|\Jone{\mathfrak{s}}\|=1$, matrices of inclusions of coordinates have norm at most $1$, and
$|\ip{v}{\Jone{\mathfrak{s}}w}|\le|v||w|$.

\begin{lemma}[the frame Gram matrix]\label{l:rfg}
Let $x,y'\in\Gone{\mathfrak{s}}$ and $a',b\le\varrho$. Then
$Z(x)_{a'b}=\ip{y'^{-1}f_{a'}}{\Jone{\mathfrak{s}}(x^{-1}y')^{-1}f_b}$, hence
$|Z(x)_{a'b}|\le|y'^{-1}f_{a'}|\cdot|(x^{-1}y')^{-1}f_b|$; $Z(sx)=Z(x)=Z(xs)$ for $s$ in the group generated by
$\mathcal{U}^0$ and $\Hl{n}{T}$; and $Z(x)=\mathrm{fm}(s)^{*}Z(sx)$ for $s\in\UT{T}$.
\end{lemma}
\begin{proof}
$h(f_{a'},xf_b)=h(y'^{-1}f_{a'},y'^{-1}xf_b)$; the rest follows from Lemma~\ref{l:fv}(i),(ii).
\end{proof}
\leanline{\lean{realFrameGrameqinner}\sep\lean{normrealFrameGramlemul}\sep\lean{realFrameGrammulinvariant}\sep\lean{realFrameGrameqframeMatrixmul}}

\begin{lemma}[the Bessel size from a dual system]\label{l:dual}
Let $x\in\Gone{\mathfrak{s}}$, $M$ a real number with $1\le M$, and $V_0,V_1,\dots$ vectors of $\CC^{r}$ such that
$x^{-1}f_a=V_a$ for $a\le\varrho$, $|V_a|\le M$ for $a\le2\varrho$, and, for all $a,b$ with $a+b\le2\varrho$,
$\ip{V_a}{\Jone{\mathfrak{s}}V_b}=-1$ if $a+b=2\varrho$ and $=0$ otherwise. Then $\bs{x}\le8(\varrho+1)M$.
\end{lemma}
\begin{proof}
Let $f'_0,\dots,f'_\varrho$ be the Gram--Schmidt orthogonalisation of the frame vectors of $x$. There is
$u\in\UT{T}$, with unitriangular frame matrix, such that $(ux)^{-1}f_a=f'_a$ (Lemma~\ref{l:fv}(ii)). The hypothesis
gives $\ip{V_{2\varrho-a}}{\Jone{\mathfrak{s}}f'_a}=-1$, so $1\le M|f'_a|$; and $|f'_a|\le|V_a|\le M$. Put
$v_a=|f'_a|^{-2}f'_a$, of norm at most $M$. The vectors $\Jone{\mathfrak{s}}v_b$, $b<\varrho$, are dual to the frame
vectors of $ux$ and span an isotropic space (the frame Gram matrix of the identity vanishes at the isotropic
positions); so by a lemma of unit 118 there is $q$ with $(ux)^{-1}\uq{q}e^*_b=\Jone{\mathfrak{s}}v_b$. The element
$y=\uq{q}^{-1}ux$ has frame vectors $f'_a$ and dual columns $\Jone{\mathfrak{s}}v_b$, all of norm at most $M$, and
$\bs{x}\le\|y\|+1\le8(\varrho+1)M$ by Lemma~\ref{l:framebox}.
\end{proof}
\leanline{\lean{besselSizeleofdual}}

\sect{The section}{s:section}
\begin{lemma}[the section is unitary; measurability]\label{l:gsec}
Let $Z\in\Ocell{\mathfrak{s}}$. Then the matrix
\[\sectionMatrix{Z}{Z^{*}}{\sqrt{\cornerDefect{Z}}}\]
of Definition~\ref{d:section} is unitary for the form $\Jone{\mathfrak{s}}$. The map $Z\mapsto\gramSection{Z}$ is measurable.
\end{lemma}
\begin{proof}
Let $C$ be the matrix of the $2n$ columns $(F_{\mathfrak{s}}\,|\,F'_{\mathfrak{s}})$; they are orthonormal, and
$\Jone{\mathfrak{s}}=1-CC^{*}+CJC^{*}$ with $J=C^{*}\Jone{\mathfrak{s}}C$. With
$\Delta=\diagLast{\cornerDefect{Z}^{-1/2}}$ and the block matrices
$\mathsf{x}=\bigl(\begin{smallmatrix}Z^{-1}&\sectionCoeff{Z^{*}}{Z}\Delta\\0&1\end{smallmatrix}\bigr)$,
$\mathsf{y}=\bigl(\begin{smallmatrix}\sectionCoeff{Z}{Z^{*}}&Z^{*-1}\Delta\\\diagLast{\sqrt{\cornerDefect{Z}}}&0\end{smallmatrix}\bigr)$,
the section matrix is $1-CC^{*}+C\,\mathsf{y}\mathsf{x}^{-1}C^{*}$, and
$\mathsf{y}^{*}J\mathsf{y}=\mathsf{x}^{*}J\mathsf{x}$; both are matrix identities that the formal text verifies by
expansion. So $\mathsf{y}\mathsf{x}^{-1}$ is unitary for $J$ and the section matrix is unitary for
$\Jone{\mathfrak{s}}$. Measurability: the map is continuous on the open set $\Ocell{\mathfrak{s}}$ and constant
outside.
\end{proof}
\leanline{\lean{measurablegramSection}}

\begin{lemma}[differentiability and size of the section matrix]\label{l:secnorm}
\textup{(i)} Let $Z,W\in\Ocell{\mathfrak{s}}$. Then $(Z',W')\mapsto\sectionMatrix{Z'}{W'}{s(Z',W')}$, with
$s(Z',W')=\sqrt{1-\invCorner{Z'}\invCorner{W'}}$, is complex differentiable at $(Z,W)$. \textup{(ii)} For
$Z\in\Ocell{\mathfrak{s}}$ the map $Z'\mapsto\sectionMatrix{Z'}{Z'^{*}}{\sqrt{\cornerDefect{Z'}}}$ is real
differentiable at $Z$. \textup{(iii)} Let $Z\in\Ocell{\mathfrak{s}}$, $P,t_0$ real numbers with $t_0>0$ and
$8\le P$, and suppose \m{\IsBoundedInCell{Y}{P}} for every $Y$ with $\|Y-Z\|<t_0$. Then the matrix of (ii) at $Z$
has norm at most $P^{12}$, and its real derivative at $Z$ has norm at most $2P^{12}/t_0$.
\end{lemma}
\begin{proof}
(i) Every entry is a polynomial in the entries of $Z'$, $W'$, of their inverses, and in $s$, $s^{-1}$; and
$1-\invCorner{Z'}\invCorner{W'}$ stays in the slit plane near $(Z,W)$, because $|\invCorner{Z}|,|\invCorner{W}|<1$.
(ii) Compose with the real-linear map $Z'\mapsto(Z',Z'^{*})$. (iii) For $(X,Y)$ within $t_0$ of $(Z,Z^{*})$ both $X$
and $Y^{*}$ are within $t_0$ of $Z$, so their norms, those of their inverses and
$\cornerDefect{X}^{-1/2}$ are at most $P$; moreover
$|1-\invCorner{X}\invCorner{Y}|\ge1-|\invCorner{X}|\ge\tfrac14\cornerDefect{X}$, which bounds $|s|^{-1}$ by $2P$.
The triangle inequality in Definition~\ref{d:section} then bounds the section matrix by $P^{12}$ on this ball, for
$P\ge8$. A holomorphic map bounded by $P^{12}$ on a ball of radius $t_0$ has derivative at most $2P^{12}/t_0$ at
the centre \Inp{J6}; restrict to $Y=X^{*}$.
\end{proof}
\leanline{\lean{differentiableAtsectionMatrixuncurry}\sep\lean{differentiableAtsectionMatrix}\sep\lean{normsectionMatrixle}}

\begin{lemma}[elements over the open cell]\label{l:double}
Let $\gamma\in\Gone{\mathfrak{s}}$ and $Z\in\Ocell{\mathfrak{s}}$ with $Z(\gamma)=Z$. Then there are parameters
$q,q'$ of $\mathcal{U}^0$ and $k,h\in\mathrm{U}(\kappa)$ with
$\gamma=\uq{q}\,\emb{k}\,\gramSection{Z}\,\bigl(\uq{q'}\,\emb{h}\bigr)^{-1}$.
\end{lemma}
\begin{proof}
For $s\in\Gone{\mathfrak{s}}$ with $Z(s)=Z$ the $2n$ columns $(F_{\mathfrak{s}}\,|\,sF_{\mathfrak{s}})$ have the Gram
matrix $\bigl(\begin{smallmatrix}N_0&Z\\Z^{*}&N_0\end{smallmatrix}\bigr)$ with $N_0=\operatorname{diag}(0,\dots,0,-1)$;
it depends on $Z$ only, and it is invertible because its determinant is a non-zero multiple of $\gd{Z}$. So the columns have a left inverse. This applies to $\gamma$ and to
$\gramSection{Z}$, and Witt's theorem \Inp{J1} gives $g\in\Gone{\mathfrak{s}}$ with $gF_{\mathfrak{s}}=F_{\mathfrak{s}}$
and $g\gamma F_{\mathfrak{s}}=\gramSection{Z}F_{\mathfrak{s}}$. An element that fixes the frame is
$\emb{k}\uq{q}^{-1}$ for some $k$, $q$ (unit 118); apply this to $g$ and to $\gramSection{Z}^{-1}g\gamma$.
\end{proof}
\leanline{\lean{existseqmulgramSectionmul}}

\begin{lemma}[staying in the open cell]\label{l:near}
Let $M,P$ be real numbers and $Z,Z'$ complex $n\times n$ matrices with $M^{-1}\le\gd{Z}$, $8(\varrho+1)\le M$,
$|Z_{ab}|\le M^{2}$ for all $a,b$, $M^{9n}\le P$ and $\|Z'-Z\|<P^{-6}$. Then $\gd{Z}/2\le\gd{Z'}$ and
\m{\IsBoundedInCell{Z'}{P}}.
\end{lemma}
\begin{proof}
On the unit ball about $Z$ the entries are at most $M^{3}$, so the squares of the leading minors are holomorphic and
bounded by a power of $M$ that is at most $P$; by \Inp{J6} they move by at most $M^{-1}/4$ between $Z$ and $Z'$, and
$\gd{\cdot}$, a difference of two of them, by at most $M^{-1}/2\le\gd{Z}/2$. So $\gd{Z'}>0$ and
$Z'\in\Ocell{\mathfrak{s}}$. From $M^{-1}\le2|\det Z'|^{2}\cornerDefect{Z'}$ and $\cornerDefect{Z'}\le1$ one gets
lower bounds for $|\det Z'|$ and for $\cornerDefect{Z'}$; Cramer's rule bounds $\|Z'^{-1}\|$ by minors over
$|\det Z'|$; all bounds are powers of $M$ with exponent at most $9n$.
\end{proof}
\leanline{\lean{memopenCellandboundsofnear}}

\sect{The chart}{s:chart}
\begin{lemma}[the logarithmic derivative of the outer factors]\label{l:logd}
Let $\mathcal{V}$ be a real normed space, $p\in\mathcal{V}$, and $\mathsf{z}=(\mathsf{z}_1,\mathsf{z}_2)$ a map from
$\mathcal{V}$ to pairs of a real $\kappa\times\kappa$ matrix and a parameter of $\mathcal{U}^0$ that is
differentiable at $p$ with derivative $\mathsf{z}'$. Then $\Phi(p')=\uq{\mathsf{z}_2(p')}\,\emb{\cay{\mathsf{z}_1(p')}}$ is
differentiable at $p$, and for every $v$, with $\mathsf{z}'(v)=(\eta,(q_1,q_2))$,
\[\mathrm{D}^{\log}\Phi(p)(v)=\btan\Bigl(\mathrm{D}^{\log}\operatorname{cay}(\mathsf{z}_1(p))(\eta),\ q_1\cay{\mathsf{z}_1(p)},\
q_2+\mathrm{sk}\bigl(\text{the skew part of }\mathsf{z}_2(p)_1\,q_1^{*}\bigr)\Bigr).\]
\end{lemma}
\begin{proof}
By the formulas of unit 118, $\uq{q}=1+\Jone{\mathfrak{s}}\beta+\tfrac12(\Jone{\mathfrak{s}}\beta)^{2}$ with
$\beta=\btan(0,q)$, and $\emb{c}=1+\btan(c-1,0)$. Both are differentiable in the parameters, the Cayley chart by
\Inp{J4}; differentiate the product by the product rule and multiply by $\Phi(p)^{*}\Jone{\mathfrak{s}}$. The matrix
identity that remains is verified by the formal text by expansion.
\end{proof}
\leanline{\lean{logDerivUunipotentOfmulembedH}}

\begin{lemma}[coordinates]\label{l:coords}
There is a real-linear isomorphism $e$ from the space of the parameters onto the space of the real $r\times r$
matrices that preserves Lebesgue measure and has $e^{-1}=\chartc\circ\operatorname{skew}$.
\end{lemma}
\begin{proof}
$\chartc\circ\operatorname{skew}$ maps the sum of the squares of the entries of a real matrix to the sum of the
squares of the real coordinates of the parameter (a regrouping of blocks), and the dimensions agree; a linear map
with this property is a measure-preserving isomorphism (the Gaussian integral is preserved).
\end{proof}
\leanline{\lean{existssymmeqchartCoordsskewOfReal}}

\begin{proof}[Proof of Theorem A\textup{(i)}]
Take $e$ from Lemma~\ref{l:coords}, and let $p=(Z,\theta,q,\eta,q')$ with $Z\in\Ocell{\mathfrak{s}}$. The chart is
a product $\Phi_1(p)\,\gramSection{Z}\,\Phi_2(p)^{-1}$ with $\Phi_1=\uq{q}\emb{\cay{\tilde\theta}}$,
$\Phi_2=\uq{q'}\emb{\cay{\eta}}$; these are differentiable by Lemma~\ref{l:logd}, and the section by
Lemma~\ref{l:secnorm}(ii) and because $\gramSection{\cdot}$ has the section matrix near $Z$. By the product rule, the logarithmic derivative of the chart is the sum of the logarithmic derivative of $\Phi_1$
conjugated by $\gramSection{Z}\Phi_2^{-1}$, of that of the section conjugated by $\Phi_2^{-1}$, and of minus that of
$\Phi_2$ conjugated by $\Phi_2^{-1}$. For $G$ invertible, the map $S\mapsto G^{*}SG$ on the skew-hermitian matrices,
read through $e^{-1}\circ\mathrm{sk}$, has a determinant of absolute value $|\det G|^{2r}$ (a lemma of unit 115). The
formal proof applies this with $G=\Phi_2(p)K$, where $\Delta=\diagLast{\cornerDefect{Z}^{-1/2}}$,
\[\mathsf{M}(A,B)=F_{\mathfrak{s}}\bigl((A\Delta-1)F_{\mathfrak{s}}^{*}+B\Delta F'^{*}_{\mathfrak{s}}\bigr),\qquad
K=1+\mathsf{M}\bigl(Z^{-1},\sectionCoeff{Z^{*}}{Z}\bigr).\]
The matrix $K$ is chosen so that $\gramSection{Z}K=\mathsf{R}+\mathsf{M}(\sectionCoeff{Z}{Z^{*}},Z^{*-1})$ for a
constant matrix $\mathsf{R}$, and $\btan(\cdot)\,\mathsf{M}(A,B)=0=\mathsf{M}(A,B)^{*}\btan(\cdot)$; and
$|\det K|=|\det(Z^{-1}\Delta)|$, so $|\det G|^{2}=|\det Z|^{-2}\cornerDefect{Z}^{-1}=\gd{Z}^{-1}$. After the
congruence by $G$ the endomorphism is block triangular in the variables $(Z;\theta,q;\eta,q')$, with these diagonal
blocks:
\begin{itemize}
\item on $Z$: the map $S\mapsto(Z^{*-1}\Delta)^{*}S\,(Z^{-1}\Delta)$ of the complex $n\times n$ matrices, because the
  $Z$-component of $\chartc$ reads $F'^{*}_{\mathfrak{s}}(\cdot)F_{\mathfrak{s}}$ and
  $Z'\mapsto F_{\mathfrak{s}}^{*}\Jone{\mathfrak{s}}\gramSection{Z'}F_{\mathfrak{s}}=Z'$ has the identity as derivative;
  its real determinant has absolute value $|\det(Z^{-1}\Delta)|^{4n}=\gd{Z}^{-2n}$;
\item on $(\theta,q)$ and on $(\eta,q')$: by Lemma~\ref{l:logd}, maps that are triangular in themselves, with the
  logarithmic derivative of a Cayley chart (for $\theta$: of its column at $k_0$) and right multiplications by
  unitary matrices on the diagonal; their determinants have absolute values $\cayd{\tilde\theta}$ and $\cayd{\eta}$
  \Inp{J4}.
\end{itemize}
So $|\det(e^{-1}\circ\mathrm{sk}\circ\mathrm{D}^{\log}\chart(p))|\cdot\gd{Z}^{-r}=\gd{Z}^{-2n}\cayd{\tilde\theta}\cayd{\eta}$,
and $r-2n=k-1$. I read the formal computation line by line; that the blocks are as stated is verified there by
expansion of the matrices and is not reproduced here.
\end{proof}

\begin{proof}[Proof of Theorem A\textup{(ii)}]
Let two parameters with $Z\in\Ocell{\mathfrak{s}}$ have the same image $g$. By Lemma~\ref{l:rfg} and
$Z(\gramSection{Z})=Z$, the frame Gram matrix of $g$ is $Z$, so the two $Z$ agree. Multiplying the equation
$g\,\Phi_2=\Phi_1\gramSection{Z}$ by $F_{\mathfrak{s}}Z^{-1}\Delta$ on the right isolates the frame and dual frame
columns of the left factor $\Phi_1=\uq{q}\emb{\cay{\tilde\theta}}$: the map $q\mapsto\uq{q}D_{\mathfrak{s}}$ is
injective, and the column at $k_0$ of a Cayley chart determines $\theta$ (unit 115). So $q$ and $\theta$ agree, then
the right factors agree, and $q'$ and $\eta$ agree because the Cayley chart is injective where it is regular.
\end{proof}

\begin{proof}[Proof of Theorem A\textup{(iii)}]
By (i), (ii) and \Inp{J4} there is a real $t$ with
\[\int_{Z\in\Ocell{\mathfrak{s}}}\frac{w(p)}{\volU{r}}f(\mathrm{e}^{\mathrm{i}t}\chart(p))\,\dd p\le\int f(\mathrm{e}^{\mathrm{i}t}y)\,\dd y=\int f\,\dd y,
\qquad w(p)=\gd{Z}^{k-1}\cayd{\tilde\theta}\cayd{\eta},\]
for Haar measures normalised as in the background library. In
the chart integral, the integral over $\eta$ with the weight $\cayd{\eta}$ is the Haar integral over
$h=\cay{\eta}\in\mathrm{U}(\kappa)$, the integrals over $q,q'$ are those over $\mathcal{U}^0$, and the integral over
$\theta$ of the normalised weight $\cayd{\tilde\theta}$ is $1$; the constants combine to $\cellc{\mathfrak{s}}$.
For each $\theta$ the hypothesis with $k=\cay{\tilde\theta}$ bounds $L$ by the inner integral; integrate over
$\theta$.
\end{proof}

\begin{lemma}[orthonormality of the full function]\label{l:orth}
Let $B$ be a real number, $f$ a natural number, and $X$ a Bessel system for $(\mathfrak{s},B,f)$ with full function
$\Phi_0$. Then for all $i,i'\in I$: $\int\Phi_{0,i}\,\overline{\Phi_{0,i'}}\,\dd\dot y=1$ if $i=i'$, and $=0$ otherwise.
\end{lemma}
\begin{proof}
The products are integrable. The matrix $\mathcal{G}$ of these integrals commutes with every $\mathrm{rep}(h)$:
substitute $y\mapsto\emb{h}y$, which preserves $\dd\dot y$ because conjugation by $\emb{h}$ preserves the Haar
measure of $\UT{T}$, and use $\Phi_0(\emb{h}y)=\mathrm{rep}(h)\Phi_0(y)$. As $\mathrm{rep}$ is irreducible,
$\mathcal{G}$ is a scalar; its trace is $\int|\Phi_0|^2\dd\dot y=|I|$.
\end{proof}
\leanline{\lean{orthonormal}\sep\lean{euclideanenormsqeqsum}}

\sect{Mean squares over the chart}{s:sobolev}
\begin{lemma}[a weight supported near one frame matrix]\label{l:fmle}
Let $w,\varphi\colon\Gone{\mathfrak{s}}\to[0,\infty]$ be measurable, $K\in[0,\infty]$ and $R$ real with $1\le R$.
Suppose $w$ is invariant under left translations by $\Ql{n}{T}$; $\int\varphi(uy)\,\dd\muz(u)\le K$ for all $y$; and
for all $y$ with $\varphi(y)\ne0$ and all $s\in\UT{T}$ with $\varphi(sy)\ne0$, all entries of $\mathrm{fm}(s)$ have
absolute value at most $R$. Then $\int\varphi\,w\,\dd y\le K\,(2R)^{2n^{2}}\int w\,\dd\dot y$.
\end{lemma}
\begin{proof}
Unfold: $\int\varphi w\,\dd y=\int w(y)\int_{\UT{T}}\varphi(uy)\,\dd u\,\dd\dot y$. Fix $y$; if the inner integral is
not $0$ choose $u_0$ with $\varphi(u_0y)\ne0$ and translate by it. Write the Haar measure of $\UT{T}$ as the image
of $\dd x\,\dd\muz(u')$ under $(x,u')\mapsto u'\,\mathrm{lift}(x)$ \Inp{J5}; the frame matrix of
$u'\,\mathrm{lift}(x)$ is the unitriangular matrix of $x$, so the integrand vanishes unless all $|x_{ab}|\le R$.
The inner integral over $u'$ is at most $K$, and the volume of the box of the $x$ is at most
$(\pi R^{2})^{n^{2}}\le(2R)^{2n^{2}}$.
\end{proof}
\leanline{\lean{lintegralmulleofframeMatrixle}}

\begin{lemma}[a mean square along the section]\label{l:sob}
Let $P,c$ be real numbers, $D\in[0,\infty]$, $\mathrm{rep}$ as before, $\Phi_0\colon\Gone{\mathfrak{s}}\to\CC^I$ with
the Schur bounds with respect to $\mathrm{rep}$, $P$ and $D$, and $Z$ a complex $n\times n$ matrix. Suppose
$8r\le P$, $0<c$, and that every $Y$ with $\|Y-Z\|<P^{-6}$ has $c\le\gd{Y}$ and \m{\IsBoundedInCell{Y}{P}}. Let
$\varphi\colon\Gone{\mathfrak{s}}\to\CC$ be measurable and integrable with $\varphi(y)=0$ whenever $\|y\|>P^{20}$.
Then
\[\begin{aligned}\cellc{\mathfrak{s}}\,c^{\,k-1}\int_{\mathrm{U}(\kappa)}\int_{\{u\in\mathcal{U}^0:\|u\|\le P^{20}\}}
\Bigl|\int\varphi(y)\,\Phi_0\bigl(\gramSection{Z}(u\,\emb{h})^{-1}y\bigr)\,\dd y\Bigr|^{2}&\dd\muz(u)\,\dd h\\
\le\ P^{2000n^{2}}\,D&\int|\varphi|^{2}\,\dd y.\end{aligned}\]
\end{lemma}
\begin{proof}
Let $V_w$ be the derivative functions of the Schur bounds. For a word $w$ in pairs of pairs of indices put
$\psi_w(y)=\varphi(y)\prod(y^{-1})_{\alpha\delta}y_{\gamma\beta}/P^{40}$ over the letters of $w$; then
$|\psi_w|\le|\varphi|$, because $\|y\|,\|y^{-1}\|\le P^{20}$ where $\varphi\ne0$. Put
$f_w(x,(h,u))=\int\psi_w(y)V_{w_1}(\gramSection{Z+\epsilon(x)}(u\emb{h})^{-1}y)\,\dd y$, for $x$ in the cube
$[0,P^{-8}]^{2n^{2}}$ of the real coordinates $\epsilon$ of the matrices, which stays within $P^{-6}/2$ of $Z$.
\emph{Derivatives.} By Lemma~\ref{l:secnorm}(iii) with $t_0=P^{-6}/2$, the curve $x_d\mapsto\gramSection{Z+\epsilon(x)}$
has derivative $g\mathsf{A}$ with $\|P^{40}\mathsf{A}\|\le P^{112}$; differentiating under the integral and using the
rule for right derivatives of $V_{w_1}$ gives $|\partial_df_w|\le P^{112}\sum_l|f_{lw}|$. \emph{The cube.} For
$|w|\le2n^{2}$ the integral of $|f_w|^{2}$ over the cube and over $(h,u)$, times $\cellc{\mathfrak{s}}c^{k-1}$, is at
most the Haar integral of $\chi(\gamma)|\int\psi_w(y)V_{w_1}(\gamma y)\dd y|^{2}$ over the $\gamma$ with $Z(\gamma)$
in the cube, by Theorem~A(iii) and $\gd{\cdot}\ge c$; here $\chi$ is a cutoff in the $\mathcal{U}^0$-direction. By
Lemma~\ref{l:fmle}, with $R=P^{2}$ from Lemma~\ref{l:rfg} and the bounds on $Z(\cdot)$ and its inverse, and by
Schur's bound for $V_{w_1}$, this is at most $(2P^{2})^{2n^{2}}P^{2|w|}D\int|\varphi|^2$. \emph{Conclusion.} The
Sobolev chain inequality \Inp{J7} bounds the value at $x=0$, $w=\emptyset$, which is the left side, by
$(2(1+(|\mathcal{A}|P^{112})^{2})P^{8})^{2n^{2}}$ times the cube bound; with $8r\le P$ the product is at most
$P^{2000n^{2}}D\int|\varphi|^{2}$. This last comparison is the formal text's arithmetic.
\end{proof}
\leanline{\lean{lintegralgramSectionintegralle}}

\sect{The tower measure and the gauge at the complex place}{s:haar}
\begin{proof}[Proof of Theorem B]
Write $\cL=\cLC{\mathfrak{h}}$, $\lambda=\haarRatio{\mathfrak{h}}$. (ii) is the choice of $\lambda$
(Definition~\ref{d:ratio}). By Theorem~A(ii) of Note~T and \Inp{J5},
$\groupMeasure{\cL}=\etab{n}*\upperMeasure{\cL}{0}$ with $\etab{n}$ a Haar measure of $\UT{T}$; so
$(u,g)\mapsto ug$ carries this Haar measure times $\upperMeasure{\cL}{0}$ to $\lambda\,\dd y$, and \Inp{J3} gives
(iii). $\lambda\ne0$ because the tower measure is not zero (Note~T), and $\lambda\ne\infty$ follows from (iv).

(iv) It suffices to test one function, the indicator $f$ of the compact set $\{\|g\|\le2\}$, of positive finite
Haar measure: we show $\frac{\volU{k}}{\volU{r}}\int f\,\dd(\lambda\,\dd y)\le\int f\,\dd y$. By the lemma of Note~T on moving a Weyl law to the left, applied to all floors below $\varrho$, the left measure is
\[\lambda\,\dd y=\etab{n}^{-1}*\floorMeasure{\cL}{n}*\bigl(\varpi_0*\dots*\varpi_{\varrho-1}*\varpi'_\varrho\bigr)*\floorMeasure{\cL}{n}*\etab{n},\]
where $\varpi'_\varrho$ is the image of the density of the boost and $\etab{n}^{-1}$ is $\etab{n}$, which is
invariant under inversion. By \Inp{J5} the middle factor is the image of $\prod_a\mathrm{dens}_a(c_a)\,\dd c$ under
$c\mapsto g_c^{-1}$, $g_c=\wl{\varrho}{c}\boost{\varrho}{c_\varrho}$, and the two Haar measures of $\UT{T}$ are images
of $\dd x\,\dd\muz$. Only the $c$ with $|c_\varrho|>1$ count. For such $c$, unitriangular $x,x'$ and
$k_1,k_2\in\mathrm{U}(\kappa)$, the element $\gamma=\mathrm{lift}(x)^{-1}\emb{k_1}g_c\emb{k_2}\mathrm{lift}(x')$ has
the frame Gram matrix $Z=x^{*}\operatorname{diag}(c)\,x'$ (unit 148), with
$\gd{Z}=\prod_{a<\varrho}|c_a|^{2}\,(|c_\varrho|^{2}-1)>0$, so $Z\in\Ocell{\mathfrak{s}}$, and by
Lemma~\ref{l:double} $\gamma=\uq{q}\emb{k}\gramSection{Z}(\uq{q'}\emb{h})^{-1}$. Since conjugation by elements of
$\Ql{n}{T}$ and translation preserve $\muz$ \Inp{J11}, and $f$ is invariant under multiplication by unitary
$\emb{h}$ on both sides, the integral of $f$ over the remaining $\mathcal{U}^0$-variables equals
$\iint f(u\,\emb{k}\gramSection{Z}(v\emb{h})^{-1})\dd\muz(u)\dd\muz(v)$ for each $k,h$. The densities match: in
Bruhat coordinates \Inp{J5} $\dd Z=\prod_a|c_a|^{4(\varrho-a)}\dd x\,\dd x'\,\dd c$, and
\[\frac{\volU{k}}{\volU{r}}\prod_{a\le\varrho}\mathrm{dens}_a(c_a)=\cellc{\mathfrak{s}}\,\gd{Z}^{k-1}\prod_{a<\varrho}|c_a|^{4(\varrho-a)}.\]
So $\frac{\volU{k}}{\volU{r}}\int f\,\dd(\lambda\,\dd y)$ is at most the right side of the hypothesis of
Theorem~A(iii), for every $k$, and that theorem gives the claim.
\end{proof}

\begin{lemma}[an orthonormal system in a kernel]\label{l:ker}
Let $s_0$ be a natural number, $L$ a real number, $\mathsf{m}$, $\iota$ finite sets and $C$ a complex
$\mathsf{m}\times\iota$ matrix with $|\iota|=s_0+1+3n$, $4\le L$, $\operatorname{rank}C\le3n$ and $|\iota|\le L$.
Then there are a subset $D\subseteq\iota$ and a complex $\iota\times D$ matrix $O$ with $CO=0$, $O^{*}O=1$ and, for
every natural number $j$,
\[\int_{\{u\in\CC^{\iota}:\,|u_i|\le L^{j}\text{ for all }i\}}\Bigl(\frac{|O^{*}u|^{2}}{36}\Bigr)^{-s_0}\dd u\le(L^{2})^{99n(j+1)}.\]
\end{lemma}
\begin{proof}
Diagonalise $C^{*}C$ by a unitary $Q$; at least $|\iota|-3n=s_0+1$ eigenvalues vanish; let $D$ be $s_0+1$ of their
indices and $O$ the corresponding columns of $Q$. Then $CO=0$ and $O^{*}O=1$. Substituting $u\mapsto Q^{*}u$, the
domain goes into a box of side $L\cdot L^{j}$ (the entries of $Q$ have modulus at most $1$ and $|\iota|\le L$). On
$\CC^{D}$ pass to $t=|w|^{2}$, which carries Lebesgue measure to $\pi^{s_0+1}t^{s_0}/s_0!\,\dd t$: the integrand
$(t/36)^{-s_0}$ times this density is the constant $36^{s_0}\pi^{s_0+1}/s_0!$, bounded independently of $s_0$; the
range of $t$ and the volume of the box in the other coordinates are bounded by powers of $L$.
\end{proof}
\leanline{\lean{existsorthonormalkerlintegralle}}

\begin{proposition}[the gauge at the complex place]\label{p:gauge}
Let $\mathfrak{h}$ be real-place sizes in $\RealPlaceSizes{\mathfrak{s}}{\rho}{s_0}{b}$ and $R$ a real number with
$\rho\le R$. Then the sets $\{|x|\le R^{i}\}$ and the constant $R^{2}$ satisfy the axioms \textup{(G1)--(G6)} of a
gauge on $\cLC{\mathfrak{h}}$ (Note~T); this is the content of Definition~\ref{d:cgauge}.
\end{proposition}
\begin{proof}
The axioms of a size hold because $R\ge\rho\ge8r$: a sum of at most $r+1$ numbers of absolute value at most $R^{i}$
is at most $R^{i+1}$, and the entries of $\Jone{\mathfrak{s}}=\Jone{\mathfrak{s}}^{-1}$ are $0,\pm1$. (G1)--(G4):
closed discs are compact; $c_{\CC}=\pi^{2}\ge1$; $\modA(u)=|u|^{2}\le(R^{2})^{j}$; the volume $\pi R^{2j}$ is at most
$(R^{2})^{j+1}$; the unit is $1$, the entries of the basis have modulus at most $1$, $\tau_{\CC}\delta_{\CC}^{-1}$
has modulus $1$; entries of a matrix are at most its norm $\le\rho\le R$; the ball has radius $\tfrac14$. (G5):
by the formula for the flat chart at $\CC$, the points of size $j$ of the cone have their two free groups of
coordinates in discs of radii $2R^{j}/|\alpha+\mathrm{i}|$ and $R^{j}$, $\alpha$ the real coordinate of the skew
parameter, whose measure is a multiple of the Cauchy density; integrating gives at most $(R^{2})^{99n(j+1)}$.
(G6): let $x$ be in the box, $y$ a position and $v$ a matrix with $\varrho$ rows. Let $\mathsf{N}$ be the matrix of
the $n$ rows of $1-\Jone{\mathfrak{s}}$ at the indices of the $e_a$ and of $z$; a vector $w'$ with $\mathsf{N}w'=0$
has $\Jone{\mathfrak{s}}w'=w'$. Apply Lemma~\ref{l:ker} to the matrix $C$ with the rows of $v$, of
$\mathsf{N}x^{-1}\orthonormalPart{\cLC{\mathfrak{h}}}$ and of $\mathsf{N}(xy)^{-1}\orthonormalPart{\cLC{\mathfrak{h}}}$:
it has $\varrho+2n\le3n$ rows, and $|\iota|=k-1=s_0+1+3n\le R$. Take $O$ from the lemma and
$\varphi(u)=36/|O^{*}u|^{2}$. For $w$ in the image of $O$ and $g\in\{x^{-1},(xy)^{-1}\}$ the vector
$w'=g\,\orthonormalPart{\cLC{\mathfrak{h}}}w$ has $\mathsf{N}w'=0$, so $|w'|^{2}=h(w',w')=|w|^{2}$. Hence, for
$t\in\CC$ with $|t|\le|O^{*}u|/9$, the vector $w=t\,|O^{*}u|^{-2}\,OO^{*}u$ has norm at most $\tfrac19$, lies in $W$
and has $u^{*}w=t$: the set $\{u^{*}w:w\in W\}$ contains a disc of radius $|O^{*}u|/9$, which gives the first
inequality of (G6). For the second, if $c\in\mathcal{C}_u$, take such a $w$ with $|t|=|O^{*}u|/9$; then
$|c|\,|t|\le|h(z_1,z'_1)|+|h(z_2,z'_2)|\le\tfrac1{16}+\tfrac14$ for close pairs, so $\mathcal{C}_u$ lies in the disc
of radius $3/|O^{*}u|$. The integral bound is that of Lemma~\ref{l:ker}.
\end{proof}

\sect{The open-cell bound}{s:oc}
\begin{proof}[Proof of Theorem C]
\emph{Set-up.} Write $Z=Z(\gamma)$ and $M=\rho^{9n}/\gd{Z}$. By Lemma~\ref{l:rfg} and the hypotheses on $y$,
$|Z_{ab}|\le\rho^{2}$; hence $\gd{Z}\le|\det Z|^{2}\le\rho^{9n-1}$ (as $n\le\rho$), so $\rho\le M$, and
$M^{-1}\le\gd{Z}$. Put $Q=M^{9n}$, $t_0=Q^{5}$ and $P=t_0^{\,n}$. Lemma~\ref{l:near} gives
$Z\in\Ocell{\mathfrak{s}}$ and, for every $Z'$ within $P^{-6}$ of $Z$, $\gd{Z'}\ge\gd{Z}/2$ and
\m{\IsBoundedInCell{Z'}{P}}. By Lemma~\ref{l:double}, $\gamma=\uq{q_1}\emb{k_1}\gramSection{Z}(\uq{q_2}\emb{k_2})^{-1}$.
Take the tube $A$, $\Omega_0$ of \Inp{J11} for $t_0$; $A$ is compact and lies in $\{\|y\|\le P^{20}\}$. For
$x=(h,u')\in\mathrm{U}(\kappa)\times\mathcal{U}^0$ put $t(x)=\uq{q_2}\,\emb{k_2h^{-1}}\,u'^{-1}$, an element of the
group generated by $\mathcal{U}^0$ and $\Hl{n}{T}$, so that
$\gamma\,t(x)=\uq{q_1}\emb{k_1}\,\gramSection{Z}\,(u'\emb{h})^{-1}$. Let $m$ be the product of the Haar probability
measure of $\mathrm{U}(\kappa)$ and of $\muz$, restricted to $\mathrm{U}(\kappa)\times\Omega_0$ and normalised, and
\[\nu=\frac{P^{2000n^{2}}\,|I|}{\deg\cdot\cellc{\mathfrak{s}}\,(\gd{Z}/2)^{k-1}\,\muz(\Omega_0)}\cdot\mathbf{1}_A\,|\Phi_0|^{2}\,\dd y.\]

\emph{(i).} The norm of an integral $\int\beta(y')\Phi_0(sgy')\,\dd y'$ does not change when $s\in\Ql{n}{T}$ is
removed (unit 118). Summing first over $i'$, the square is the squared norm of
$\int\overline{\varphi(y')\Phi_{0,i}(\gamma t(x)y')}\,\mathbf{1}_A(y')\Phi_0(t(x)y')\,\dd y'$, in which $t(x)$ may be
removed from the last factor; summing then over $i$, the sum is
$\sum_{i'}|\int\varphi_{i'}(y')\Phi_0(\gamma t(x)y')\,\dd y'|^{2}$ with
$\varphi_{i'}=\mathbf{1}_A\,\varphi\,\overline{\Phi_{0,i'}}$, which does not depend on $x$; and
$\uq{q_1}\emb{k_1}$ may be removed from $\gamma t(x)$. For each $i'$, Lemma~\ref{l:sob} applies to $\varphi_{i'}$,
with $c=\gd{Z}/2$, the bound $P$ (the Bessel system has the Schur bounds at $P\ge B^{5}$ with $D=|I|/\deg$) and the
domain $\{\|u'\|\le P^{20}\}\supseteq\Omega_0$. Summing over $i'$, with
$\sum_{i'}|\varphi_{i'}|^{2}=\mathbf{1}_A|\varphi|^{2}|\Phi_0|^{2}$, gives (i).

\emph{(ii).} By Lemma~\ref{l:fmle} with $w=|\Phi_0|^{2}$, $\varphi=\mathbf{1}_A$, $K=t_0\,\muz(\Omega_0)$ from the
tube, and $R=P^{40}$ (for $y,sy\in A$ the entries of $\mathrm{fm}(s)$ are at most $\|sy\|\|y^{-1}\|$),
$\int_A|\Phi_0|^{2}\dd y\le t_0\,\muz(\Omega_0)(2R)^{2n^{2}}\int|\Phi_0|^{2}\dd\dot y$, and the last integral is $|I|$.
So $\nu(\Gone{\mathfrak{s}})$ is at most a power of $M$ times $|I|^{2}/(\deg\,\cellc{\mathfrak{s}}(\gd{Z}/2)^{k-1})$;
that the power is at most $M^{10^{5}n^{4}}$ is the formal text's arithmetic.

\emph{(iii).} Let $x=(h,u')$ with $u'\in\Omega_0$, and $y'\in\frbox{\rho}$ with $\bs{\gamma t(x)y'}\le\rho$. Put
$y''=(u'\emb{h})^{-1}y'$; its frame vectors are those of $y'$ (Lemma~\ref{l:fv}(i)), and
$\bs{\gramSection{Z}y''}\le\rho$ by invariance of the Bessel size. By Lemma~\ref{l:bsnorm}(i) there is
$w\in\UT{T}$ with $\|w\,\gramSection{Z}y''\|\le\rho$. The explicit section matrix expresses the columns
$y''^{-1}e^*_a$ by the frame vectors of $y''$ and of $w\gramSection{Z}y''$, with coefficients among the entries of
$Z^{-1}$, of $\cornerDefect{Z}^{-1/2}$ and of the inverse of $\mathrm{fm}(w)$, which is bounded through
Lemma~\ref{l:rfg} and Lemma~\ref{l:near} for $Z(w\gramSection{Z})$; so these columns have norm at most $Q^{4}$-times
a constant, and Lemma~\ref{l:framebox} gives $\|y''\|\le t_0$. Hence
$y'=u'\emb{h}y''\in\Omega_0\cdot\{\|y\|\le t_0\}\subseteq A$.
\end{proof}

\sect{The Bessel cut}{s:cut}
\begin{proof}[Proof of Theorem D\textup{(i)}]
Write $e=\deg_l\psi_l$; by \Inp{J10} it is a bounded integrable idempotent kernel. Let $\beta$ be a smooth compactly
supported function whose smoothing of $\psi_l$ is the generator of the module of \Inp{J10}; for a word $w$ let
$\gamma_w$ be the corresponding iterated derivative of $\beta$ (twisted by the frame) and
$X_w=\deg_l\cdot(\text{the smoothing of }\psi_l\text{ by }\gamma_w)$, a function on $\Gone{\mathfrak{s}}$. Then $X_w$
is integrable, $e*X_w=X_w$, and $F*X_w$ is the smoothing of $F=F*e$ by $\gamma_w$. Put $V_w=F*X_w$. It is continuous,
$V_\emptyset=F$ (the smoothing by $\beta$ reproduces, and $F*e=F$), it is a Bessel function for $\mathrm{rep}$
because left translations pass through the convolution, and the right derivative of $V_w$ along the elementary
matrix at $(\alpha,\beta)$ is $V_{(\alpha,\beta)w}$ by differentiation of the smoothing. For the Schur bound, the
last inequality of \Inp{J10} gives the constant $\deg_l^{-1}\|X_w\|_2^{2}\,\|F\|_2^{2}$ with
$\|X_w\|_2^{2}=\deg_l\cdot\|x_w\|^{2}$ for the vector $x_w$ of the module, and
$\|x_w\|\le((r+1)^{4}(L+r+1+|w|))^{|w|}\le B^{|w|}$ for $|w|\le2n^{2}$, by the hypothesis on $B$.
\end{proof}

\begin{proof}[Proof of Theorem D\textup{(ii)}]
Let $M=8r^{2}+1$, $R=b^{6n+2}$, $\phi(g)=(2/\|g\|)^{2M}$ and
$\Phi^{\sharp}(y)=\int\phi(x^{-1}sy)\,\dd\nu_S(s)$. \emph{Normalisation.} $Y$ is integrable for $\dd\dot y$, by the
peak bound and the integrability of $\psi_l$; let $\mathsf{q}=\int|Y|^{2}\dd\dot y$, finite and not $0$ as
$Y(x)\ne0$, choose $c$ with $c^{2}\mathsf{q}=|I|$, and put $\Phi_0=cY$; then $\int|\Phi_0|^{2}\dd\dot y=|I|$ and
$\Phi_0*(\deg_l\psi_l)=\Phi_0$; $\Phi_0$ is continuous because it is a convolution. \emph{Cut.} Let
$\chi(t)=\max(0,\min(1,R-t))$ and $\Phi(y)=\chi(\bs{y})\Phi_0(y)$, continuous by Lemma~\ref{l:bsnorm}(iii), a Bessel
function because the Bessel size is invariant under $\Ql{n}{T}$, with $|\Phi|\le|\Phi_0|$ and $\Phi(y)=0$ unless
$\bs{y}\le R$. \emph{The peak against the $L^{2}$ norm.} From $Y=Y*e$, $|e(g)|\le\deg_l\phi(g)$, the Cauchy--Schwarz
inequality with the weight $\phi$, the unfolding through $\nu_S$ and $\Phi^{\sharp}\le C^{2}b^{N}|Y(x)|$ (second
peak bound, $C=16^{M}$),
\[|Y(x)|^{2}\le\deg_l^{2}\int\phi\cdot\int|Y|^{2}\,\Phi^{\sharp}\,\dd\dot y\le\deg_l^{2}\int\phi\cdot C^{2}b^{N}|Y(x)|\cdot\mathsf{q},\]
so $|Y(x)|\le\deg_l^{2}\,C^{2}b^{N}\,\mathsf{q}\int\phi$.
\emph{The majorant.} Take $\mathsf{A}=|c|\,\Phi^{\sharp}$ and
$\mu=|c|\cdot\frac{C}{\vol\{\|g-1\|\le1\}}\,b^{N}|Y(x)|\cdot(2/\|g\|)^{M}\dd g$, $C=16^{M}$. Then:
$|\Phi_0|,|\Phi|\le\mathsf{A}$ by the first peak bound and the decay of $\psi_l$; where $\Phi_0(y)\ne\Phi(y)$,
that is $\bs{y}>R-1$, every $s$ has $\|x^{-1}sy\|>2b$ (Lemma~\ref{l:bsnorm}(ii) and $\bs{x}\le b^{5n}$), and the
decay with the larger exponent $2M+3r^{3}f\le2M+A_0$ gains the factor $b^{-3r^{3}f}$; $\mathsf{A}$ is
$\UT{T}$-invariant by the invariance of $\nu_S$; $\mathsf{A}(y)\le\mu\{g:\bs{yg}\le R\}$ by the second peak bound
and a comparison of $\phi$ on a unit ball; $\mu$ is finite, by the integral of the decay; and
$\mu(\Gone{\mathfrak{s}})\int\mathsf{A}\,\dd\dot y\le|c|^{2}C^{2}b^{N}|Y(x)|\int\phi\le|I|\,(\deg_l C^{2}\int\phi)^{2}b^{2N}
\le|I|\,b^{2r^{3}}$, by the bound for the peak, $\deg_l\int\phi\cdot C^{2}\le b^{99r^{2}}$ \Inp{J10} and
$N+99r^{2}\le r^{3}$. \emph{Closeness.} Where $\Phi\ne\Phi_0$ we have $b^{3r^{3}f}|\Phi-\Phi_0|\le b^{3r^{3}f}|\Phi_0|\le
\mathsf{A}$, and $\mathsf{A}\le\mu(\Gone{\mathfrak{s}})$; so
$b^{6r^{3}f}\int|\Phi-\Phi_0|^{2}\dd\dot y\le\mu(\Gone{\mathfrak{s}})\int\mathsf{A}\,\dd\dot y\le b^{2r^{3}}|I|$, and
$2r^{3}-6r^{3}f\le-10r^{2}$ because $f\ge1$ and $r\ge99$, which follows from $N+99r^{2}\le r^{3}$. The comparisons of exponents are those of the formal text.
\end{proof}

\sect{The volume weight}{s:vol}
\begin{lemma}[the volume weight]\label{l:volpos}
For $\gamma,y\in\Gone{\mathfrak{s}}$: \textup{(i)} $\frameVolume{\gamma}>0$; \textup{(ii)} if
$\gamma\in\mathcal{U}^0$, then $\frameVolume{\gamma y}=\frameVolume{y}$.
\end{lemma}
\begin{proof}
(i) If the matrix had a kernel vector $v\ne0$, the vector $w=\gamma^{-1}F_{\mathfrak{s}}v$ would be orthogonal, for
the Euclidean product, to all $e_a-e^*_a$ and to $z$; so $w_{e_a}=w_{e^*_a}$, $w_z=0$ and
$h(w,w)=2\sum|w_{e_a}|^{2}+\sum|w_{x_i}|^{2}\ge0$, while $h(w,w)=h(F_{\mathfrak{s}}v,F_{\mathfrak{s}}v)=-|v_\varrho|^{2}\le0$.
Hence $w$ has only coordinates at the $e_a,e^*_a$, equal in pairs and with $h(w,w)=0$, so $w=0$ and $v=0$. (ii)
$\gamma^{-1}=\uq{q}$, and the explicit matrix of $\uq{q}$ changes the matrix of the definition by a factor of
determinant $1$.
\end{proof}
\leanline{\lean{frameVolumepos}\sep\lean{frameVolumemulleft}}

\begin{proof}[Proof of Theorem E]
By Lemma~\ref{l:bsnorm}(i) take $\omega\in\UT{T}$ with $\|\omega y\|\le\theta$. Let $\mathsf{T}$ be $\theta$ times the
leading $j\times j$ block of $\mathrm{fm}(\omega)$, unitriangular up to the factor, so $\det\mathsf{T}=\theta^{j}$
and its minor on $\sigma$ has determinant $\theta^{|\sigma|}$. Let $\mathsf{B}$ be the inverse of the matrix in
Definition~\ref{d:vol} for $y$, times $((1-\Jone{\mathfrak{s}})F_{\mathfrak{s}})^{*}$; let
$\mathsf{p}=Z\mathsf{I}_jZ_j^{-1}$, $\mathsf{C}=\mathsf{T}\mathsf{p}^{*}\mathsf{B}$,
$\mathsf{c}=(VZ_j^{-1}\mathsf{T}^{*}-W_{\mathfrak{s}}^{*}\mathsf{C}^{*})\mathsf{E}$ and
$\mathsf{S}=R\,(\mathsf{T}\mathsf{E})^{*}\mathsf{E}$, an invertible $\sigma\times\sigma$ matrix with
$|\det\mathsf{S}|=\theta^{|\sigma|}|\det R|$. Define
\[K(a)=\Bigl(\frac{(2\theta^{2})^{j}}{|\det Z_j|}\Bigr)^{k+n}\det\bigl(1+(a\mathsf{S}+\mathsf{c})^{*}(a\mathsf{S}+\mathsf{c})\bigr)^{-(k+n)}.\]
The bound for $\int K$ is Hua's \Inp{J8}, with $M=1$, the $k$ rows of $a$ and $r'=n\ge|\sigma|$.

Now let $\gamma$, $A$ be as in the statement, $\mathsf{t}=\gamma F_{\mathfrak{s}}\mathsf{I}_jZ_j^{-1}$ and
$\mathsf{G}=A\mathsf{S}+\mathsf{c}$. \emph{The ratio as a determinant.} The normal form at $j$ gives, for each frame
position $b\ge j$, vectors $X_b$, $Y_b$ of the flags of the frame and of $\gamma f$, triangular with unit leading
coefficients up to a number of modulus $1$; with them and with
$\mathsf{M}=\mathsf{B}\mathsf{t}Z_j-\mathsf{I}_j$ the formal text shows
\[\frac{\frameVolume{y}}{\frameVolume{\gamma^{-1}y}}=\bigl|\det(1+\mathsf{M}\mathsf{p}^{*})\bigr|^{-1}=\bigl|\det(1+\mathsf{p}^{*}\mathsf{M})\bigr|^{-1}
=\frac{|\det\mathsf{T}|^{2}}{|\det Z_j|\,|\det(\mathsf{C}\mathsf{t}\mathsf{T}^{*})|}.\]
\emph{Two positivity facts.} Put $\mathsf{Y}=\Jone{\mathfrak{s}}\mathsf{t}\mathsf{T}^{*}-\mathsf{C}^{*}$,
$\mathsf{W}=W_{\mathfrak{s}}^{*}\mathsf{Y}$ and $\Theta=-\mathsf{C}\mathsf{t}\mathsf{T}^{*}$. First,
$1+\mathsf{W}^{*}\mathsf{W}\le\Theta+\Theta^{*}$: the difference is $\mathsf{D}(1-\mathsf{Q}\mathsf{Q}^{*})\mathsf{D}^{*}$
plus a matrix of the form $N^{*}N$, where $\mathsf{D}=\mathsf{C}y$ and
$\mathsf{Q}=\theta^{-1}(\omega y)^{-1}F_{\mathfrak{s}}\mathsf{I}_j$ has norm at most $1$ and
$\mathsf{D}\mathsf{Q}=1$; it is positive by the matrix Bessel inequality \Inp{J9}. Second,
$\mathsf{G}=\mathsf{W}\mathsf{E}$ by the hypothesis on $\frameXPart{\gamma}$, so
$\mathsf{W}\mathsf{W}^{*}-\mathsf{G}\mathsf{G}^{*}=\mathsf{W}(1-\mathsf{E}\mathsf{E}^{*})\mathsf{W}^{*}\ge0$, again by
\Inp{J9}. By $\det(1+XX^{*})=\det(1+X^{*}X)$, the monotonicity of the determinant on positive matrices and the
Ostrowski--Taussky inequality \Inp{J9},
\[\det(1+\mathsf{G}^{*}\mathsf{G})=\det(1+\mathsf{G}\mathsf{G}^{*})\le\det(1+\mathsf{W}\mathsf{W}^{*})=\det(1+\mathsf{W}^{*}\mathsf{W})
\le2^{j}\,|\det(\mathsf{C}\mathsf{t}\mathsf{T}^{*})|.\]
Combining, $\frameVolume{y}/\frameVolume{\gamma^{-1}y}\le(2\theta^{2})^{j}|\det Z_j|^{-1}\det(1+\mathsf{G}^{*}\mathsf{G})^{-1}$,
and the power $k+n$ is $K(A)$. I read this proof line by line; the determinant identity of the first step is the
formal text's computation with the vectors $X_b$, $Y_b$.
\end{proof}

\subhead{Four lemmas on lattice translates (for the unit on the second moment).} The remaining declarations of the
file of the volume weight concern the adelic objects of unit 148 and are used there. Let $E$ be a CM field with ring
of finite adeles $\mathbb{A}_{E,f}\supseteq\widehat{O}_E$, and $\mathsf{S}$ a \emph{filling} over $E$ and
$\mathfrak{s}$ (unit 148), with its natural number $\mathsf{g}$, its classes $c$, the elements $h_c$, the frame
Gram matrix $Z(\gamma)$ of an element $\gamma$ of the rational unitary group, the \emph{frame image}
$\gamma F_{\mathfrak{s}}$, its rows at the $x_i$ and $z$ (the \emph{middle part}), the matrix
$\mathrm{pl}_j(Z)$ whose rows $a<j$ are those of $Z$ and whose other rows are those of the identity, the first
$\varrho$ rows $\xi(Z)$ of $Z$, and the coordinates $\mathsf{x}_c$ of the standard basis of $\langle x_i\rangle$ in
the lattice of the class $c$. An element $\gamma$ \emph{has level $j$} if it is in normal form at $j$, the leading
minor of order $j$ of $Z(\gamma)$ is not $0$, and $\gd{Z(\gamma)}\ne0$ exactly when $j=n$.

\begin{lemma}[counting by congruence classes]\label{l:count}
Let $E$ be a CM field that is stably unramified in the sense of unit 148, $\alpha$ a set, $\pi,\sigma,\tau$ finite
sets, $x\in[0,\infty]$, $c$ a finite idele with $c\in\widehat{O}_E$, $B\subseteq\alpha$, and $Y$ a family of
$\sigma\times\pi$ matrices over $\mathbb{A}_{E,f}$, integral for $a\in B$. Let $m$ be a primitive $\pi\times\tau$
matrix over $\widehat{O}_E$, $s$ a natural number with a bijection between $\tau\sqcup\{1,\dots,s\}$ and $\pi$, $K$
a lattice containing $c\,\widehat{O}_E^{\sigma}$, $\mathsf{n}\colon\alpha\to\NN$ with the index of
$K\cap\Lambda_a$ in $\Lambda_a$ dividing $\mathsf{n}(a)$ for $a\in B$, where $\Lambda_a$ is the lattice spanned by the
columns of $Y(a)m$, and $l$ a natural number with $|\tau|\le l$. Suppose every subset $D\subseteq B$ on which the
columns of $Y(a)-Y(a')$ lie in $K$ for all $a,a'\in D$, and which is not empty, has at most $x$ elements. Then
\[\sum_{a\in B}\mathsf{n}(a)^{-l}\le\bigl([\widehat{O}_E:c\widehat{O}_E]^{|\sigma|}\bigr)^{|\sigma|+s}\cdot x.\]
\end{lemma}
\begin{proof}
Complete $m$ to an invertible matrix $U$ over $\widehat{O}_E$, and colour $a\in B$ by the columns of $Y(a)U$ modulo
$c\widehat{O}_E^{\sigma}$: there are at most $[\widehat{O}_E^{\sigma}:c\widehat{O}_E^{\sigma}]^{|\pi|}$ colours, and
it suffices to bound the sum over one colour class. In a class a Hermite form of $(Y(a)m)$ gives generators of
$\Lambda_a$, and a second colouring by these generators modulo $K$ has at most $\mathsf{n}(a)^{|\tau|}\le
\mathsf{n}(a)^{l}$ colours relevant for $a$; inside one sub-class the columns of $Y$ agree modulo $K$, so it has at
most $x$ elements. Summing $\mathsf{n}(a)^{-l}$ over the sub-classes gives the bound. The count of the colours in
terms of the index $[\widehat{O}_E:c\widehat{O}_E]^{|\sigma|}$ and of $s$ is the formal text's.
\end{proof}
\leanline{\lean{tsuminvindexpowleofcongr}}

\begin{lemma}[three lemmas on frame images]\label{l:frames}
Let $\mathsf{S}$ be a filling over $E$ and $\mathfrak{s}$, $c$ a class and $j$ a natural number.
\textup{(i)} Let $\gamma$ have level $j$, and $b\le\varrho$ a frame position with $b=\varrho$ if $b<j$. Then the
vector $\mathsf{y}=\gamma F_{\mathfrak{s}}\cdot(\text{column }b\text{ of }\mathrm{pl}_j(Z(\gamma))^{-1})$ has
coordinates $0$ at all $e^*_a$, and $5\sum_{i\in\kappa}|\mathsf{y}_{x_i}|_v^{2}\le1$ at every infinite place $v$
of~$E$.
\textup{(ii)} Let $\mathsf{l}$ be a matrix over $E$ with rows indexed by $\kappa$ such that $\mathsf{x}_c\mathsf{l}$ is
integral, and $i$ one of its columns with $\sum_{a\in\kappa}|\mathsf{l}_{ai}|_v^{2}<1$ at every infinite place $v$.
Then $\mathsf{l}_{ai}=0$ for every $a\in\kappa$.
\textup{(iii)} Let $\mathsf{S}$ be admissible, $t$ a label, $p,p'$ labelled terms of level $j$ over $c$ with the
label $t$, and $u,u'$ elements of the group $\mathcal{U}^0_c$ of the class; if the frame images of $u\,p_2$ and
$u'p'_2$ are equal, then $p=p'$.
\end{lemma}
\begin{proof}
(i) If $b\ge j$, the definition of the level gives a vector common to the flags of the frame and of $\gamma f$,
which is $\mathsf{y}$; its coordinates on $\langle x_i\rangle$ vanish. If $b<j$, then $b=\varrho$ and $j=n$, the
column is that of $Z(\gamma)^{-1}$, and a lemma of unit 148 computes the norm of the middle part of $\mathsf{y}$ at
each place and bounds it by the constant of the filling. (ii) The column is an integral vector of the definite
lattice of the class with all archimedean norms less than $1$, so it is $0$ (a lemma of unit 148, from: an algebraic
integer with all conjugates of absolute value less than $1$ is $0$, unit 121). (iii) $u\,p_2$ and $u'p'_2$ are
again terms over $c$; equal frame images put the quotient of the corresponding compact elements into
$\Ql{n}{T}$, and the uniqueness of the labelled term in its double coset (unit 148) gives $p=p'$.
\end{proof}
\leanline{\lean{middlesharedvectorle}\sep\lean{coleqzeroofnormltone}\sep\lean{eqofframeImageeq}}

\begin{lemma}[equal Gram blocks and congruent middle parts]\label{l:translate}
Let $E$ be stably unramified, $\mathsf{S}$ a filling over $E$ and $\mathfrak{s}$, $c,c',c''$ classes, $j$ a natural
number, $\gamma,\gamma'$ of level $j$ with $Z(\gamma')=Z(\gamma)$, $\iota$ a finite set and $\mathsf{q}$ an
injection of $\iota$ into $\{0,\dots,\varrho-1\}$ such that every $a$ outside its image has $j\le a$. Let
$\xi$ be the rows $\mathsf{q}(\iota)$ of $\xi(Z(\gamma))$ and
$\mathsf{t}=\mathsf{g}^{-1}\mathrm{pl}_j(Z(\gamma))^{-1}$. Suppose that every column of the conjugate transpose of
the difference of the middle parts of the positions $h_c^{-1}\gamma'h_{c''}$ and $h_c^{-1}\gamma h_{c'}$ lies in the
sum of $\mathsf{g}$ times the lattice spanned by the columns of $\xi^{*}$ and of the lattice of the vectors that
$\mathsf{t}^{*}$ maps to integral vectors. Then there is a $\kappa\times\iota$ matrix $\mathsf{l}$ over $E$ with
$\mathsf{x}_c\mathsf{l}$ integral such that the frame images of $\gamma'$ and $\gamma$ agree at the rows $e^*_a$ and
$z$ and differ by $\mathsf{g}\,\mathsf{l}\,\xi$ at the rows $x_i$ (the rows other than the $e_a$).
\end{lemma}
\begin{proof}
Let $\mathsf{x}$ be the difference of the middle parts of the frame images at the rows $x_i$. The hypothesis, read
in the lattice of the class, gives an adelic matrix $\Lambda'$ with $\mathsf{x}_c\Lambda'$ integral and
$\mathsf{x}_c(\mathsf{x}-\mathsf{g}\Lambda'\xi)\mathsf{t}$ integral; by density of $E$ in the finite adeles choose a
matrix $\mu$ over $E$ with the same two properties. Let $\mathsf{L}=(\mathsf{x}-\mathsf{g}\mu\xi)\mathsf{t}$. For a
column $b$ with $b\ge j$ or $b=\varrho$, Lemma~\ref{l:frames}(i) for $\gamma$ and for $\gamma'$ gives
$5\sum_i|\cdot|_v^{2}\le1$ for the middle parts of the two columns $b$ of the frame images times
$\mathrm{pl}_j(Z(\gamma))^{-1}$; so the column $b$ of $\mathsf{x}\mathsf{t}$ has $\sum_i|\cdot|_v^{2}\le\tfrac45\mathsf{g}^{-2}<1$
at every infinite place. The column $b$ of $\xi\mathsf{t}$ vanishes, by the first assertion of that lemma, because
the rows of a frame image at the $e^*_a$ are rows of $Z(\gamma)$. So the column $b$ of $\mathsf{L}$ has norms less
than $1$ and is $0$ by Lemma~\ref{l:frames}(ii). Hence $\mathsf{L}$ lives on the columns $\mathsf{q}(\iota)$, and
$\mathsf{l}=\mathsf{L}|_{\mathsf{q}(\iota)}+\mu$ satisfies
$\mathsf{x}=\mathsf{g}\,\mathsf{l}\,\xi$, because $\mathsf{t}^{-1}$ restricted to these columns is
$\mathsf{g}\,\xi$. The rows at $e^*_a$ and $z$ of the two frame images agree because $Z(\gamma')=Z(\gamma)$.
\end{proof}
\leanline{\lean{existsframeImagetoRowseqadd}}

\sect{Where the results are used}{s:used}
Theorem~B and Proposition~\ref{p:gauge} supply the level, the gauges and the Haar hypothesis for Note~L at the
complex place. Theorem~C is the open-cell bound of a Bessel system, Theorem~D produces the cut pair and the Schur
bounds of such a system from a peaked Bessel function of the unit on the discrete series (unit 133), and
Theorem~E and Lemmas~\ref{l:count}--\ref{l:translate} are used in the unit on the second moment (unit 148), in the
Schur test for the incoherent and degenerate pairs.

\sect{Notes on the formal text}{s:notes}
\textbf{A declaration outside the census.} The definition of the level at the complex place
(Definition~\ref{d:clevel}; lines 31 to 41 of the file of the gauge at $\CC$) is not listed by the census of the
artifact, because its attribute runs over two lines; so it is a member of no unit, and it is not in the
concordance. It is stated here from the Lean text, which the second reader also recovered from the pinned source.

\textbf{Proofs inside definitions.} Two definitions contain proofs: that the section matrix is unitary
(Lemma~\ref{l:gsec}) and that the data of Definition~\ref{d:cgauge} form a gauge (Proposition~\ref{p:gauge}); the
level at $\CC$ contains the verification of the axioms of a level. The dossier has no structured proofs for them;
Lemma~\ref{l:gsec} and Proposition~\ref{p:gauge} are written from the Lean terms.

\textbf{Argument lists not confirmed by Lean.} For the open cell, the section, and for the Lean theorems behind
Lemma~\ref{l:near}, Lemma~\ref{l:sob}, Theorem~D(i), Theorem~E and Lemma~\ref{l:translate}, the reader's dossier
derived the list of arguments by Lean's rules, because the census misreads names bound inside the declaration or
stops at a \texttt{let} in the statement (missing the shape, the frame instance and the class). Nothing was compiled,
so these lists are not confirmed by Lean. The statements of Lemma~\ref{l:logd}, Theorem~D(i),
Lemma~\ref{l:frames}(i), Lemma~\ref{l:translate} and Theorem~E continue after a \texttt{let} and were read to their
end in the Lean text.

\textbf{Differences from the natural statements.} (1) The inverse of a singular matrix is $0$ in the formal text;
this matters only in Definition~\ref{d:cell}. (2) The chart is defined on all parameters; where $Z$ is not in the
open cell the section is the identity. (3) Theorem~A(i) is stated with the real determinant of an endomorphism of
the parameter space, obtained from the logarithmic derivative through $\mathrm{sk}$ and $e^{-1}$. (4) In Theorem~B
the constant is a choice; (ii) holds by this choice, and the content is (i), (iii), (iv). (5) Theorem~C asks
$y\in\frbox{\rho}$ and $\gamma^{-1}y\in\frbox{\rho}$ only to bound the entries of $Z(\gamma)$. (6) In Theorem~D(ii)
the number $N$ is the Lean variable \texttt{n} and $n=\varrho+1$ is \texttt{𝔰.n}. (7) In Theorem~E the set $\sigma$
lies in the lowest universe, and the leading block is taken through the inclusion of $\{0,\dots,j-1\}$ given by
$j\le n$, an implicit argument. (8) Lemma~\ref{l:ker} is stated for any finite sets and uses the shape only through
$n$. (9) Lemma~\ref{l:count} allows $x=\infty$ and uses real powers of indices. (10) In Lemma~\ref{l:dual} the
family $V$ is indexed by all natural numbers and constrained up to $2\varrho$.

\textbf{Small results.} Twenty-seven declarations have no statement of their own; each is cited where it is used:
continuity and norms of frame vectors (Lemma~\ref{l:fv}), the two rules for the Bessel size
(Lemma~\ref{l:bsnorm}), norms of $\Jone{\mathfrak{s}}$ and of inclusions (Lemma~\ref{l:framebox}), the four formulas
for the frame Gram matrix (Lemma~\ref{l:rfg}), the properties of $\gd{\cdot}$ and of the open cell
(Definition~\ref{d:cell}), the identities for the section matrix (Definitions~\ref{d:section}, \ref{d:gsec}), the
bound $\rho\ge4$ (Definition~\ref{d:rps}), a sum of squares (Lemma~\ref{l:orth}) and two properties of the volume
weight (Definition~\ref{d:vol}).

\subhead{How this note was written.} From the structured proofs of the reader's dossier and from the Lean text. I
read all statements and definitions of the ten files in the Lean text, with their comments. I read line by line the
formal proofs behind Theorem~A(i), Theorem~B, Theorem~C, both parts of Theorem~D and Theorem~E, which the reader
marked as the least certain, and the proofs inside the definitions of the section, of the level and of the gauge.
The proofs of Sections~\ref{s:frame} and \ref{s:section}, of Theorem~A(ii),(iii), Lemmas~\ref{l:orth},
\ref{l:fmle}, \ref{l:sob}, \ref{l:ker}, \ref{l:volpos} and of the four lemmas on lattice translates are written
from the structured proofs, whose steps a program compared with the source; I did not read those Lean proofs line
by line. Not reproduced, and named where they occur: the matrix identities in Lemma~\ref{l:gsec} and
Lemma~\ref{l:logd}; the evaluation of the block determinants in Theorem~A(i); the determinant identity in
Theorem~E; the constants in the proofs of Theorems C and D and of Proposition~\ref{p:gauge}(G5); and the count of
colours in Lemma~\ref{l:count}. In these places the text says what the formal proof establishes and by which
inputs; it is an account of the formal proof, not an independent derivation. The references were taken from the
second reader's verdicts and the catalogue of the write-up, as cited; none was opened for this note. The program
check of the unit's dossier covered \UnitSteps{} leaf steps of \UnitProofs{} proofs.

\head{Appendix A. Concordance with the Lean artifact}
The letters of this note and the names of the formal text: $\Gone{\mathfrak{s}}$, $\Jone{\mathfrak{s}}$ are
\texttt{G₁}, \texttt{J₁}; $F_{\mathfrak{s}}$, $F'_{\mathfrak{s}}$, $W_{\mathfrak{s}}$, $D_{\mathfrak{s}}$ are
\texttt{frameCols}, \texttt{dualFrameCols}, \texttt{definiteCols}, \texttt{dualCols}; $Z(y)$ is
\texttt{realFrameGram}; $\mathrm{fm}$ is \texttt{frameMatrix}; $\gd{\cdot}$ is \texttt{gramDiscReal};
$\Ocell{\mathfrak{s}}$ is \texttt{Shape.openCell}; $\uq{q}$ is \texttt{unipotentOf}; $\emb{h}$ is \texttt{embedH};
$\muz$ is \texttt{besselHaar₀}; $\dd\dot y$ is \texttt{besselQuotMeasure}; $\bs{\cdot}$ is \texttt{besselSize};
$\frbox{\rho}$ is \texttt{frameBox}; $\operatorname{cay}$, $\mathsf{j}$ are \texttt{cayleyChart},
\texttt{cayleyDensity}; $\chart$, $\chartc$, $\btan$ are \texttt{openCellChart}, \texttt{chartCoords},
\texttt{besselTangent}; $\mathfrak{h}$ is \texttt{hF}; $\cLC{\mathfrak{h}}$ is \texttt{RealPlaceSizes.complexLevel}
(not in the census); $s_0$ is \texttt{s}; $R$ is \texttt{L}; $\Phi_0$ is \texttt{F}, \texttt{W .full} or \texttt{Y};
$N$ is \texttt{n}; $A_0$ is \texttt{A}; $\mathsf{S}$, $\mathsf{g}$ are \texttt{S}, \texttt{S.g}. Two statements as
the artifact has them:
\quoted{haarRatiospec}
\quoted{opencell}
{\fontsize{8.6}{11.6}\selectfont\renewcommand{\leanfont}{\ttfamily\fontsize{7.7}{10}\selectfont\color{leancol}}
\setlength{\LTleft}{0pt}\setlength{\LTright}{\fill}\setlength{\tabcolsep}{0pt}
\begin{longtable}{@{}p{30mm}p{56mm}p{43mm}r@{}}
\textbf{In this note} & \textbf{Lean declaration} & \textbf{File} & \textbf{Lines}\\\midrule\endhead
\ConcordanceRows
\end{longtable}}

\head{Notation}
\begingroup\fontsize{8.6}{10.9}\selectfont\raggedright
\setlength{\columnsep}{9mm}\setlength{\parskip}{0pt}
\begin{multicols}{2}
\nt{$\mathfrak{s}=(\varrho,\kappa,k_0)$, $n$, $k$, $r$}{the shape; $\varrho+1$; $|\kappa|$; $2\varrho+k+1$}
\nt{$V$, $e_a$, $e^*_a$, $x_i$, $z$, $f_a$, $h$}{the space, its basis, the frame, the hermitian form}
\nt{$\Jone{\mathfrak{s}}$, $\Gone{\mathfrak{s}}$, $\dd y$, $\dd\dot y$}{Gram matrix, unitary group, Haar measure, measure of $\UT{T}\backslash\Gone{\mathfrak{s}}$}
\nt{$F_{\mathfrak{s}}$, $F'_{\mathfrak{s}}$, $W_{\mathfrak{s}}$, $D_{\mathfrak{s}}$}{matrices of the frame, the dual frame, the $x_i$, the $e^*_a$}
\nt{$T$, $\Hl{l}{T}$, $\Ql{l}{T}$, $\UT{T}$, $\mathcal{U}^0$}{the form at $\CC$ and the subgroups of Note~T; the unipotent group fixing the frame}
\nt{$\uq{q}$, $\emb{h}$, $\muz$, $\mathrm{lift}$}{elements of $\mathcal{U}^0$, of $\Hl{n}{T}$; Haar measure of $\mathcal{U}^0$; a lift of unitriangular matrices}
\nt{$Z(y)$, $\mathrm{fm}(y)$, $\gd{Z}$, $D_\varrho$}{frame Gram matrix, frame matrix, discriminant, leading minor}
\nt{$\bs{y}$, $\frbox{\rho}$, $\|X\|$, $X^{*}$}{Bessel size, frame box, operator norm, conjugate transpose}
\nt{$\operatorname{ic}$, $\operatorname{cd}$, $\Ocell{\mathfrak{s}}$, $\operatorname{Cell}_{\le P}$}{corner, defect, open cell, bounded part}
\nt{$\operatorname{d}_\varrho$, $\mathsf{c}$, $\mathsf{f}$, $\mathsf{m}$, $\gamma(Z)$, $\Delta$}{the matrices of the section; the section; a diagonal matrix (in Section~\ref{s:vol}: the volume weight)}
\nt{$\chart$, $\chartc$, $\btan$, $\mathrm{D}^{\log}$, $\mathrm{sk}$, $\operatorname{skew}$}{chart, its coordinates, tangent map, logarithmic derivative, real coordinates of matrices}
\nt{$\operatorname{cay}$, $\mathsf{j}$, $\theta$, $\tilde\theta$, $\eta$, $q$, $q'$, $p$, $e$}{Cayley chart, its weight, parameters of the chart, the isomorphism of coordinates}
\nt{$\mathfrak{h}$, $\operatorname{RPS}$, $\rho$, $s_0$, $b$}{real-place sizes and their numbers}
\nt{$\cLC{\mathfrak{h}}$, $\lambda$, $\Gamma(\mathfrak{h};\rho\le R)$, $R$}{level, Haar ratio, gauge, its radius}
\nt{$\operatorname{Lev}$, $\operatorname{Gau}$, $\mathsf{M}$, $\mathsf{M}_{\ge l}$, $\theta_l$, $\etab{j}$, $\varpi_l$, $\mathsf{O}$}{notions of Note~T: levels, gauges, tower and upper measures, floor measures, Heisenberg measures, Weyl laws, orthonormal part}
\nt{$\boost{\varrho}{c}$, $\wl{\varrho}{c}$, $g_c$, $\mathrm{dens}_a$}{boost, Weyl element, their product, densities of the floors}
\nt{$\cellc{\mathfrak{s}}$, $\volU{N}$}{$\frac{\volU{k}}{\volU{r}}\frac{2\pi^k}{(k-1)!}$; $\prod_{i<N}2\pi^{i+1}/i!$}
\nt{$I$, $\ip{v}{w}$, $\mathrm{rep}$, $\Phi_0$, $\Phi$, $Y$}{index set, inner product, representation, full and cut functions, a peaked function}
\nt{$\deg$, $\deg_l$, $\psi_l$, $l$, $L$, $A_0$, $\varepsilon$, $\mathsf{p}$}{formal degrees, coefficient, parameter, its bound and gap, signs, frame}
\nt{$\nu_S$, $\sd$, $*$, $V_w$, $X_w$, $\beta$, $\gamma_w$}{measure of unit 133, superdiagonal sum, convolution, derivative functions, smoothing data}
\nt{$\Delta(\gamma)$, $\mathsf{x}_0(\gamma)$, $\mathsf{I}_j$, $\mathsf{E}$, $Z_j$, $K$}{volume weight, part in the positive space, inclusions, leading block, the kernel}
\nt{$E$, $\mathbb{A}_{E,f}$, $\widehat{O}_E$, $\mathsf{S}$, $\mathsf{g}$, $h_c$, $\mathsf{x}_c$, $\mathrm{pl}_j$, $\xi$}{CM field, finite adeles, integral adeles, filling and its objects}
\nt{$\sigma$, $\tau$, $\pi$, $\alpha$, $\iota$, $\mathsf{m}$}{finite index sets ($\pi$ also the number); $\alpha$, $\gamma$, $\delta$ also indices}
\nt{$a$, $b$, $c$, $d$, $i$, $j$, $l$, $m$, $t$, $s$, $f$, $N$, $M$, $P$, $C$, $D$, $B$, $A$}{indices, numbers and constants of the statements; matrices and sets of the proofs}
\nt{$g$, $u$, $v$, $w$, $x$, $y$, $\omega$, $\gamma$}{group elements and vectors}
\nt{$\varphi$, $\chi$, $\phi$, $\mu$, $\nu$, $\Omega_0$, $\epsilon$, $\Lambda$, $\Lambda_a$, $\mathcal{G}$, $\mathcal{A}$, $\mathcal{V}$, $\mathcal{C}$}{functions, measures, sets, lattices and auxiliary objects of the proofs}
\nt{$\mathsf{A}$, $\mathsf{B}$, $\mathsf{C}$, $\mathsf{G}$, $\mathsf{L}$, $\mathsf{N}$, $\mathsf{R}$, $\mathsf{T}$, $\mathsf{W}$, $\mathsf{Y}$, $\mathsf{l}$, $\mathsf{n}$, $\mathsf{q}$, $\mathsf{r}$, $\mathsf{t}$, $\mathsf{y}$, $\mathsf{z}$, $\Theta$}{matrices, numbers and maps of the proofs}
\nt{$G$, $H$, $J$, $O$, $Q$, $S$, $U$, $X$, $Z$, $F$, $K$, $k$, $h$, $q$, $o$, $n$, $r$, $z$, $\mathbf{1}$, $\infty$, $\mathrm{e}$, $\mathrm{i}$}{further matrices and elements; letters inside names; the indicator; infinity; the numbers $\mathrm{e}$ and $\mathrm{i}$}
\nt{$\varrho$, $\kappa$, $\delta$, $\zeta$, $\psi$}{letters of names of Note~T and of the place}
\end{multicols}
\endgroup

\end{document}
